% !TEX root = ../Dissertation.tex

% ist ist anzunehmen / lässt sich daraus schließen 
% mit einer gewissen vorsicht
% gemessen an...
% nur moderierende Faktoren, die mit mehr oder weniger XY in Verbindung gebracht werden können
% (keine) kausalen, also ursächlichen Faktoren / kausaler Zusammenhang
% XY lässt Vorhersagen zu, ob...
% XY kann mit YZ in Verbindung gebracht werden
% Es handelt sich dabei zunächst um eine Korrelation, keine Kausalität
% Nichtsdestotrotz liegt die Hypothese nahe, dass...
% hart und belastbar
% interessiert an relativen Interaktionen bzw. Unterschieden: Satz A wird signifikant höher gerankt als Satz B
% auf Floor-/Ceiling-Effekte eingehen
% skalenbasierte Beschreibung
%Im Verfahren des Chi-Quadrats wurde die Nullhypothese höchstsignikant (p < 0,0005) widerlegt, ersichtlich
% Bezüglich des empirischen Vorhabens der vorliegenden Arbeit wurden zwei Ziele formuliert. Erstens sollten dXY statistisch geprüft werden, zweitens sollten statistische Tests auch den Einfluss XY auf YZ feststellen




\chapter{{Experimentelle Studien -- Teil II}} \label{Experimente}
In diesem Kapitel wird eine weitere experimentelle Untersuchung präsentiert, die darauf abzielt, mögliche Unterschiede zwischen den Mitgliedern der Intensiviererklassen empirisch zu testen. So ist im Sinne von {\citet[87]{Biedermann.1969}} und \citet[408]{Breindl.2007} von einem wahrnehmbaren Unterschied in den Stärkegraden von Intensivierern auszugehen, der sich skalar erfassen lässt. Auf dieser theoretischen Basis ist implizit anzunehmen, dass sich Intensivierer in Abhängigkeit ihrer Semantik in Bezug auf den durch sie ausgedrückten Intensitätsgrad voneinander unterscheiden und sich die Expressivität der Ausdrücke auf deren Skalaritätsgrad auswirkt, wie \citet[133]{Gutzmann.2019a} in seiner Definition nahelegt. Diese potenziellen Abstufungen zu überprüfen, ist das Hauptziel des Experiments.


Die dargelegte Untersuchung setzt sich aus zwei separaten Teilen zusammen: Während im ersten Rating-Experiment der Grad an Skalarität erfragt wird (Experiment 2a), geht es im zweiten um den Grad an Sprechereinstellung (Experiment 2b). Um die Ausprägung der bezeichneten Merkmale zu erheben, werden Sätze mit deskriptiven und expressiven Intensivierern mit intuitiv unterschiedlich wahrnehmbarer Intensität wie im vorherigen Experiment von Versuchspersonen auf einer Skala eingeschätzt. In den ersten Abschnitten werden der theoretische Rahmen, der Untersuchungsaufbau und die zu überprüfende Hypothese vorgestellt. Der darauffolgende Part zerfällt einerseits in die Untersuchung von Skalarität und andererseits in die Untersuchung von Sprechereinstellung. In diesem Zusammenhang wird jeweils auf den Untersuchungsgegenstand, die Durchführung, die Personenstichprobe und die Auswertungskriterien eingegangen. Danach werden deskriptiv- und inferenzstatistische Datenanalysen vorgenommen, an die sich eine erste Ergebnisinterpretation anschließt. Im Rahmen einer Gesamtdiskussion werden die Ergebnisse der beiden Experimente in Hinsicht auf die zu untersuchende Hypothese begutachtet, woraufhin in einem Fazit die durch die Untersuchung erbrachten Erkenntnisse resümiert und {bewertet werden}.

%irgendwo Ergebnisse der letzten Analyse einbauen, die jetzt noch aussteht




\section{Theoretischer Rahmen}\label{abs:theo}
Die Untersuchung wurde aus der Annahme heraus durchgeführt, dass Sätzen mit expressiven Intensivierern sowohl hinsichtlich Skalarität als auch Sprechereinstellung grundsätzlich eine höhere Skalenposition zugeschrieben wird als Sätzen mit deskriptiven Intensivierern (vgl. {\citealt[315]{Lang.1983}}; \citealt[22]{Paradis.1997}; \citealt[57]{Ortner.2014}; \citealt[133]{Gutzmann.2019a}). Dabei ist {\citet[87]{Biedermann.1969}} und {\citet[408]{Breindl.2007}} zufolge jedoch anzunehmen, dass sich auch innerhalb der Intensiviererklassen ebensolche graduellen Unterschiede beobachten lassen, die auf die Semantik der Ausdrücke zurückzuführen sind. Dieser Umstand wird in der theoretischen Literatur bislang als gültig erachtet, obwohl dessen Zutreffen, soweit mir bekannt, nicht empirisch abgesichert wurde. Dementsprechend liegt das Ziel dieser Untersuchung darin, die potenziellen Abstufungen zwischen Intensivierern unterschiedlicher Stärkegrade zum ersten Mal experimentell zu testen. Vor diesem Hintergrund wird ebenfalls überprüft, ob die beiden Merkmale Skalarität und Sprechereinstellung in einer Korrelation stehen und die Expressivität der Ausdrücke, wie von {\citet[133]{Gutzmann.2019a}} vermutet, auf deren Skalaritätsgrad einwirkt, der infolgedessen als stärker wahrgenommen wird.

%Skala
%Adjektive
%Intensivierertypen (aber eher nicht, weil die eigentlich schon in Kapitel 2 thematisiert sein sollten)

%Gottfried R. Marschall in Schmale (2011: 106): "Was man nachprüfbar messen und physikalisch bestimmen kann, ist der Inbegriff der Objektivität. --> Idee: etwas Subjektives wie Sprecherbewertung in Form einer ordinalen Skala objektiv quantifizierbar zu machen



\section{Hypothesen und Untersuchungsaufbau}\label{abs:Hypothese und Forschungsaufbau}
Diese experimentelle Untersuchung strebt einen Vergleich von Intensivierern mit unterschiedlich wahrnehmbarer Intensität in Bezug auf Skalarität und Sprechereinstellung an. Wie eingangs angeführt, wird in der Literatur die These vertreten, dass expressive Intensivierer zum einen mit einem höheren Skalaritätsgrad assoziiert und in der Folge auf einer einschlägigen Skala höher verortet werden als deskriptive Intensivierer (vgl. diesbezüglich auch {Abschnitt \ref{Intensivierertypen}}): \glqq{}[O]ne semantic difference between \textit{very} and \textit{sau} is that \textit{sau G} expresses an even higher degree of \textit{G}-ness than {\textit{very G}}.\grqq{} ({\citealt[133]{Gutzmann.2019a}}). Deswegen nimmt {\citet[133]{Gutzmann.2019a}} für die Intensiviererklassen eine nach Intensität organisierte Rangfolge an:\footnote{\citet[133]{Gutzmann.2019a} ordnet die Typen gemäß der traditionellen Konvention von Horn-Skalen (\citealt[]{Horn.1972}) in umgekehrter, abnehmender Intensitätsabfolge von expressiven Intensivierern über deskriptive Intensivierer zum Positiv an. In Analogie zu der dieser Untersuchung zugrunde gelegten Skala wurde die Reihenfolge hier abgeändert.} \smallskip


\textit{$\emptyset$ cool}\textsubscript{[POS]} $\prec$ \textit{sehr cool}\textsubscript{[DI]} $\prec$ \textit{sau cool}\textsubscript{[EI]} \smallskip

\noindent Dass diese Einschätzung im Kern korrekt ist, konnte durch das in {Kapitel \ref{Experiment0}} besprochene Experiment gezeigt werden. Zum anderen wird in der Literatur die Auffassung vertreten, dass expressive Intensivierer zudem die persönliche Einstellung des Sprechers ausdrücken, die auf propositionaler Ebene nicht overt realisiert ist: \glqq{}Mit der expressiven \textelp{} Funktion der Sprache wird eine Haltung des Sprechenden, mit anderen Worten die Wertung des Mitgeteilten bzw. die zum Ausdruck gebrachte Emotion als wichtiger Anteil einer Äußerung erfasst.\grqq{} (\citealt[57]{Ortner.2014}). Daher wird davon ausgegangen, dass ebenso der ausgedrückte Grad an Sprechereinstellung bei expressiven Intensivierern als stärker empfunden wird, weswegen sie in Hinsicht auf dieses Merkmal gleichfalls höher positioniert werden.\footnote{Der Definition nach \citet[133]{Gutzmann.2019a} folgend drücken deskriptive Intensivierer im Unterschied zu expressiven Intensivierern überhaupt keine Sprechereinstellung aus und sollten deshalb am unteren Ende einer Expressivitätsskala {positioniert werden}.} Auch dies konnte auf Grundlage der bisherigen Untersuchungsergebnisse für gültig erklärt werden. Ungeklärt ist jedoch die theoretische Annahme, dass sich ebensolche graduellen Abstufungen nicht nur zwischen, sondern auch innerhalb der Intensiviererklassen abzeichnen. So teilen bspw. {\citet[]{Biedermann.1969}} und {\citet[]{Breindl.2007}} deskriptive Intensivierer in Abhängigkeit ihres Intensitätsgrades auf einer Skala von `minimal' bis `absolut' \citep[vgl.][87]{Biedermann.1969} bzw. `gemäßigt' bis `extrem' {\citep[vgl.][408]{Breindl.2007}} ein. Darauf basierend lässt sich für meine Untersuchung folgende {Hypothese ableiten}:

\begin{itemize}
\item[\textbf{\textsc{H6}}] Innerhalb der Intensiviererklassen lassen sich sowohl im Hinblick auf Skalarität als auch auf Sprechereinstellung graduelle Abstufungen zwischen den Ausdrücken beobachten, die semantisch bedingt sind.
\end{itemize}

\noindent Um diese Hypothese empirisch zu überprüfen, wurden zwei Rating-Experimente mit parallelem Untersuchungsaufbau durchgeführt: Während Experiment 2a den Skalaritätsgrad erfragte, ermittelte Experiment 2b den Grad an Sprechereinstellung. Den Experimenten lag eine Stichprobe aus je vier deskriptiven und expressiven Intensivierern zugrunde, die sich gemäß der Einteilung von Intensivierern nach {\citet[87]{Biedermann.1969}} und {\citet[408]{Breindl.2007}} in dem durch sie ausgedrückten Intensitätsgrad voneinander unterschieden. Die Aufgabe der Versuchspersonen lag darin, den ausgedrückten Ausprägungsgrad auf einer Skala einzuschätzen. Die numerischen Skalenwerte standen in Experiment 2a für den durch den Intensivierer ausgedrückten Grad an Skalarität, in Experiment 2b für den ausgedrückten Grad an Sprechereinstellung. Diese Werte konnten im Anschluss an die Untersuchung gemappt und miteinander verglichen werden, was zur Überprüfung der von \citet[133]{Gutzmann.2019a} vermuteten Korrelation zwischen Skalarität und Sprechereinstellung diente. Dieser vertritt die Ansicht, dass sich die Expressivität auf den Skalaritätsgrad der Ausdrücke auswirkt, letzterer infolge der Expressivität also als stärker wahrgenommen wird.


Im Folgenden werden die beiden Experimente der Reihe nach vorgestellt. Begonnen wird mit Experiment 2a, das auf die Ermittlung des Skalaritätsgrades verschiedener deskriptiver und expressiver Intensivierer abzielte. 


%Seitenzahlen bei Lang und Paradis (1997: 22?) genau überprüfen und überall einbauen


\section{Experiment 2a -- Skalarität}\label{abs:Skalarität}
Die nächsten Abschnitte präsentieren das Experiment, das mit dem Ziel durchgeführt wurde, den Skalaritätsgrad von Intensivierern mit unterschiedlich wahrnehmbarer Intensität zu erheben, um den potenziellen klasseninternen Abstufungen nachzugehen und erste Rückschlüsse auf die Gültigkeit der zur Diskussion stehenden Hypothese H6 ziehen zu können.



\subsection{Untersuchungsgegenstand}\label{abs:Untersuchungsgegenstand63}
Um erste Hinweise auf das Zutreffen der oben beschriebenen {Hypothese H6} zu erhalten, wurde eine Rating-Aufgabe mit 56 Testsätzen konstruiert. Diese Items beinhalteten verschiedene Nomen-Adjektiv-Kombinationen (wie z.\,B. \textit{Die Pfütze ist groß} und \textit{Der Flur ist lang}) und wurden im Experiment jeweils mit vier deskriptiven und expressiven Intensivierern intuitiv unterschiedlicher Stärkegrade gekreuzt: \textit{äußerst}, \textit{sehr}, \textit{überaus} und \textit{ziemlich} (deskriptiv) sowie \textit{mega}, \textit{sau}, \textit{super} und \textit{voll} (expressiv). Bei den deskriptiven Intensivierern handelte es sich ausnahmslos um Ausdrücke, die sowohl {\citet[87]{Biedermann.1969}} als auch \citet[408]{Breindl.2007} im Rahmen ihrer Gradeinteilung als Vertreter einer Intensitätsgruppe (bspw. `gemäßigt', `hoch' oder `extrem') nennen und sich in Abhängigkeit des durch sie ausgedrückten Intensitätsgrades auf einer kontinuierlichen Skala wie folgt anordnen lassen: \medskip

\noindent \textit{ziemlich} (gemäßigt) $\prec$ \textit{sehr} (gemäßigt - hoch) $\prec$ \textit{überaus} (hoch) $\prec$ \textit{äußerst} (extrem) \par  \medskip

\noindent Analog zur theoretisch fundierten deskriptiven Intensiviererabfolge wurde bei der Stichprobe der expressiven Intensivierer nach persönlicher Einschätzung und unter Zuhilfenahme der von \citet[408]{Breindl.2007} vorgeschlagenen Klassifikation die folgenden Abstufungen vermutet: \medskip

\noindent \textit{voll} (gemäßigt) $\prec$ \textit{super} (gemäßigt - hoch) $\prec$ \textit{sau} (hoch) $\prec$ \textit{mega} (extrem) \medskip

\noindent Unter der Annahme, dass die soeben angeführten graduellen Abstufungen im Kern zutreffend sind, wird im Vorfeld der Untersuchung also erwartet, dass die expressiven Intensivierer jeder Stufe von den Versuchspersonen als stärker skalar beurteilt werden als die deskriptiven Intensivierer.


Auf welcher Basis die für das Experiment erforderlichen Adjektive und Nomen ausgewählt wurden, ist Gegenstand der {anschließenden Ausführungen}.


\subsubsection{Auswahl der Adjektive}
In der Untersuchung wurde auf Dimensionsadjektive zurückgegriffen {(n = 7)}. \citet[471]{Wurzel.1987} zufolge handelt es sich dabei um Ausdrücke, die durch ihre semantischen Eigenschaften zur Graduierung -- und somit auch zur Intensivierung -- geradezu prädestiniert sind. Um implizite Negation zu vermeiden, wurden gemäß der in \citet[59]{Os.1989} beschriebenen \glqq{}>\,-\,Relation\grqq{} bloß Adjektive mit positiver, nach oben gerichteter Interpretationsrichtung ausgewählt, während vom Gebrauch negativer, abwärtsgerichteter Adjektive abgesehen wurde.\footnote{In der Semantik gelten aufwärtsgerichtete bzw. Pluspol-Adjektive wie \textit{groß} oder \textit{lang}, die in der Regel mit Maßangaben stehen, als unmarkiert, während bei abwärtsgerichteten Adjektiven Negation angenommen wird, sodass die Werte auf der gegebenen Skala umgedreht werden müssen: $\lsem$klein$\rsem$ = $\neg$\,\textit{groß} (vgl. {\citealt[107]{Bierwisch.1987b}}). Durch die Umkehrung gilt die kognitive Verarbeitung von Minuspol-Adjektiven als komplexer als die von Pluspol-Adjektiven (vgl. {\citealt[502]{Wurzel.1987}}; \citealt[602]{Blutner.1987}; {\citealt[678]{Bierwisch.1987c}}).} Nachdem eine Reihe von Adjektiven mit Blick auf diese Kriterien auf ihre Tauglichkeit hin überprüft wurden, sind die folgenden in das Experiment eingeflossen: \textit{alt}, \textit{breit}, \textit{dick}, \textit{groß}, \textit{heiß}, \textit{hoch} und \textit{lang}.\footnote{Die meisten dieser Pluspol-Adjektive bezieht auch \citet[519]{Doelling.1987} in seine Untersuchung zu den logischen Eigenschaften von Dimensionsadjektiven ein.}


%Zu erkennen ist, dass sich die berücksichtigten Dimensionsadjektive trotz der besagten semantischen Eigenschaften auf unterschiedliche Größen beziehen: Mit \textit{alt} liegt ein Adjektiv vor, das zeitliche Verhältnisse beschreibt. Die Adjektive \textit{breit}, \textit{dick}, \textit{groß}, \textit{hoch} und \textit{lang} betrachten dagegen räumliche Verhältnisse,\footnote{Dieselben räumlichen Verhältnisse bzw. Objekte betreffende Pluspol-Adjektive bezieht \citet[519]{Doelling.1987} in seine Untersuchung zu den logischen Eigenschaften von Dimensionsadjektiven ein.} während \textit{heiß} eine physikalische Größe einordnet {\citep[vgl.][287]{Lang.1987}}.


%Dies bedeutet, dass durch die ausgewählten Adjektive in der Untersuchung zugleich verschiedene Größen abgedeckt werden
Wie in Abschnitt \ref{Intensivierertypen} geschildert, ist die Interpretation von Intensivausdrücken grundsätzlich stark kontextabhängig und wird größtenteils vom nachfolgenden Adjektiv bestimmt \citep[vgl.][27]{Paradis.1997}. Daher war es für diese Untersuchung wichtig, dass sowohl die Adjektive als auch die Nomen möglichst neutral waren. Dadurch konnte sichergestellt werden, dass die Unterschiede in Bezug auf den Skalaritätsgrad in erster Linie auf den Intensivierer zurückzuführen und nicht dem Einfluss des Adjektivs geschuldet sind. Dieses Kriterium bezeichne ich im Anschluss als Neutralitätskriterium. 
%Interessant ist dieses Kriterium ebenfalls hinsichtlich der von {\citet[102]{Schmidt.1983}} beschriebenen Beobachtung, dass Intensivierer präferiert mit \glqq{}positiv wertenden Lexemen oder \textelp{} Wörtern mit negativer Wertung\grqq{} auftauchten, wobei die Polarität ebendieser beibehalten werde. Entsprechend schließt sich die Frage an, was passiert, wenn der Kontext bzw. das zu intensivierende Adjektiv gänzlich neutral gehalten ist.

%umschreiben, wenn Aufsatz da und ich weiß, was ich damit sagen wollte



\subsubsection{Auswahl der Nomen}
Bei der Auswahl der Nomen (n = 56) waren zwei Kriterien relevant: Diese sollten einerseits im Kontext des verwendeten Adjektivs natürlich klingen. Eine mögliche natürliche Kombination aus Nomen und Adjektiv ist z.\,B. \textit{Flur} und \textit{lang} wie in \textit{Der Flur ist sehr lang}, wohingegen die Kombination aus \textit{Flur} und \textit{dick} wie in \textit{Der Flur ist sehr dick} nicht plausibel ist. Andererseits sollten sie gemäß des Neutralitätskriteriums möglichst neutral sein. Durch dieses Vorgehen wurde ausgeschlossen, dass die Nomen durch ihre Semantik und eine potenziell damit verbundene Konnotation emotionale Reaktionen bei den Versuchspersonen hervorriefen. In diesem Sinne wurden Sätze mit neutralem Nomen, bspw. \textit{Scheune} wie in \textit{Die Scheune ist sehr groß}, in die Untersuchung miteinbezogen, während positiv konnotierte Nomen wie \textit{Geschenk} oder negativ konnotierte Nomen wie \textit{Schmerz} aus besagten Gründen vermieden wurden. Um einen grammatischen Satz zu bilden, wurden die ausgewählten Nomen im Experiment mit dem jeweiligen {Definitartikel präsentiert.}



\subsubsection{Untersuchungsmaterial}
Beispiele für die im Experiment abgefragten Kopulasätze mit verschiedenen Kombinationen aus Nomen, Intensivierer und Adjektiv sind in (\ref{itemse1}) gelistet.

\ea  \label{itemse1}
\ea Die Pfütze ist äußerst groß.
\ex Der Flur ist sehr lang.
\ex Die Mauer ist überaus hoch.
\ex Das Gemälde ist ziemlich alt.
\ex Die Straße ist mega breit.
\ex Das Wasser ist sau heiß.
\ex Der Stamm ist super dick.
\ex Das Schiff ist voll groß.
\z
\z
Im Experiment wurde eine unabhängige Variable mit 8 Stufen untersucht: 4 deskriptive und 4 expressive Intensivierer. Bei den Nomen und Adjektiven handelte es sich um feste Kombinationen, die sich über die Listen hinweg nicht unterschieden. So trat z.\,B. \textit{Die Pfütze} in jeder Liste zusammen mit \textit{groß} auf. Was je Liste jedoch variierte, war der verwendete Intensivierer. Demzufolge rotierte jede konstante Kombination aus Nomen und Adjektiv (n = 7) im Experiment über die Intensivierer. Dadurch konnte für jeden Intensivierer die gleiche Anzahl von Datenpunkten gesammelt werden. Die abzufragenden Items wurden nach dem lateinischen Quadrat auf acht Versuchspersonengruppen verteilt, sodass jede Person eine dieser acht Listen mit insgesamt 56 Testsätzen (d.\,h. 28 deskriptive und 28 expressive Items) zu beurteilen hatte. Aufgrund des gewählten Within Subjects-Designs sah jede Versuchsperson jedes mögliche Item genau ein Mal. Die Abfolge der Items wurde nach dem Adjektiv pseudorandomisiert. Durch diese scheinbar zufällige Randomisierung wurde verhindert, dass das gleiche Adjektiv im Experiment unmittelbar aufeinander folgt, was im äußersten Fall zu einer ähnlichen Skaleneinschätzung geführt hätte.


Zusätzlich zu den für die Auswertung relevanten Items wurden vier Übungssätze mit gleicher syntaktischer Struktur, jedoch anderem lexikalischen Material erstellt. Wie schon im vorangegangenen Experiment konnten sich die Versuchspersonen mittels dieser Sätze zu Beginn des Experiments mit der Aufgabe vertraut machen und ein Gespür für die Betätigung des Schiebereglers sowie den Umgang mit der Skala entwickeln. Distraktoren wurden in das Experiment nicht miteinbezogen. Erneut bestand die {Intention} darin, die Intensivierer relativ zueinander in Beziehung setzen und den durch sie ausgedrückten Grad an Skalarität evaluieren zu lassen. Deswegen wurden Ablenkersätze als irrelevant eingestuft.





\subsection{Untersuchungsdurchführung}\label{abs:Untersuchungsdurchführung112}
Die Aufgabe der Versuchspersonen lag darin, am Computer kurze und syntaktisch einfach konstruierte Kopulasätze  wie die in (\ref{itemse1}) in Bezug auf den darin ausgedrückten Grad an Skalarität zu bewerten. In diesen wurden verschiedene Objekte beschrieben, wie bspw. \textit{Der Berg ist super hoch}. Um ihre Einschätzung abzugeben, sollten sie einen Schieberegler auf diejenige Skalenposition setzen, die ihrer Ansicht nach das bezeichnete Objekt im Vergleich zu einem normal ausgeprägten Objekt am treffendsten repräsentiert. Zur Veranschaulichung ist die in {Experiment 2a} verwendete Skala in {Abbildung \ref{Bewertungsskala} dargestellt}.

%
\begin{figure}

  \includegraphics[width=\textwidth]{Skala1_2}
  \caption[Darstellung der in Experiment 2a verwendeten Bewertungsskala zu Skalarität] {Darstellung der in Experiment 2a verwendeten Bewertungsskala.}  
  \label{Bewertungsskala}
\end{figure}


\noindent Wie in {Abbildung \ref{Bewertungsskala}} zu sehen ist, unterschied sich die Skala in ihrer Gestaltung nicht von den in Kapitel \ref{Experiment0} beschriebenen Skalen zu Experiment 1: Während sie vor dem Betätigen des Schiebereglers farblos war, färbte sie sich mit dem Betätigen vom Nullpunkt bis zum gewählten Punkt hin ein, um den Versuchspersonen die zugeschriebene Intensität visuell zu verdeutlichen. Beschriftet war die Skala am linken Rand mit `normal' und am rechten Rand mit `maximal'. Lediglich im Hinblick auf die Aufgabenstellung unterschied sich diese Skala von den vorherigen: In Experiment 2a stand über der Skala der von den Versuchspersonen zu bewertende Satz, darunter die Aufgabenstellung, die im Wortlaut dem Itemsatz angepasst und in dem das jeweilige Adjektiv aus Gründen der Salienz fett hervorhoben war. Am Beispiel also: \textit{Geben Sie mit dem Schieberegler an, \textbf{wie hoch} der Berg ist}.
\noindent Um die Aufgabe zu erfüllen, war der Schieberegler durch Klicken oder Ziehen mit der Maus an diejenige Position zu bewegen, die der persönlichen Einschätzung nach am besten der Bedeutung des Satzes entspricht. Der gewählte Skalenwert steht dabei für den Ausprägungsgrad, den die Person mit der im Satz beschriebenen Eigenschaft, bspw. der Höhe des Bergs, {verbindet}.


Vor Beginn des Experiments wurden die Versuchspersonen in einer ausführlichen Instruktion über den Anlass, die Dauer und den Ablauf des Experiments sowie datenschutzrechtliche Aspekte informiert. Darüber hinaus wurden sie dazu aufgefordert, das Experiment an einem Endgerät mit physischer Tastatur durchzuführen. Durch diesen Hinweis sollte verhindert werden, dass die Skalen aufgrund unterschiedlicher Bildschirmgröße (Smartphone/Tablet vs. Notebook/Desktop-PC) nicht korrekt angezeigt oder nicht wie vorgesehen bedient werden können.\footnote{Zur Kontrolle wurde auch hier die Bildschirmgröße der Versuchspersonen erfasst.} Nach dem Einleitungstext erschienen zunächst die Übungssätze, die ebenso wie die Items einzeln abgefragt wurden und nach der Beantwortung unmittelbar und ohne weitere Erklärung in die für die Auswertung relevanten Testsätze übergingen.




\subsection{Versuchspersonen}\label{abs:Versuchspersonen2}
Die Personenstichprobe setzte sich aus 63 deutschen Muttersprachler:innen im Alter von 18 bis 40 Jahren zusammen, die keine Kenntnis über das Ziel der Untersuchung hatten. Diese wurden per \textit{Clickworker} rekrutiert und (bemessen an dem zu der Zeit gültigen deutschen Mindestlohn) mit je 1,70\,Euro entlohnt.
Auch dieses Experiment wurde mittels \textit{LimeSurvey} erstellt und von den Versuchspersonen individuell online bearbeitet. Die Dauer des Experiments betrug im Schnitt ca. zehn Minuten. Die Teilnahme erfolgte vollkommen anonym und ohne Erhebung personenbezogener Daten wie Alter, Geschlecht oder Herkunft.


Um zu gewährleisten, dass von jedem Intensivierer pro Liste sieben Datenpunkte und somit eine gleichmäßige Verteilung vorliegt, wurde die Teilnahme von 64 Versuchspersonen angestrebt. Nach dem Abschluss des Experiments kam es allerdings zu einer Komplikation vonseiten \textit{Clickworkers}, in deren Folge bloß Daten von 63 Personen zum Export zur Verfügung standen. Dadurch konnten letztlich fünf Listen mit Daten von je acht Versuchspersonen, zwei Listen mit Daten von je sieben Versuchspersonen sowie eine Liste mit Daten von neun Versuchspersonen exportiert werden, was in einem leichten Ungleichgewicht resultierte. Wie bereits beim vorangegangenen Experiment wurde dies aufgrund der Robustheit der gewählten statistischen Auswertungsmethode sowie dem häufig unvermeidbaren Ausschluss einzelner Versuchspersonen jedoch als nicht weiter problematisch erachtet.




\subsection{Ausschlusskriterien}
Auch für diese Untersuchung wurden im Vorfeld Ausschlusskriterien festgelegt. Wie schon in Experiment 1 beinhalteten diese vor allem die Kombination aus einer überschnellen Reaktionszeit und einem uniformen Antwortmuster, die ein bloßes Durchklicken der Person durch das Experiment nahelegt.


Von den 63 Versuchspersonen, deren Antworten exportiert werden konnten, gingen 61 in die Auswertung mit ein, wohingegen zwei Personen von der Analyse ausgeklammert wurden. Der Grund für den Ausschluss ebendieser lag darin, dass deren Ergebnis ein völlig einheitliches Muster ohne Variation in der Skalenbewertung aufwies. So positionierten sie den Schieberegler bei allen Intensivierern gleichermaßen am linken oder rechten Skalenrand, was ein Indiz dafür ist, dass sie die Aufgabe nicht aufmerksam bearbeitet, sondern sich lediglich durch das Experiment durchgeklickt haben. Diese Vermutung konnte durch eine Begutachtung der durchschnittlichen Reaktionszeit je Item erhärtet werden: Betrug die Reaktionszeit aller Versuchspersonen pro Item über das Experiment hinweg im Schnitt etwa 5,07 Sekunden ({\textit{SD} = 13,1}), lag die der ausgeschlossenen Personen deutlich darunter, nämlich bei 1,38 Sekunden ({\textit{SD} = 1,1}) und 1,91 Sekunden (\textit{SD} = 1,0). Dies erscheint kaum ausreichend, um einen gelesenen Satz korrekt zu verarbeiten und zu bewerten. Gemessen an den vergleichsweise niedrigen Reaktionszeiten und der uniformen Skaleneinschätzung war der Ausschluss dieser Personen legitim.
%1.38 Sekunden bei Proband 585636.4 bzw. bei 1.91 Sekunden bei Proband 




\subsection{Ergebnisse}\label{abs:Ergebnisse147}
Nachfolgend werden die Resultate von Experiment 2a vorgestellt, das auf die Ermittlung des Skalaritätsgrades von Intensivierern mit unterschiedlich wahrnehmbarer Intensität abzielte. Dieses diente der Überprüfung der Hypothese H6, nach der sich möglicherweise nicht nur zwischen, sondern auch innerhalb der Intensiviererklassen Unterschiede bzw. graduelle Abstufungen in Bezug auf den durch die Intensivierer ausgedrückten Intensitätsgrad zeigen, wie \citet[87]{Biedermann.1969} und \citet[408]{Breindl.2007} postulieren.


Wie beim vorangegangenen Experiment wird mit einer deskriptivstatistischen Auswertung begonnen, an die sich eine mittels R (Version 1.1.456, \citealt[]{R.2018}) und den Paketen \textit{lme4} (\citealt[]{Bates.2015}) und \textit{lmerTest} (\citealt[]{Kuznetsova.2017}) vorgenommene inferenzstatistische Auswertung anschließt. Die Visualisierung der Daten erfolgte zudem mit \textit{ggplot2} {(\citealt[]{Wickham.2016}}).

%Wie zuvor wurde letztere mittels R (Version 1.1.456, \citealt[]{R.2018}) und den Paketen \textit{lme4} (\citealt[]{Bates.2015}) und \textit{lmerTest} (\citealt[]{Kuznetsova.2017}) vorgenommen. Die Visualisierung der Daten erfolgte weiterhin mit \textit{ggplot2} (\citealt[]{Wickham.2016}). 


\subsubsection{Vergleich der Stufen}\label{abs:Deskriptivstatistik}


%Im Rahmen der deskriptivstatistischen Analyse werden zunächst die beiden Intensiviererklassen miteinander verglichen. Daran schließt sich eine Detailbetrachtung der einzelnen Intensivierer an, was einen besseren Überblick über die von den Versuchspersonen zugewiesenen Skalenwerte ermöglicht.



%\subsubsection{Vergleich der Intensiviererklassen} $~$ \\
\noindent Tabelle \ref{tab:Skalenwert1} bildet die gemittelten Skalenwerte und die Standardabweichungen für die beiden Intensiviererklassen, basierend auf 427 Datenpunkten pro Intensivierer, ab. An dieser Stelle werden die jeweils vier deskriptiven und expressiven Intensivierer zu einer Stufe zusammengefasst. Zusätzlich ist der Mittelwert für beide Stufen zusammen angegeben. 
%Die Anordnung erfolgt in alphabetischer Reihenfolge.

%
\begin{table}
%\setlength{\tabcolsep}{28.75pt}
\begin{tabularx}{.8\textwidth}{lrrrr}
\lsptoprule
Intensiviererklasse & Skalenwert &  (\textit{SD}) & Mittelwert &  \textit{SD}\\
\midrule
deskriptiv          & 58,18      &  (23,5)        & 60,43      & 23,8\\
expressiv           & 62,67      & (23,9)         & 60,43      & 23,8\\
\lspbottomrule
\end{tabularx}
\caption[Gemittelte Skalenwerte der Stufen bei Experiment 2a]{Gemittelte Skalenwerte und Standardabweichungen der Intensiviererklassen bei Experiment 2a -- Skalarität (in alphabetischer Reihenfolge).}
\label{tab:Skalenwert1}
\end{table}

\noindent Die Werte aus Tabelle \ref{tab:Skalenwert1} geben insgesamt Hinweise darauf, dass Sätze mit expressiven Intensivierern den Erwartungen gemäß von den Versuchspersonen hinsichtlich Skalarität auf der Skala um 4,50 Punkte höher eingestuft wurden als Sätze, die mit deskriptiven Intensivierern modifiziert wurden. Demnach wird der Skalaritätsgrad von expressiven Intensivierern bereits auf den ersten Blick als höher wahrgenommen als der von deskriptiven Intensivierern. Auch wenn der Unterschied zwischen den beiden Intensiviererklassen nicht übermäßig stark ausgeprägt ist, zeigt sich dennoch eine entsprechende Tendenz. Dieses Ergebnis deutet darauf hin, dass das Skalenrating als Untersuchungsmethode wie beabsichtigt funktioniert hat.


Die Tendenz zugunsten der expressiven Stufe ist ebenso anhand der in Abbildung \ref{Skalenwert2} dargestellten Boxplots zu erkennen.


\begin{figure}
\includegraphics[width=1\textwidth]{Bedingung_E1_neu2}  
\caption[Gemittelte Skalenwerte der Stufen bei Experiment 2a] {Gemittelte Skalenwerte der Intensiviererklassen bei Experiment 2a -- Skalarität (x-Achse = Intensiviererklasse, y-Achse = Skalenwert).}
\label{Skalenwert2}
\end{figure}



\subsubsection{Regressionsanalyse zum Vergleich der Stufen}\label{LiGeMo1}


\noindent Nachfolgend werden die Ergebnisse der Regressionsanalyse zur Untersuchung von Skalarität vorgestellt. Das finale Modell wurde ausgehend vom vollen Modell per Rückwärtsselektion und Modellvergleich mithilfe der Funktion \textit{anova} in R schrittweise berechnet. Dazu wurde zunächst ein volles Regressionsmodell mit den Effekten aller unabhängigen Variablen und deren Zwei-Wege-Interaktionen berechnet, die möglicherweise Einfluss auf die abhängige Variable, d.\,h. den Skalaritätsgrad, hatten. Das volle Modell beinhaltete \textit{Fixed Effects} für den \textsc{Skalenwert} als metrische abhängige Variable sowie für die kategorialen Prädiktoren \textsc{Intensiviererklasse} und \textsc{Position} sowie deren Zwei-Wege-Interaktionen. In Experiment 2a steht die Variable \textsc{Skalenwert} für den Grad an Skalarität, den die Versuchspersonen mit der jeweiligen Intensiviererklasse bzw. dem jeweiligen Intensivierer assoziiert haben. Um einen potenziellen graduellen Effekt zwischen den Intensiviererklassen zu finden, wurde die nominale Variable \textsc{Intensiviererklasse} dummy-gecoded, wobei den Ausprägungsformen `Deskriptiv' und `Expressiv' die Werte $0$ und $1$ zugewiesen wurden. Diese wurden daraufhin miteinander verglichen. Außerdem wurden im vollen Modell versuchspersonenweise \textit{Random Intercepts} und \textit{Random Slopes} für die Variablen \textsc{Skalenwert}, \textsc{Intensiviererklasse} und \textsc{Position} miteingerechnet. Das bedeutet, dass auch mögliche Unterschiede zwischen den Versuchspersonen berücksichtigt wurden.\footnote{Da das volle Regressionsmodell unter Einbezug der Adjektive nicht konvergierte, konnten bei dieser Berechnung lediglich versuchspersonenweise \textit{Random Intercepts} und \textit{Random Slopes} miteingebunden werden.} Ausgehend vom vollen Modell wurden durch Rückwärtsselektion und Modellvergleich nach und nach all diejenigen Effekte getilgt, die den Modell-Fit nicht signifikant verbesserten. Das volle und finale Modell ist im Anhang \ref{RegX} zu sehen. Das finale Regressionsmodell enthält signifikante Haupteffekte für die Prädiktoren \textsc{Intensiviererklasse} und \textsc{Position}. Des Weiteren ist deren Zwei-Wege-Interaktion ebenfalls signifikant. Die berechneten Koeffizienten der Regressionsanalyse sind in {Tabelle \ref{LinReg01}} zusammengetragen.

%
\begin{table}
%\setlength{\tabcolsep}{8.75pt}
\begin{tabularx}{\textwidth}{Xrrrrr}
\lsptoprule
Prädiktor & \textit{Est.} & \textit{SE} & \textit{$\chi$$^{2}$} & \textit{p} & \\
\midrule

(Intercept) & 60,07 & 2,06 & 170,70 & < 0,001 & $\star$$\star$$\star$ \\
Intensiviererklasse & \textcolor{white}{-}4,47 & 1,47 & \textcolor{white}{-}\textcolor{white}{-}8,81 & < 0,01 & $\star$$\star$ \\
Position & -2,90 & 0,69 & \textcolor{white}{-}16,20 & < 0,001 & $\star$$\star$$\star$ \\
Intensiviererklasse\,:\,Position & -1,23 & 0,60 & \textcolor{white}{-}\textcolor{white}{-}4,17 & < 0,05 & \multicolumn{1}{r}{$\star$} \\
\lspbottomrule

\end{tabularx}
\caption[Finales Regressionsmodell von Experiment 2a (mit `Deskriptiv' als Referenzlevel)]{Finales Regressionsmodell von Experiment 2a -- Skalarität (mit `Deskriptiv' als Referenzlevel).}
\label{LinReg01}
\end{table}



%% IK: Wie verändert sich die Rating in Abhängigkeit von IK
%% Position numerisch
%% \textcolor{white}{-} So kann man einfach ein "weißes" Minus davorschreiben, das man nicht sieht, für ums unnernanner zu setzen


\noindent Wie der Übersicht in Tabelle \ref{LinReg01} zu entnehmen ist, haben zwei unabhängige Variablen sowie deren Zwei-Wege-Interaktion einen signifikanten Effekt auf den Skalenwert: \textsc{Intensiviererklasse} und \textsc{Position}. Demzufolge werden Sätze mit expressiven Intensivierern auf der Skala signifikant höher verortet als Sätze mit deskriptiven Intensivierern ($\chi$$^{2}$\,(1) = 8,81, \textit{p} < 0,01). Dieses Ergebnis steht im Einklang mit der in Experiment 1 untersuchten Hypothese H3 nach \citet[22]{Paradis.1997} und \citet[133]{Gutzmann.2019a} (vgl. Kapitel \ref{Experiment0}). Hinzu kommt ein signifikanter Positionseffekt: Je weiter fortgeschritten die Position eines Items im Experiment ist, desto niedriger wird es bezüglich Skalarität auf der Skala beurteilt {($\chi$$^{2}$\,(1) = 16,20,} {\textit{p} < 0,001}). Ferner ist die Interaktion zwischen beiden unabhängigen Variablen signifikant ($\chi$$^{2}$\,(1) = 4,17, \textit{p} < 0,05), was darauf hindeutet, dass sich die beiden Intensiviererklassen im Laufe des Experiments einander angenähert haben und der Unterschied zwischen deskriptiven und expressiven Intensivierern folglich kleiner geworden ist. Dies legt einen Gewöhnungseffekt nahe, der dem uniformen Untersuchungsdesign geschuldet sein kann. Somit besteht die Möglichkeit, dass sich die Versuchspersonen während des Experiments an die Intensivierer gewöhnt haben, weswegen sich der zu Beginn als stärker empfundene Unterschied zwischen den Intensiviererklassen sukzessive verringert hat.


Kurz gefasst belegen die Ergebnisse von {Experiment 2a} den in der Literatur vertretenen Unterschied zwischen deskriptiven und expressiven Intensivierern (vgl. \citealt[22]{Paradis.1997}; {\citealt[133]{Gutzmann.2019a}}). So wurden Sätze mit expressiven Intensivierern von den Versuchspersonen in der Tat mit einem höheren Grad an Skalarität assoziiert als Sätze, die mit einem deskriptiven Intensivierer präsentiert wurden. Der Umstand, dass sich der Unterschied zwischen den beiden Intensiviererklassen in aller Deutlichkeit zeigt, lässt außerdem den Schluss zu, dass das Skalenrating als Untersuchungsmethode wie intendiert funktioniert hat.



\subsubsection{Vergleich der Intensivierer}\label{LiGe}


\noindent Im nächsten Analyseschritt werden die gemittelten Skalenwerte der einzelnen Intensivierer betrachtet. Durch dieses Vorgehen können die untersuchten Ausdrücke in Bezug auf den ihnen von den Versuchspersonen zugewiesenen Skalaritätsgrad miteinander verglichen werden. Diese Werte finden sich inklusive Standardabweichungen in Tabelle \ref{Skalenwert3}. Die Anordnung der Intensivierer je Klasse erfolgt nach ansteigendem Skalenwert.

%
\begin{table}
%\setlength{\tabcolsep}{27pt}
\begin{tabularx}{.8\textwidth}{Xlrr}
\lsptoprule
Intensiviererklasse & Intensivierer  & Skalenwert &  \textit{SD}\\
\midrule
\rowcolor[gray]{.90}
deskriptiv & \textit{ziemlich} & 51,22 & 23,6\\
           & \textit{äußerst}  & 59,46 & 23,4\\
 \rowcolor[gray]{.90}
           & \textit{sehr}     & 59,69 & 23,3\\
           & \textit{überaus}  & 62,36 & 22,0\\
 \rowcolor[gray]{.90}
expressiv  & \textit{voll}     & 56,38 & 23,4\\
           & \textit{super}    & 63,15 & 22,2\\
 \rowcolor[gray]{.90}
           & \textit{sau}      & 63,87 & 24,4\\
           & \textit{mega}     & 67,27 & 24,4\\
\lspbottomrule

\end{tabularx}
\caption[Gemittelte Skalenwerte der einzelnen Intensivierer von Experiment 2a]{Gemittelte Skalenwerte und Standardabweichungen der einzelnen Intensivierer von Experiment 2a -- Skalarität (angeordnet nach ansteigendem Skalenwert).}
\label{Skalenwert3}
\end{table}

%
\noindent Aus den in Tabelle \ref{Skalenwert3} aufgeführten Werten lässt sich ableiten, dass den deskriptiven Intensivierern auch im Einzelnen betrachtet von den Versuchspersonen eine niedrigere Skalenposition zugesprochen wurde als den expressiven Intensivierern. So liegen beinahe alle deskriptiven Intensivierer auf der Skala durchschnittlich im Bereich zwischen 50 und 60 Punkten, wohingegen die expressiven Intensivierer vermehrt in einem Skalenbereich zwischen 60 und 70 Punkten angesiedelt sind. Nichtsdestotrotz scheint es in beiden Intensiviererklassen eine Ausnahme zu geben, nämlich \textit{überaus} aufseiten der deskriptiven Intensivierer und \textit{voll} aufseiten der expressiven Intensivierer. Das Ergebnis dieser Ausdrücke ist insofern überraschend, als dem deskriptiven Intensivierer ein verhältnismäßig hoher, dem expressiven Intensivierer hingegen ein verhältnismäßig niedriger Skalenwert zugewiesen wurde. Darüber hinaus manifestieren sich in beiden Intensiviererklassen tatsächlich graduelle Abstufungen, wie die Übersicht in Tabelle \ref{schwach0} nahelegt.


\begin{table}
\small
%\setlength{\tabcolsep}{8.25pt}
%\begin{tabular}{p{0.15\textwidth} p{0.1\textwidth} p{0.15\textwidth} p{0.1\textwidth} p{0.15\textwidth}}
\begin{tabularx}{\textwidth}{X@{} lrr lrr }
\lsptoprule
 &  \multicolumn{6}{c}{\textbf{Intensiviererklasse}}  \\
\cmidrule(lr){2-7}\cmidrule(lr){5-7}
& \multicolumn{3}{c}{{deskriptiv}} & \multicolumn{3}{c}{{expressiv}}  \\
\cmidrule(lr){2-4}\cmidrule(lr){5-7}
{Stärkegrad}  & Intensivierer & Skalenwert& \textit{SD}  & Intensivierer & Skalenwert& \textit{SD}   \\
\midrule
schwach & \textit{ziemlich} & 51,22& 23,6 & \textit{voll} & 56,38 & 23,4 \\
\multicolumn{1}{l}{mittel} & \textit{äußerst} & 59,46& 23,4  & \textit{super} & 63,15 & 22,2\\
 & \textit{sehr} & 59,69& 23,3 & \textit{sau} & 63,87& 24,4\\
stark & \textit{überaus} & 62,36& 22,0 & \textit{mega} & 67,27 & 24,4 \\
\lspbottomrule

\end{tabularx}
\caption[Klasseninterne Abstufungen bei Experiment 2a]{Klasseninterne Abstufungen bei Experiment 2a -- Skalarität (angeordnet nach ansteigendem Skalenwert).}
\label{schwach0}
\end{table}

\noindent Wie in Tabelle \ref{schwach0} zu beobachten ist, zeichnet sich innerhalb der deskriptiven Intensiviererklasse im Hinblick auf die von den Versuchspersonen wahrgenommene Skalarität folgendes dreistufiges Muster ab: {\textit{ziemlich} $\prec$ \textit{äußerst}|\textit{sehr} $\prec$ \textit{überaus}}. Ähnliches findet sich innerhalb der expressiven Intensiviererklasse: \textit{voll} $\prec$ \textit{super}|\textit{sau} $\prec$ \textit{mega}. Demnach wurde den deskriptiven Intensivierern \textit{äußerst} und \textit{sehr} von den Versuchspersonen ein nahezu identischer Skalenwert zugeschrieben ({$\Delta$ = ca. 0,23 Punkte}), genauso wie die expressiven Intensivierer \textit{super} und \textit{sau} auf der Skala ähnlich hoch positioniert wurden ($\Delta$ = ca. 0,72 Punkte). Diese Abstufungen lassen sich auch mithilfe des in Abbildung \ref{Bar1} dargestellten Säulendiagramms erkennen. Im Anhang \ref{Aja1} ist in Tabelle \ref{Adj0} außerdem eine Übersicht über die Ratings der Adjektive beigefügt. Darin ist zu sehen, dass deren Mittelwerte vergleichbar sind: Während das am niedrigsten eingestufte Adjektiv \textit{breit} auf der Skala durchschnittlich bei 58,06 Punkten angesiedelt ist, befindet sich das am höchsten verortete Adjektiv \textit{heiß} im Schnitt bei 64,80 Punkten. Die nur geringe Differenz von ca. 6,74 Punkten deutet auf ein Funktionieren der Methode hin. Andernfalls ließe sich spekulieren, dass der verzeichnete Unterschied zwischen den Intensiviererklassen ebenso den im Experiment verwendeten Adjektiven geschuldet sein könnte. Dies ist allerdings nicht {der Fall}.


%Wäre diese bspw. größer, so könnte der beobachtete Unterschied zwischen den Intensiviererklassen durch die verwendeten Adjektive geschuldet sein.



 
%
\begin{figure}
\includegraphics[width=1\textwidth]{Rplot_Bar_E1}
\caption[Säulendiagramm zu den einzelnen Intensivierern von Experiment 2a] {Säulendiagramm zu den einzelnen deskriptiven und expressiven Intensivierern von Experiment 2a -- Skalarität (x-Achse = Intensivierer, y-Achse = Skalenwert).}
\label{Bar1}
\end{figure}


\subsubsection{Paarweise Vergleiche der Intensivierer}\label{Paarweiser VerBie dengleich}


Um die beobachteten Unterschiede zwischen den einzelnen Ausdrücken auf Signifikanz hin zu überprüfen, wurden in einem weiteren Analyseschritt Paarvergleiche vorgenommen. Im Zuge dessen wurden die zugeschriebenen Skalenwerte der Intensivierer unter Zuhilfenahme des Inferenztests \textit{t}-Test gepaart miteinander verglichen. Exemplarisch am deskriptiven Intensivierer \textit{ziemlich} bedeutet die Gegenüberstellung, dass der Ausdruck erst mit \textit{äußerst} verglichen wurde, darauffolgend mit \textit{sehr}, mit \textit{überaus}, mit \textit{voll} etc. Durch dieses schrittweise Vorgehen konnte eine Matrix erstellt werden, in der das Ergebnis bzw. der Mittelwertunterschied des jeweils kontrastierten Intensiviererpaars festgehalten ist. Im Anschluss an die \textit{t}-Tests wurde eine Bonferroni-Korrektur angewandt, in deren Folge das Signifikanzniveau adjustiert wurde. Dieses Korrekturverfahren dient dazu, die bei multiplen Vergleichen mit gleicher Nullhypothese mögliche Alphafehler-Kumulierung zu vermeiden (vgl. {\citealt[249]{Wollschlaeger.2017}}; \citealt[176f.]{Winter.2020}). Beim Alphafehler wird die eigentlich korrekte Nullhypothese fälschlicherweise verworfen. Denn \glqq{}[m]it der Anzahl der Testwiederholungen steigt die Wahrscheinlichkeit (Gefahr), dass ein Unterschied als signifikant erscheint, auch wenn in Wirklichkeit keiner der Unterschiede signifikant ist.\grqq{} ({\citealt[186]{Backhaus.2018}}). Bei der Adjustierung des Signifikanzniveaus nach Bonferroni wurde der \textit{p}-Wert eines jeden gegenübergestellten Intensiviererpaars mit der Gesamtzahl der Testwiederholungen  multipliziert. Weil alle Intensivierer paarweise miteinander gekreuzt wurden, wurden insgesamt 28 Vergleiche durchgeführt: \textit{ziemlich} vs. \textit{äußerst}, \textit{ziemlich} vs. \textit{sehr}, \textit{ziemlich} vs. \textit{überaus}, \textit{ziemlich} vs. \textit{voll}, \textit{ziemlich} vs. \textit{super}, \textit{ziemlich} vs. \textit{sau}, \textit{ziemlich} vs. \textit{mega}, \textit{äußerst} vs. \textit{sehr}, \textit{äußerst} vs. \textit{überaus} usw. Infolge der Multiplikation mit der Anzahl der Testwiederholungen wird der jeweilige \textit{p}-Wert größer und ist gegebenenfalls weniger oder gar nicht mehr signifikant, wodurch dem Risiko des Alphafehlers {entgegengewirkt wird}.


Die aus den Berechnungen resultierende Matrix befindet sich in {Tabelle \ref{Matrix_1_neu}}.  In ihr sind die jeweiligen Testprüfgrößen der \textit{t}-Tests und die Bonferroni-korrigierten \textit{p}-Werte der einander gegenübergestellten Intensivierer gelistet. In der Tabelle erfolgen zunächst die Vergleiche der deskriptiven Intensivierer (in ansteigendem Skalenwert), hinterher die Vergleiche der expressiven Intensivierer (ebenfalls nach ansteigendem Skalenwert sortiert). Anzumerken ist, dass bei den darin gelisteten \textit{t}-Werten aus Gründen der Darstellbarkeit auf eine ausführliche Dokumentation mit eingeklammerter Angabe der Gesamtzahl der Freiheitsgrade (englisch: \textit{number of degrees of freedom}, hier und im Folgenden abgekürzt mit \textit{df}), die auf der Größe der jeweiligen Stichprobe basiert, verzichtet wird (\textit{t}\,(\textit{\textit{df}}) = \textelp{}). Bei den vorgenommenen Paarvergleichen entspricht der Freiheitsgrad dem Wert 3346. Aus diesem Grund gilt es zu beachten, dass die in der Tabelle jeweils angegebene verkürzte Form (bspw. \textit{t} = 6,84, {\textit{p} < 0,001}) streng genommen Folgendes meint: \textit{t}\,(3346) = \textelp{}. Angewandt auf das eingeklammerte Beispiel also {\textit{t}\,(3346) = 6,84,} {\textit{p} < 0,001}). Die Formel der Paarvergleiche findet sich außerdem im Anhang \ref{RegX}.



%\cellcolor[gray]{0.90}
%\cellcolor[gray]{0.75}
%\cellcolor[gray]{0.59}


%
\begin{sidewaystable}

%\setlength{\tabcolsep}{8pt}
\begin{tabularx}{\textwidth}{l@{~~}QQQQQQ@{}Q@{}Q@{}}
\lsptoprule
Intensivierer & \textit{ziemlich} & \textit{äußerst} & \textit{sehr} & \textit{überaus} & \textit{voll} & \textit{super} & \textit{sau} & \textit{mega} \\
\midrule

\multirow{3}{*}{\textit{ziemlich}} &  & \cellcolor[gray]{0.59}\textit{t} = 6,84, & \cellcolor[gray]{0.59} \textit{t} = 6,93, & \cellcolor[gray]{0.59} \textit{t} = 8,98, & \cellcolor[gray]{0.59} \textit{t} = 4,30,  & \cellcolor[gray]{0.59} \textit{t} = 9,64, & \cellcolor[gray]{0.59} \textit{t} = 10,39, & \cellcolor[gray]{0.59} \textit{t} = 12,82,  \\ 
& \hfil \ding{55} & \cellcolor[gray]{0.59}\textit{p} < 0,001  & \cellcolor[gray]{0.59}\textit{p} < 0,001   & \cellcolor[gray]{0.59}\textit{p} < 0,001   &   \cellcolor[gray]{0.59}\textit{p} < 0,001   &  \cellcolor[gray]{0.59}\textit{p} < 0,001  & \cellcolor[gray]{0.59}\textit{p} < 0,001  &  \cellcolor[gray]{0.59}\textit{p} < 0,001   \\ 
 &  & \cellcolor[gray]{0.59} $\star$$\star$$\star$ & \cellcolor[gray]{0.59} $\star$$\star$$\star$ & \cellcolor[gray]{0.59} $\star$$\star$$\star$  & \cellcolor[gray]{0.59} $\star$$\star$$\star$ & \cellcolor[gray]{0.59} $\star$$\star$$\star$ & \cellcolor[gray]{0.59} $\star$$\star$$\star$ &  \cellcolor[gray]{0.59} $\star$$\star$$\star$ \\ 
  
\multirow{3}{*}{\textit{äußerst}} & \cellcolor[gray]{0.59} \textit{t} = -6,84, &  & \textit{t} = 0,09, &  \textit{t} = 2,14, &  \textit{t} = -2,55,  &  \textit{t} = 2,80, & \cellcolor[gray]{0.90} \textit{t} = 3,55, &  \cellcolor[gray]{0.59} \textit{t} = 5,99,  \\ 
& \cellcolor[gray]{0.59} \textit{p} < 0,001 & \hfil \ding{55} & \textit{p} > 0,05   &  \textit{p} > 0,05 &   \textit{p} > 0,05   &  \textit{p} > 0,05  & \cellcolor[gray]{0.90} \textit{p} < 0,05   & \cellcolor[gray]{0.59} \textit{p} < 0,001    \\ 
& \cellcolor[gray]{0.59} $\star$$\star$$\star$  &  &  &    &  &  & \cellcolor[gray]{0.90} $\star$ &  \cellcolor[gray]{0.59} $\star$$\star$$\star$ \\ 

\multirow{3}{*}{\textit{sehr}} & \cellcolor[gray]{0.59} \textit{t} = -6,93, & \textit{t} = -0,09, &  &  \textit{t} = 2,06, &  \textit{t} = -2,64,  &  \textit{t} = 2,72, & \cellcolor[gray]{0.90} \textit{t} = 3,47, & \cellcolor[gray]{0.59} \textit{t} = 5,91,  \\ 
& \cellcolor[gray]{0.59}\textit{p} < 0,001  & \textit{p} > 0,05  & \hfil \ding{55}  &  \textit{p} > 0,05  &   \textit{p} > 0,05   & \textit{p} > 0,05  & \cellcolor[gray]{0.90} \textit{p} < 0,05    & \cellcolor[gray]{0.59}\textit{p} < 0,001   \\ 
 & \cellcolor[gray]{0.59} $\star$$\star$$\star$ & &  &    &  &  & \cellcolor[gray]{0.90} $\star$ &  \cellcolor[gray]{0.59} $\star$$\star$$\star$  \\ 

\multirow{3}{*}{\textit{überaus}} & \cellcolor[gray]{0.59} \textit{t} = -8,98, & \textit{t} = -2,14, & \textit{t} = -2,06, &   & \cellcolor[gray]{0.59}  \textit{t} = -4,69,  &  \textit{t} = 0,66, & \textit{t} = 1,42, & \cellcolor[gray]{0.75} \textit{t} = 3,85,  \\ 
& \cellcolor[gray]{0.59}\textit{p} < 0,001 & \textit{p} > 0,05  & \textit{p} > 0,05  &  \hfil \ding{55}  &   \cellcolor[gray]{0.59}\textit{p} < 0,001  &  \textit{p} > 0,05 & \textit{p} > 0,05  & \cellcolor[gray]{0.75}\textit{p} < 0,01   \\ 
 & \cellcolor[gray]{0.59} $\star$$\star$$\star$ &  &  &   & \cellcolor[gray]{0.59} $\star$$\star$$\star$  &  & &  \cellcolor[gray]{0.75} $\star$$\star$ \\ 

\multirow{3}{*}{\textit{voll}} & \cellcolor[gray]{0.59} \textit{t} = -4,30, & \textit{t} = 2,55, & \textit{t} = 2,64, &  \cellcolor[gray]{0.59} \textit{t} = 4,69, &    &  \cellcolor[gray]{0.59} \textit{t} = 5,36, &  \cellcolor[gray]{0.59}  \textit{t} = 6,11, & \cellcolor[gray]{0.59}  \textit{t} = 8,55,  \\ 
& \cellcolor[gray]{0.59}\textit{p} < 0,001 & \textit{p} > 0,05  & \textit{p} > 0,05  &  \cellcolor[gray]{0.59}\textit{p} < 0,001 &   \hfil \ding{55}  &  \cellcolor[gray]{0.59}\textit{p} < 0,001  & \cellcolor[gray]{0.59}\textit{p} < 0,001   & \cellcolor[gray]{0.59} \textit{p} < 0,001    \\ 
 & \cellcolor[gray]{0.59} $\star$$\star$$\star$ &  &  &  \cellcolor[gray]{0.59} $\star$$\star$$\star$  &  & \cellcolor[gray]{0.59} $\star$$\star$$\star$ &\cellcolor[gray]{0.59} $\star$$\star$$\star$ & \cellcolor[gray]{0.59} $\star$$\star$$\star$  \\ 

\multirow{3}{*}{\textit{super}} & \cellcolor[gray]{0.59} \textit{t} = -9,64, & \textit{t} = -2,80, & \textit{t} = -2,72, &  \textit{t} = -0,66, & \cellcolor[gray]{0.59}  \textit{t} = -5,36,  &   & \textit{t} = 0,76, & \cellcolor[gray]{0.90} \textit{t} = 3,19,  \\ 
& \cellcolor[gray]{0.59} \textit{p} < 0,001 & \textit{p} > 0,05 & \textit{p} > 0,05   & \textit{p} > 0,05  &   \cellcolor[gray]{0.59} \textit{p} < 0,001   &  \hfil \ding{55}  & \textit{p} > 0,05   & \cellcolor[gray]{0.90} \textit{p} < 0,05    \\ 
 & \cellcolor[gray]{0.59} $\star$$\star$$\star$ &  &  &    & \cellcolor[gray]{0.59} $\star$$\star$$\star$ &  & &    \cellcolor[gray]{0.90} $\star$\\ 
 
\multirow{3}{*}{\textit{sau}} & \cellcolor[gray]{0.59} \textit{t} = -10,39, & \cellcolor[gray]{0.90} \textit{t} = -3,55, & \cellcolor[gray]{0.90} \textit{t} = -3,47, &  \textit{t} = -1,42, &  \cellcolor[gray]{0.59} \textit{t} = -6,11,  &  \textit{t} = -0,76, &  & \textit{t} = 2,43,  \\ 
& \cellcolor[gray]{0.59}\textit{p} < 0,001  &  \cellcolor[gray]{0.90} \textit{p} < 0,05  & \cellcolor[gray]{0.90} \textit{p} < 0,05  &  \textit{p} > 0,05  &  \cellcolor[gray]{0.59}\textit{p} < 0,001   &  \textit{p} > 0,05   & \hfil \ding{55}  & \textit{p} > 0,05   \\ 
 & \cellcolor[gray]{0.59} $\star$$\star$$\star$  & \cellcolor[gray]{0.90} $\star$ & \cellcolor[gray]{0.90} $\star$ &   & \cellcolor[gray]{0.59} $\star$$\star$$\star$ &  & &   \\ 

\multirow{3}{*}{\textit{mega}} & \cellcolor[gray]{0.59} \textit{t} = -12,82, &  \cellcolor[gray]{0.59} \textit{t} = -5,99, & \cellcolor[gray]{0.59}  \textit{t} = -5,91, &  \cellcolor[gray]{0.75} \textit{t} = -3,85, &  \cellcolor[gray]{0.59} \textit{t} = -8,55,  &  \cellcolor[gray]{0.90} \textit{t} = -3,19, & \textit{t} = -2,43, &  \\ 
& \cellcolor[gray]{0.59}\textit{p} < 0,001  & \cellcolor[gray]{0.59}\textit{p} < 0,001  & \cellcolor[gray]{0.59}\textit{p} < 0,001   &  \cellcolor[gray]{0.75}\textit{p} < 0,01  &   \cellcolor[gray]{0.59}\textit{p} < 0,001   &  \cellcolor[gray]{0.90} \textit{p} < 0,05   & \textit{p} > 0,05  &  \hfil \ding{55}  \\ 
 & \cellcolor[gray]{0.59} $\star$$\star$$\star$  & \cellcolor[gray]{0.59} $\star$$\star$$\star$ & \cellcolor[gray]{0.59} $\star$$\star$$\star$  &  \cellcolor[gray]{0.75} $\star$$\star$ & \cellcolor[gray]{0.59} $\star$$\star$$\star$  &  \cellcolor[gray]{0.90} $\star$ & &   \\
\lspbottomrule
\end{tabularx}
\caption[Ergebnisse der Paarvergleiche der einzelnen Intensivierer von Experiment 2a]{Ergebnisse der Paarvergleiche der einzelnen Intensivierer bei Experiment 2a -- Skalarität (nach Bonferroni-Korrektur; angeordnet nach ansteigendem Skalenwert).}
\label{Matrix_1_neu}
\end{sidewaystable}





 Wie Tabelle \ref{Matrix_1_neu} zu entnehmen ist, führten die Paarvergleiche zu einigen diskutablen Feststellungen, die anschließend kurz beschrieben werden. Für einen besseren Überblick werden bloß diejenigen Ergebnisse aufgeführt, die signifikant sind, während die restlichen an dieser Stelle außen vor gelassen werden.\footnote{Der signifikante \textit{p}-Wert wird bloß bei der ersten Erwähnung eines miteinander verglichenen Intensiviererpaars angegeben (bspw. \textit{ziemlich} vs. \textit{äußerst}), wohingegen er bei der Zweitnennung aus Gründen der Übersichtlichkeit außen vor gelassen wird (\textit{äußerst} vs. \textit{ziemlich}).} Da den deskriptiven Intensivierern von den Versuchspersonen in Bezug auf Skalarität die niedrigsten Skalenwerte zugewiesen wurden und sie in der Tabelle daher den Ausgangspunkt bilden, werden zunächst die Werte von \textit{ziemlich} betrachtet. Den Ergebnissen der Paarvergleiche zufolge unterscheidet sich \textit{ziemlich} gänzlich von den restlichen Intensivierern, sind bei diesem Ausdruck doch alle Kontraste gleichermaßen signifikant: \textit{äußerst} {(\textit{t}\,(3346) = 6,84,} {\textit{p} < 0,001),} \textit{sehr} {(\textit{t}\,(3346) = 6,93,} {\textit{p} < 0,001)}, \textit{überaus} {(\textit{t}\,(3346) = 8,98}, {\textit{p} < 0,001)}, \textit{voll} {(\textit{t}\,(3346) = 4,30,} \textit{p} < 0,001), \textit{super} ({\textit{t}\,(3346) = 9,64,} \textit{p} < 0,001), \textit{sau} {(\textit{t}\,(3346) = 10,39,} \textit{p} < 0,001) und \textit{mega} {(\textit{t}\,(3346) = 12,82}, {\textit{p} < 0,001)}. Als Nächstes werden die Ergebnisse von \textit{äußerst} inspiziert. Hier finden sich signifikante Unterschiede bei den Vergleichen mit \textit{ziemlich}, \textit{sau} ({\textit{t}\,(3346) = 3,55, \textit{p} < 0,05)} und \textit{mega} {(\textit{t}\,(3346) = 5,99}, {\textit{p} < 0,001}). Gleichsam zeigen sich bei \textit{sehr} signifikante Unterschiede bei den Kontrasten mit \textit{ziemlich}, \textit{sau} (\textit{t}\,(3346) = 3,47, {\textit{p} < 0,05}) und \textit{mega} ({\textit{t}\,(3346) = 5,91}, {\textit{p} < 0,001)}. Bei \textit{überaus} sind darüber hinaus die Vergleiche mit \textit{ziemlich}, \textit{voll} ({\textit{t}\,(3346) = -4,69}, {\textit{p} < 0,001}) und \textit{mega} ({\textit{t}\,(3346) = 3,85,} {\textit{p} < 0,01}) signifikant.


Daneben weisen auch die Vergleiche der expressiven Intensivierer auf signifikante Unterschiede hin. So treten bei \textit{voll} signifikante Unterschiede bei den Kontrasten mit \textit{ziemlich}, \textit{überaus}, \textit{super} (\textit{t}\,(3346) = 5,36, \textit{p} < 0,001), \textit{sau} (\textit{t}\,(3346) = 6,11, \textit{p} < 0,001) und \textit{mega} (\textit{t}\,(3346) = 8,55, \textit{p} < 0,001) auf. Demnach unterscheidet sich \textit{voll} zur Gänze von den restlichen expressiven Intensivierern. Der nächste expressive Intensivierer in der Reihe ist \textit{super}. Bei diesem Ausdruck sind die Vergleiche mit \textit{ziemlich}, \textit{voll} und \textit{mega} (\textit{t}\,(3346) = 3,19, {\textit{p} < 0,05)} signifikant. Wie schon erwähnt, unterscheidet sich \textit{sau} signifikant von \textit{ziemlich}, \textit{äußerst}, \textit{sehr} und \textit{voll}. Aufgrund dessen, dass die Werte des expressiven Intensivierers \textit{mega} der Natur der Matrix gemäß bereits beschrieben wurden, werden die Resultate lediglich kurz wiederholt: Bei \textit{mega} sind signifikante Unterschiede bei den Vergleichen mit \textit{ziemlich}, \textit{äußerst}, \textit{sehr}, \textit{überaus}, \textit{voll} und \textit{super} zu beobachten. Das bedeutet, dass sich \textit{mega} mit Ausnahme von \textit{sau} ebenfalls gänzlich von den restlichen Intensivierern unterscheidet.


Zusammengefasst demonstrieren die Ergebnisse der Vergleiche, dass sich vor allem drei Ausdrücke klar von den anderen abheben: der deskriptive Intensivierer \textit{ziemlich} sowie die beiden expressiven Intensivierer \textit{mega} und \textit{voll}. Unterscheidet sich ersterer von jedem einzelnen anderen Ausdruck, verhält sich \textit{mega} anders als sechs und \textit{voll} anders als fünf der restlichen Intensivierer. Diese Befunde werden im nächsten Abschnitt \ref{abs:Diskussion} ausführlich besprochen.





\subsection{Diskussion} \label{abs:Diskussion}
In diesem Abschnitt werden die Ergebnisse von Experiment 2a diskutiert und unter Berücksichtigung der eingangs formulierten {Hypothese H6} evaluiert. Diese geht auf \citet[87]{Biedermann.1969} und {\citet[408]{Breindl.2007}} zurück, die von intraklassischen Unterschieden ausgehen. In diesem Sinne ist anzunehmen, dass sich die Ausdrücke in Abhängigkeit des durch sie ausgedrückten Intensitätsgrades auf einer kontinuierlichen Skala anordnen lassen. Ferner ist möglich, dass die Expressivität der Ausdrücke deren Skalaritätsgrad beeinflusst und letzterer in der Folge als stärker wahrgenommen wird {\citep[vgl.][133]{Gutzmann.2019a}}. Weil dieses Experiment, das dazu diente, den Skalaritätsgrad von Intensivierern mit unterschiedlich wahrnehmbarer Intensität zu ermitteln, und das anschließend vorgestellte {Experiment 2b}, das den Grad an Sprechereinstellung erfragte, als Ganzes betrachtet werden müssen, werden zunächst keine finalen {Schlüsse gezogen}.


In Experiment 2a ging es für die Versuchspersonen darum, auf einer Skala von `normal' bis `maximal' die Merkmalsausprägung eines im Satz beschriebenen Gegenstands anzugeben. Die Fragestellung war im Wortlaut jeweils an den Itemsatz angepasst, bspw. \textit{Wie hoch ist der Berg?}


Als erstes Ergebnis kann festgehalten werden, dass sich die beiden Intensiviererklassen hinsichtlich des ihnen zugeschriebenen Skalaritätsgrades deutlich voneinander unterscheiden (vgl. Tabelle \ref{LinReg01}). So wurde die Skalarität von Sätzen mit expressiven Intensivierern als signifikant stärker eingeschätzt als die Skalarität von Sätzen mit deskriptiven Intensivierern. Wie schon in dem in {Kapitel \ref{Experiment0}} präsentierten Experiment 1 spricht dieser Befund klar für die Gültigkeit der grundlegenden Annahme eines wahrnehmbaren Unterschiedes in der Skalarität von deskriptiven und expressiven Intensivierern (vgl. die {Hypothese H3} nach {\citealt[22]{Paradis.1997}} und {\citealt[133]{Gutzmann.2019a}}). Den Ergebnissen der beiden Untersuchungen zufolge verschieben expressive Intensivierer den Ausprägungsgrad der bezeichneten Eigenschaft auf einer Skala also tatsächlich weiter nach rechts als deskriptive Intensivierer, was als weitere Evidenz für die Hypothese H3 herangezogen {werden kann}.


Da in diesem Experiment vor allem die intraklassischen Abstufungen im Fokus standen, wurden die Intensiviererklassen im nächsten Analyseschritt aufgeschlüsselt und ihre Mitglieder einzeln inspiziert. Dabei trat sowohl bei den deskriptiven als auch den expressiven Intensivierern eine dreiteilige Abstufung im Hinblick auf den mit den Ausdrücken verbundenen Skalaritätsgrad auf, die sich wie folgt äußert: \textit{ziemlich} $\prec$ \textit{äußerst}|\textit{sehr} $\prec$ \textit{überaus} sowie {\textit{voll} $\prec$ \textit{super}|\textit{sau} $\prec$ \textit{mega}}. Demzufolge gibt es nach dem Empfinden der Versuchspersonen Intensivierer, die sich in Bezug auf den durch sie ausgedrückten Skalaritätsgrad stark ähneln, nämlich \textit{äußerst} und \textit{sehr} aufseiten der deskriptiven Intensivierer und \textit{super} und \textit{sau} aufseiten der expressiven Intensivierer. Dass die angenommenen Unterschiede auftreten, erhärtet das Zutreffen der zu testenden Hypothese H6. Die Ergebnisse decken sich in Teilen auch mit der in diesem Kontext dargelegten Rangfolge. Wie eingangs erläutert, verbinden {\citet[87]{Biedermann.1969}} und {\citet[408]{Breindl.2007}} den deskriptiven Intensivierer \textit{ziemlich} mit einem eher gemäßigten Intensitätsgrad. Gleiches zeigt sich in diesem Experiment, schätzten  die Versuchspersonen dessen Skalaritätsgrad doch ebenso vergleichsweise niedrig ein. Den beiden Intensivierern \textit{sehr} und \textit{überaus} sprechen die Autor:innen einen gemäßigten bis hohen bzw. hohen Intensitätsgrad zu, wohingegen sie \textit{äußerst} mit einem extremen Intensitätsgrad assoziiert sehen. In diesem Punkt zeichnen sich bei den durch das Experiment gewonnenen Ergebnissen Unterschiede ab, wurde \textit{äußerst} nicht auf höchster Stufe, sondern zusammen mit \textit{sehr} in der Mitte der Skala positioniert. Am stärksten wahrgenommen haben die Versuchspersonen den Skalaritätsgrad bei \textit{überaus}. Dies ist insofern bemerkenswert, als sich beide Ausdrücke aus semantischer Sicht kaum voneinander unterscheiden. So handelt es sich sowohl bei \textit{äußerst} als auch bei \textit{überaus} um native Ausdrücke, die qua Semantik angeben, dass ein kontextuell determinierter Schwellenwert überschritten wird. Im Gegensatz zu den deskriptiven Intensivierern entsprechen die Ergebnisse der expressiven Intensivierer gänzlich der im Vorfeld getroffenen persönlichen Einschätzung. Bei dieser Klasse bestand die Erwartung, dass das dem Anschein nach abgenutzte \textit{voll} nicht mehr ausdrucksstark genug ist, um als übermäßig skalar empfunden zu werden (vgl. die Resultate zu {Experiment 1}). Diese Annahme hat sich wiederholt bestätigt. Obendrein wurde davon ausgegangen, dass \textit{super} mit einem gemäßigten Skalaritätsgrad einhergeht, während er bei \textit{sau} und \textit{mega} mutmaßlich als am höchsten eingeschätzt wird. Den Ergebnissen nach sind auch diese Vermutungen im Wesentlichen korrekt, spiegelt sich innerhalb dieser Intensiviererklasse doch ausnahmslos die zuvor erwartete Rangfolge wider.


Mit dem Ziel die beobachteten Unterschiede auf Signifikanz hin zu überprüfen, wurden die Intensivierer anschließend paarweise miteinander verglichen (vgl. {Tabelle \ref{Matrix_1_neu}}). Wie vor diesem Hintergrund erläutert, sind dabei vor allem die Ergebnisse des deskriptiven Intensivierers \textit{ziemlich} auffällig. Dieser zeigt signifikante Unterschiede bei den Kontrasten mit allen anderen untersuchten Intensivierern -- unabhängig von der Intensiviererklasse. Schaut man sich die gemittelten Skalenwerte der Intensivierer an (vgl. {Tabelle \ref{Skalenwert3}}), so fällt auf, dass dem Ausdruck von den Versuchspersonen der mit Abstand niedrigste Grad an Skalarität zugeschrieben wurde, was höchstwahrscheinlich durch dessen inhärente Semantik bedingt ist: Anders als die restlichen Intensivierer drückt \textit{ziemlich} semantisch kein Übermaß einer bestimmten Eigenschaft aus, sondern approximiert lediglich einen bestimmten Referenzpunkt. So gibt bspw. \citet[1610]{Pfeifer.1995} dessen Bedeutung als `in verhältnismäßig hohem, reichlichem Maße, annähernd, ungefähr' an. Wie von {\citet[87]{Biedermann.1969}} und \citet[408]{Breindl.2007} korrekterweise gemutmaßt, scheint \textit{ziemlich} daher in der Tat als bloß moderat skalierend wahrgenommen zu werden, was der Hauptgrund für die signifikante Abweichung von den restlichen im Experiment berücksichtigten Intensivierern sein kann.


Ein weiteres auffälliges und somit diskutables Ergebnis ist beim expressiven Intensivierer \textit{mega} auszumachen, der von den Versuchspersonen mit dem höchsten Skalaritätsgrad assoziiert wurde. Bei den Paarvergleichen zeigte sich, dass sich der Ausdruck neben allen deskriptiven Intensivierern auch von den niedrig eingeschätzten expressiven Intensivierern \textit{voll} und \textit{super} signifikant unterscheidet. Auf Grundlage der bisherigen Resultate dieser Untersuchung ist diese Beobachtung kaum überraschend. So deuteten zum einen auch die Befunde der in {Kapitel \ref{Korpus}} präsentierten Korpusanalyse auf die Neuheit und anhaltende Beliebtheit dieses Intensivierers hin. Zum anderen deckt sie sich mit dem Ergebnis, das im vorangegangenen Experiment bei {Aufgabe I} auftrat: Auch hier war \textit{mega} der von den Versuchspersonen bezüglich Skalarität am höchsten beurteilte Intensivierer. Diese Ergebnisse bieten in Summe Anhaltspunkte dafür, dass \textit{mega} grundsätzlich mit einem hohen Skalaritätsgrad in Verbindung gebracht wird. Interessanterweise stimmt dies mit {\citet[]{Calpestrati.2017}} überein, der in seiner Korpusanalyse zu einem ähnlichen Schluss kommt. Demnach hat \textit{mega} anscheinend auch ein halbes Jahrzehnt später nichts von seiner wahrnehmbaren Intensität {eingebüßt}.


Ebenso bemerkenswert ist das Ergebnis von \textit{voll}, dessen Skalaritätsgrad innerhalb der expressiven Intensiviererklasse von den Versuchspersonen mit Abstand am niedrigsten empfunden wurde. Aus diesem Grund unterscheidet sich der Ausdruck bei den paarweisen Vergleichen signifikant von allen expressiven Intensivierern, von dem am höchsten eingeschätzten deskriptiven Intensivierer \textit{überaus} sowie dem am niedrigsten eingeschätzten deskriptiven Intensivierer \textit{ziemlich}. Eine mögliche Erklärung hierfür bietet die Überlegung, dass sich \textit{voll} unterdessen pragmatisch abgeschliffen hat und deswegen prinzipiell als nicht mehr stark skalierend {wahrgenommen wird}.


Auch der hinsichtlich Skalarität am zweithöchsten evaluierte expressive Intensivierer \textit{sau} unterscheidet sich signifikant von vier der restlichen Ausdrücke: den vergleichsweise niedrig eingeschätzten deskriptiven Intensivierern \textit{ziemlich}, \textit{äußerst} und \textit{überaus} sowie dem am niedrigsten bewerteten expressiven Intensivierer \textit{voll}. Dass \textit{sau} mit einem dem Anschein nach hohen Skalaritätsgrad assoziiert zu werden scheint, ist eventuell dessen klar erkennbarer Semantik geschuldet, die in Kombination mit der damit verbundenen derben Konnotation zur hohen Gradeinschätzung beigetragen haben kann.


Bei den Paarvergleichen war weiterhin festzustellen, dass sich \textit{äußerst} und \textit{sehr} parallel verhalten, die aufgrund ihres als beinahe identisch wahrgenommenen Skalaritätsgrades signifikante Unterschiede bei drei der restlichen Ausdrücke zeigen: dem am niedrigsten verorteten deskriptiven Intensivierer \textit{ziemlich} sowie den beiden als hoch eingeschätzten expressiven Intensivierern \textit{sau} und \textit{mega}. Im Gegensatz dazu unterscheidet sich der expressive Intensivierer \textit{super} ebenso signifikant von drei der anderen Ausdrücke: dem deskriptiven Intensivierer \textit{ziemlich}, dem am niedrigsten bewerteten expressiven Intensivierer \textit{voll} sowie dem am höchsten bewerteten expressiven Intensivierer \textit{mega}. Daraus lässt sich schließen, dass \textit{super} und \textit{sau} trotz der nahezu identischen Skaleneinschätzung bei den Paarvergleichen nicht die gleichen signifikanten Unterschiede aufweisen.


Kurz gefasst liefern die Ergebnisse von Experiment 2a erste stichhaltige Hinweise auf die Gültigkeit der zu testenden Hypothese H6 nach \citet[87]{Biedermann.1969} und \citet[408]{Breindl.2007}, finden sich bei der Untersuchung von Skalarität doch die vermuteten klasseninternen Abstufungen. Dies deutet darauf hin, dass sich die Ausdrücke im Hinblick auf den durch sie ausgedrückten Skalaritätsgrad de facto wahrnehmbar voneinander unterscheiden. Nachfolgend wird das parallel aufgebaute {Experiment 2b} vorgestellt. Dieses wurde separat durchgeführt und zielte darauf ab, den durch die Intensivierer ausgedrückten Grad an Sprechereinstellung zu ermitteln.

%Dies dient der Überprüfung, ob die im Kontext von Skalarität beobachtete Rangfolge in gleicher Weise bei Expressivität zutreffend ist und ob die Merkmale Skalarität und Sprechereinstellung im Sinne von {\citet[133]{Gutzmann.2019a}} miteinander korrelieren.
%Dabei ist anzunehmen, dass die Expressivität der Ausdrücke, wie von \citet[133]{Gutzmann.2019a} behauptet, Einfluss auf deren Skalaritätsgrad hatte. So drückten die expressiven Intensivierer relativ zu den deskriptiven Intensivierern auf jeder Stufe einen höheren Grad an Skalarität aus. 

%In diesem Sinne waren bei beiden Intensiviererklassen graduelle Abstufungen in Bezug auf den mit den Ausdrücken einhergehenden Skalaritätsgrad auszumachen: schwach $\prec$ mittel $\prec$ stark. Demzufolge können deskriptive Intensivierer entweder schwach, mittel oder stark skalierend sein, ebenso wie auch die expressiven Intensivierer als schwach, mittel oder stark skalierend wahrgenommen werden können. Hierbei sind die Ausdrücke in Abhängigkeit ihrer Intensiviererklasse auf zwei unterschiedlichen Skalen nebeneinander geordnet, die nach dem bspw. von {\citet[87]{Biedermann.1969}} und {\citet[408]{Breindl.2007}} vermuteten Intensitätsgrad abgestuft sind. Diese Idee wird in Tabelle \ref{schwach} genauer vorgeführt. Die Anordnung der Intensiviererklassen erfolgt in ansteigendem Skalenwert.



%%
%\begin{table}[h]
%
%%\setlength{\tabcolsep}{8.25pt}
%%\begin{tabularx}{\textwidth}{p{0.15\textwidth} p{0.1\textwidth} p{0.15\textwidth} p{0.1\textwidth} p{0.15\textwidth}}
%\begin{tabularx}{.8\textwidth}{Xlclc}
%\lsptoprule
% & & \multicolumn{2}{c}{\textbf{Intensiviererklasse}}  \\ [0.5em]
%\multirow{2}{*}{Stärkegrad} & \multicolumn{2}{c}{\textsc{deskriptiv}} & \multicolumn{2}{c}{\textsc{expressiv}}  \\ 
% & Intensivierer & Skalenwert & Intensivierer & Skalenwert  \\ 
%  &  & (\textit{SD})  &  & (\textit{SD})  \\
%\midrule 
%\\
%schwach & \textit{ziemlich} & 51,22 & \textit{voll} & 56,38  \\
% & & (23,6) & & (23,4)  \\
%\multicolumn{1}{l}{mittel} & \textit{äußerst} & 59,46 & \textit{super} & 63,15 \\
%& & (23,4) &  & (22,2)\\
% & \textit{sehr} & 59,69 & \textit{sau} & 63,87\\
%  & & (23,3) & & (24,4)\\
%stark & \textit{überaus} & 62,36 & \textit{mega} & 67,27  \\
% & & (22,0) & & (24,4)  \\
%\end{tabularx}
%\caption[Klasseninterne graduelle Abstufungen bei Experiment 2a]{Klasseninterne graduelle Abstufungen bei Experiment 2a -- Skalarität (in ansteigendem Skalenwert).}
%\label{schwach}
%\end{table}









%Nachfolgend wird das parallel aufgebaute {Experiment 2b} vorgestellt. Dieses wurde separat durchgeführt und zielte darauf ab, den durch die Intensivierer ausgedrückten Grad an Sprechereinstellung zu ermitteln.
%%Dies dient der Überprüfung, ob die im Kontext von Skalarität beobachtete Rangfolge in gleicher Weise bei Expressivität zutreffend ist und ob die Merkmale Skalarität und Sprechereinstellung im Sinne von {\citet[133]{Gutzmann.2019a}} miteinander korrelieren.



%So lieferten die im Rahmen der Korpusuntersuchung ermittelten diachronen Häufigkeitsverteilungen Hinweise darauf, dass der Ausdruck schon in den Sechzigerjahren vermehrt in intensivierender Funktion verwendet wurde. Daher ist davon auszugehen, dass er sich in der gesprochenen Sprache noch früher etabliert hat, was offenbar abträglich für die mit ihm einhergehende Intensität war. Nichtsdestotrotz stellt sich an der Stelle die Frage, aus welchen Gründen dem im Korpus ähnlich frequenten \textit{sau} von den Versuchspersonen der zweithöchste Skalaritätsgrad zugeschrieben wurde. Diesbezüglich lässt sich spekulieren, dass dies dessen klar erkennbarer Semantik geschuldet war, die in Kombination mit der damit verbundenen derben Konnotation zum hohen Skalenwert beigetragen haben kann. Ob diese Vermutung der Realität entspricht, lässt sich allerdings nicht eindeutig beantworten, zumal der Skalaritätsgrad von \textit{sau} im vorherigen Experiment nur unwesentlich höher eingestuft wurde als der von \textit{voll}.


%Bei den paarweisen Vergleichen der deskriptiven Intensivierer war festzustellen, dass sich \textit{äußerst} und \textit{sehr} parallel verhielten, die aufgrund beinahe identischer Skalenwerte (\textit{äußerst} = 59,46 Punkte und \textit{sehr} = 59,69 Punkte) gleichermaßen signifikante Unterschiede bei den Kontrasten mit \textit{ziemlich}, \textit{mega} und \textit{sau} anzeigten. Die möglichen Gründe hierfür wurden oben im Kontext der dreiteiligen Abstufung der Intensiviererklassen erörtert. Vor diesem Hintergrund wurde erklärt, dass \textit{äußerst} mit 19\,723 absoluten Treffern im Korpus nahezu doppelt so häufig zu finden ist wie \textit{überaus}, das darin nur 10\,534 absolute Treffer aufweist. Der vergleichsweise niedrige Skalenwert von \textit{überaus} (62,36 Punkte) war vermutlich nicht semantisch bedingt, unterschied sich dieser Ausdruck dahingehend doch kaum von anderen deskriptiven Intensivierern wie z.\,B. \textit{äußerst}. Eine bessere Erklärung bietet dagegen die Gebrauchsfrequenz: Da \textit{überaus} erwiesenermaßen seltener verwendet wird als die anderen deskriptiven Intensivierer, wurde er von den Versuchspersonen mutmaßlich als markierter wahrgenommen und in der Folge als stärker skalierend beurteilt.


%Zusammengenommen deuteten auch die Ergebnisse der Paarvergleiche auf einen deutlichen Unterschied zwischen den Intensiviererklassen hin.




%Dabei zeigte sich, dass die mittleren Werte der deskriptiven Intensivierer auch im Detail überwiegend unter denen der expressiven Intensivierer lagen. Dennoch war auf beiden Seiten eine Ausnahme auszumachen, nämlich der deskriptive Intensivierer \textit{überaus} und der expressive Intensivierer \textit{voll}. So wurde \textit{überaus} deutlich über den restlichen deskriptiven Intensivierern, \textit{voll} dagegen deutlich unter den restlichen expressiven Intensivierern eingestuft. Ob und inwiefern es sich bei diesen Intensivierern um Ausreißer handelt, hängt im Wesentlichen davon ab, wie eng man die oben skizzierte Hypothese H3 fasst, sind diesbezüglich doch zwei Interpretationsweisen denkbar: eine starke und eine schwache Interpretation. Gemäß ersterer, bei der sich deskriptive und expressive Intensivierer auf der gleichen Skala befinden, ist anzunehmen, dass der schwächste expressive Intensivierer im Hinblick auf Skalarität grundsätzlich höher positioniert wird als der stärkste deskriptive Intensivierer: {deskriptive Intensivierer $\prec$ expressive Intensivierer}. Dementsprechend gibt es Abstufungen in den Stärkegraden, die unabhängig von der Intensiviererklasse sind. Entscheidend ist alleine, dass die expressiven Intensivierer eine höhere Position auf der Skala einnehmen als die deskriptiven Intensivierer.
%Nach dieser Interpretation würde es sich bei \textit{überaus} und \textit{voll} tatsächlich um Ausreißer handeln. 
%In dem Fall ergibt sich die Frage, wieso sich die beiden Intensivierer anders verhielten als die restlichen Klassenmitglieder. Im Kontext von \textit{überaus} bietet die Annahme, dass es als Intensivierer vergleichsweise seltener -- und möglicherweise zu selten -- ist, eine mögliche Erklärung für vergleichsweise hoch eingeschätzten Skalaritätsgrad. Um diese Annahme zu überprüfen, wurde das absolute Vorkommen der deskriptiven Intensivierer im Nachgang anhand des gleichen Korpus, das im Fokus der in Kapitel \ref{Korpus} vorgestellten Korpusuntersuchung stand, ermittelt. Dabei stellte sich heraus, dass \textit{überaus} mit 10\,523 absoluten Treffern die niedrigste Belegdichte der deskriptiven Intensiviererklasse aufweist.\footnote{Die Belegtreffer der restlichen deskriptiven Intensivierer sehen wie folgt aus: \textit{äußerst}: 19\,723 absolute Treffer, \textit{ziemlich}: 54\,379 absolute Treffer und \textit{sehr}: 165\,161 absolute Treffer.} Bei diesem Ausdruck kann die niedrigere Frequenz also dazu geführt haben, dass er von den Versuchspersonen als markierter wahrgenommen wurde als die restlichen deskriptiven Intensivierer, was sich letztlich in dem ihm zugesprochenen Skalaritätsgrad widergespiegelt hat. Der expressive Intensivierer \textit{voll} ist hingegen eventuell schlichtweg zu abgenutzt. So konnten im Zeitschriftenkorpus schon um 1960 herum erste Belege für sein vermehrtes Vorkommen gefunden werden, was zeigt, dass es \textit{voll} schon lange Zeit in intensivierender Funktion gibt. Auffällig ist, dass dieses Ergebnis dem Befund entspricht, der sich im vorangegangenen Experiment 1 bei Aufgabe I zeigte: Auch hier wurde \textit{voll} von den Versuchspersonen mit einem verhältnismäßig niedrigen Skalaritätsgrad assoziiert ({vgl. {Tabelle \ref{SkaWe2}}}). Dies erhärtet den Verdacht, dass das intensivierende \textit{voll} inzwischen verbraucht ist, weswegen ihm grundsätzlich eine niedrigere Skalenposition zugesprochen wird als anderen expressiven Intensivierern.


%Neben der starken Interpretationsart der Hypothese H3 ist allerdings auch eine gemäßigtere vorstellbar. Nach dieser gibt es bei beiden Intensiviererklassen graduelle Abstufungen in Bezug auf den mit den Ausdrücken einhergehenden Skalaritätsgrad: schwach $\prec$ mittel $\prec$ stark. Dieser Einteilung nach können deskriptive Intensivierer entweder schwach, mittel oder stark skalierend sein, ebenso wie auch die expressiven Intensivierer als schwach, mittel oder stark skalierend wahrgenommen werden können. Hierbei sind die Intensiviererklassen auf zwei unterschiedlichen Skalen nebeneinander geordnet, die nach der Stärke der Skalierung abgestuft sind. Dadurch, dass die Grenzen der Stärkegrade parallel verschiebbar sind und sich die Abstufungen nur innerhalb einer Klasse abzeichnen, gibt es anders als bei der soeben umrissenen starken Interpretation keine Ausreißer, sondern Überlappungen zwischen den Intensiviererklassen. Daher würde es sich bei \textit{überaus} und \textit{voll} um Überlappungen handeln, die auf das Verschieben von zwei parallelen Skalen zurückgehen. Diese Idee wird in Tabelle \ref{schwach} genauer vorgeführt. Die Anordnung der Intensiviererklassen erfolgt in ansteigendem Skalenwert.



%
%\begin{table}[h]
%
%%\setlength{\tabcolsep}{8.25pt}
%%\begin{tabular}{p{0.15\textwidth} p{0.1\textwidth} p{0.15\textwidth} p{0.1\textwidth} p{0.15\textwidth}}
%\begin{tabularx}{.8\textwidth}{Xlclc}
%\lsptoprule
% & & \multicolumn{2}{c}{\textbf{Intensiviererklasse}}  \\ [0.5em]
%\multirow{2}{*}{Stärkegrad} & \multicolumn{2}{c}{\textsc{deskriptiv}} & \multicolumn{2}{c}{\textsc{expressiv}}  \\ 
% & Intensivierer & Skalenwert & Intensivierer & Skalenwert  \\ 
%  &  & (\textit{SD})  &  & (\textit{SD})  \\
%\midrule 
%\\
%schwach & \textit{ziemlich} & 51,22 & \textit{voll} & 56,38  \\
% & & (23,6) & & (23,4)  \\
%\multicolumn{1}{l}{mittel} & \textit{äußerst} & 59,46 & \textit{super} & 63,15 \\
%& & (23,4) &  & (22,2)\\
% & \textit{sehr} & 59,69 & \textit{sau} & 63,87\\
%  & & (23,3) & & (24,4)\\
%stark & \textit{überaus} & 62,36 & \textit{mega} & 67,27  \\
% & & (22,0) & & (24,4)  \\
%\end{tabularx}
%\caption[Schwache Interpretationsart als Erklärung für die Ausreißer von Experiment 2a]{Schwache Interpretationsart als Erklärung für die potenziellen Ausreißer von Experiment 2a -- Skalarität (in ansteigendem Skalenwert).}
%\label{schwach}
%\end{table}

%Abgesehen von diesem Befund konnte ein weiteres Ergebnis beobachtet werden, das es zu diskutieren gilt. So  deutete sich nämlich sowohl innerhalb der deskriptiven als auch der expressiven Intensiviererklasse eine dreiteilige Abstufung im Hinblick auf den gemittelten Skalenwert an, die sich wie folgt ausdrückte: \textit{ziemlich} $\prec$ \textit{äußerst}|\textit{sehr} $\prec$ \textit{überaus} sowie {\textit{voll} $\prec$ \textit{super}|\textit{sau} $\prec$ \textit{mega}}. Somit schien es nach Einschätzung der Versuchspersonen Intensivierer zu geben, die sich in Bezug auf den durch sie ausgedrückten Skalaritätsgrad stark ähnelten, nämlich \textit{äußert} (59,46 Punkte) und {\textit{sehr} (59,69 Punkte}) aufseiten der deskriptiven Intensivierer ($\Delta$ = ca. 0,23 Punkte) und {\textit{super} (63,15 Punkte)} und \textit{sau} (63,87 Punkte) aufseiten der expressiven Intensivierer {($\Delta$ = ca. 0,72 Punkte)}. Semantisch könnten die beiden Intensiviererpaare nicht unterschiedlicher sein: So handelt es sich bei \textit{äußerst} bspw. um einen transparenten Ausdruck, dessen Ursprungslexem deutlich wahrnehmbar ist. Anders als bei \textit{sehr}: Wie in Abschnitt \ref{Intensivierertypen} erörtert, ist der lexikalische Ursprung von \textit{sehr} unterdessen vollkommen opak. Dessen ungeachtet ist anzunehmen, dass sich die beiden Ausdrücke hinsichtlich ihres Lexikalisierungsgrades nicht sonderlich stark voneinander unterscheiden: Während es sich bei letzterem um den klassischen Vertreter der deskriptiven Intensiviererklasse handelt, der sowohl im schriftlichen wie auch im mündlichen Sprachgebrauch entsprechend frequent ist, kann \textit{äußerst} jedoch als ebenso stark lexikalisiert betrachtet werden. Diese Vermutung lässt sich auf Grundlage der erhobenen Korpusfrequenzen belegen. Im Korpus ist \textit{äußerst} mit 19\,723 absoluten Treffern bspw. erheblich frequenter als \textit{überaus}, das darin lediglich 10\,534 absolute Treffer aufweist. Demnach ist es möglich, dass der Lexikalisierungsgrad von \textit{äußerst} im Hinblick auf Skalarität zu einer vergleichbaren Einschätzung geführt hat wie die von \textit{sehr}. Ähnlich verhält es sich bei \textit{super} und \textit{sau}. Deren vergleichbarer Skalenwert ist insofern überraschend, als beide Ausdrücke ebenfalls über kaum semantische Gemeinsamkeiten verfügen: So handelt es sich bei \textit{super} um einen aus dem Griechischen entlehnten Ausdruck, der \textit{qua} Semantik mit einer dimensionalen Größe einhergeht. Dieser war mit ingesamt 391,26 Treffern pro Million Textwörter der am häufigsten verwendete Intensivierer im Korpus. Dagegen ist \textit{sau} ein nativer Ausdruck, dessen Ursprung im Tierreich zu verorten ist und der möglicherweise aufgrund seiner heutigen derben Konnotation den Weg in die Intensiviererrolle gefunden hat \citep[vgl.][360]{Claudi.2006}. Dieser Ausdruck lässt sich im Korpus mit 60,41 Treffern pro Million Textwörter nachweisen. Was die beiden Ausdrücke trotz dieser Unterschiede allerdings miteinander gemein haben, ist ihr persistenter Gebrauch, wie sich anhand der in {Abbildung \ref{SyDia}} dargestellten Histogramme herausstellte: So fanden sich sowohl für \textit{super} als auch für \textit{sau} schon zu Beginn des Untersuchungszeitraums erste Korpustreffer, was darauf hindeutet, dass sich die Intensivierer in der gesprochenen Sprache frühzeitig etabliert haben. Dieser persistente Gebrauch hat gegebenenfalls dazu geführt, dass die beiden Ausdrücke in der Zwischenzeit mit einem ähnlichen Skalaritätsgrad assoziiert werden und von den Versuchspersonen in der Konsequenz auf der Skala vergleichbar hoch {eingestuft wurden}.











%\subsubsection{Regressionsanalyse} $~$ \\
%Das in Tabelle \ref{LinReg1} präsentierte lineare gemischte Modell zeigte, dass die Prädiktoren \textsc{Intensiviererklasse} und \textsc{Position} sowie deren Interaktion signifikanten Einfluss auf den Skalenwert hatten. Demzufolge wurden expressive Intensivierer in Übereinstimmung mit der Hypothese H4 auf der Skala signifikant höher positioniert als deskriptive Intensivierer. Hinzu kam ein signifikanter Positionseffekt: Je weiter fortgeschritten die Position eines Items im Experiment war, desto niedriger wurde es im Hinblick auf Skalarität auf der Skala beurteilt. Die ebenfalls signifikante Zwei-Wege-Interaktion zwischen den beiden Variablen bedeutete außerdem, dass sich die Intensiviererklassen mit fortschreitender Position einander angenähert und der Unterschied zwischen deskriptiven und expressiven Intensivierern folglich kleiner geworden ist. Dies legt einen Gewöhnungseffekt nahe, der dem uniformen Untersuchungsdesign geschuldet sein kann. Somit besteht die Möglichkeit, dass sich die Versuchspersonen im Laufe des Experiments an die Intensivierer gewöhnt haben, weswegen sich der zu Beginn als stärker empfundene Unterschied zwischen den Intensiviererklassen sukzessive verringert hat.  
%Zusammengefasst deuten die Ergebnisse des Regressionsmodells \textit{de facto} auf einen signifikanten Unterschied zwischen beiden Intensiviererklassen hin, der sich durch eine höhere Skalenposition der expressiven Intensivierer ausdrückt, was in Einklang mit der zu testenden Hypothese steht.



%-	Die Wahrscheinlichkeit, dass XY dem Zufall geschuldet ist, ist sehr klein, da die Korrelation hochsignifikant ist
%-	Außerdem finden wir bei XY eine starke Korrelation
%-	XY ist XY und dies ist nicht zufallsbedingt 
%-	Mit einem Mean von





%%%%%%%%%%%%%%%%%%%%%%%%%%%%%%%%%%%%%%%%%%%%%%%%%%%%%%%%%%%%%%%%%%%%%%%%%%%%%%%%%%%%%%%%%%%%%%%%%%%%%%%%%%%%%%%
%%%%%%%%%%%%%%%%%%%%%%%%%%%%%%%%%%%%% EXPERIMENT II %%%%%%%%%%%%%%%%%%%%%%%%%%%%%%%%%%%%%%%%%%%%%%%%%%%%%%%%%%%
%%%%%%%%%%%%%%%%%%%%%%%%%%%%%%%%%%%%%%%%%%%%%%%%%%%%%%%%%%%%%%%%%%%%%%%%%%%%%%%%%%%%%%%%%%%%%%%%%%%%%%%%%%%%%%%





\section{Experiment 2b -- Sprechereinstellung}\label{abs:Sprechereinstellung}
In diesem Teilkapitel steht das zweite Experiment im Fokus. Wie {Experiment 2a} wurde auch dieses aus der Intention heraus umgesetzt, die in Abschnitt \ref{abs:Hypothese und Forschungsaufbau} beschriebene Hypothese H6 auf ihre Korrektheit hin zu überprüfen. In diesem Kontext ist davon auszugehen, dass sich nicht nur zwischen, sondern auch innerhalb der Intensiviererklassen graduelle Abstufungen beobachten lassen, die auf die Semantik der Ausdrücke zurückzuführen sind. Im Vorfeld dieser Untersuchung wird erwartet, dass die expressiven Intensivierer jeder Stufe von den Versuchspersonen als stärker bewertend beurteilt werden als die deskriptiven Intensivierer. Dabei stellt sich zudem die Frage, ob die in Experiment 2a beobachteten Abstufungen in gleicher Weise bei Sprechereinstellung zutage treten, die Intensivierer also in Bezug auf dieses Merkmal eine ähnliche Rangfolge wie bei Skalarität aufweisen. Diesbezüglich ist im Sinne von \citet[133]{Gutzmann.2019a} denkbar, dass die Expressivität der Ausdrücke deren Skalaritätsgrad beeinflusst, letzterer infolge der Expressivität als stärker wahrgenommen wird und es einen korrelativen Zusammenhang zwischen beiden Merkmalen gibt.





\subsection{Untersuchungsgegenstand}\label{abs:Untersuchungsgegenstand2509}
Bei dem in Experiment 2b verwendeten Testmaterial handelte es sich um die gleichen Items wie in Experiment 2a, in denen feste Kombinationen aus verschiedenen Nomen (n = 56) und skalaren Adjektiven (n = 7) mit jeweils vier deskriptiven und expressiven Intensivierern gekreuzt wurden: \textit{äußerst}, \textit{sehr}, \textit{überaus} und \textit{ziemlich} (deskriptiv) sowie \textit{mega}, \textit{sau}, \textit{super} und \textit{voll} (expressiv). Durch den parallelen Untersuchungsaufbau konnte maximale Vergleichbarkeit gewährleistet und die Ergebnisse der beiden Experimente in Bezug auf eine mögliche Korrelation hin überprüft werden. Zur Erinnerung werden die in (\ref{itemse1}) beispielhaft gelisteten Items erneut aufgeführt:

\ea
\ea Die Pfütze ist äußerst groß.
\ex Der Flur ist sehr lang.
\ex Die Mauer ist überaus hoch.
\ex Das Gemälde ist ziemlich alt.
\ex Die Straße ist mega breit.
\ex Das Wasser ist sau heiß.
\ex Der Stamm ist super dick.
\ex Das Schiff ist voll groß.
\z
\z
Somit wurde auch in diesem Experiment eine unabhängige Variable mit 8 Stufen untersucht: 4 deskriptive und 4 expressive Intensivierer. Die Kombination aus Nomen und Adjektiv wurde erneut nicht variiert. Stattdessen rotierte jede feststehende Nomen-Adjektiv-Kombination (n = 7) im Experiment über die Intensivierer, wodurch für jeden Intensivierer gleich viele Datenpunkte gesammelt werden konnten. Wie zuvor wurden die abzufragenden Items nach dem lateinischen Quadrat auf acht Versuchspersonengruppen verteilt, sodass jede Person eine dieser acht Listen mit insgesamt 56 Testsätzen mit je 28 deskriptiven und expressiven Items zu evaluieren hatte. Auch am Randomisierungsverfahren änderte sich nichts: Wiederholt wurde die Abfolge der Items nach dem Adjektiv pseudorandomisiert, um zu verhindern, dass das gleiche Adjektiv im Experiment direkt hintereinander folgt, was sich gegebenenfalls unerwünscht auf die Skaleneinschätzung ausgewirkt hätte.


Als Ergänzung zu den für die Auswertung relevanten Items wurden die zuvor konstruierten Übungssätze in das Experiment miteinbezogen. Diese verfügten über die gleiche syntaktische Struktur wie die Items, wiesen jedoch andere Intensivierer, Nomen und Adjektive auf. Die Übungssätze dienten im Wesentlichen dazu, die Versuchspersonen mit der Aufgabe, der Skala und der Bedienung des Schiebereglers vertraut zu machen. Wie schon in Experiment 2a wurden auch in diesem Experiment keine Distraktoren verwendet, da sie für die Zielsetzung der Untersuchung als überflüssig erachtet wurden.





%Neben der neuen Aufgabenstellung, die im nächsten Abschnitt vorgestellt wird, wurden auch die vier Trainingssätze modifiziert. Wie zuvor dienten diese im Wesentlichen dazu, die Versuchspersonen mit der Aufgabe und der Bedienung des Schiebereglers vertraut zu machen.
%, ohne dies zu explizieren. 
%Die neuen Übungssätze waren wie folgt konzipiert:
%
%\ex. \label{Übung2}
%\a. Das Fenster ist recht groß.
%\b. Das Chili ist extrem scharf.
%\b. Die Rolltreppe ist ganz steil.
%\b. Der Kühlschrank ist übelst neu.



\subsection{Untersuchungsdurchführung}\label{abs:Untersuchungsdurchführung2541}
Erneut wurden den Versuchspersonen am Computer kurze und syntaktisch einfach konstruierte Kopulasätze präsentiert, in denen verschiedene Objekte beschrieben wurden, wie bspw. \textit{Der Berg ist äußerst hoch}. Daraufhin wurden sie dazu aufgefordert, mit einem Schieberegler auf einer kontinuierlichen Skala anzugeben, wie bewertend die jeweilige Aussage ist (vgl. {Experiment 2a}: Wo lässt sich das im Satz beschriebene Objekt im Vergleich zu einem normal ausgeprägten Objekt auf der Skala verorten?). Um den Grad der ausgedrückten Sprechereinstellung zu ermitteln, wurde in {Experiment 2b} also die Aufgabenanweisung an die neue Zielsetzung angepasst. Das bedeutet, dass sich der Wortlaut der Fragestellung sowohl von der in {Experiment 2a} untersuchten Skalarität als auch von der in {Experiment 1} in Aufgabe II erfragten Sprechereinstellung unterschied. Während der Grund für die unterschiedlichen Fragestellungen zwischen den {Experimenten 2a und 2b} auf das jeweils zu untersuchende Merkmal, d.\,h. Skalarität und Sprechereinstellung, zurückgeht, ist der Unterschied zwischen Experiment 2b und {Experiment 1 -- Aufgabe II} durch die zeitliche Abfolge der Experimente begründet. So wurde letzteres erst nach der in diesem Kapitel dargelegten Untersuchung durchgeführt und konnte in Bezug auf die Fragestellung optimiert werden. Abbildung \ref{Bewertungsskala2} illustriert die in Experiment 2b verwendete Skala.

%
\begin{figure}[h] 

  \includegraphics[width=1\textwidth]{Skala2_2}
 \caption[Darstellung der in Experiment 2b verwendeten Bewertungsskala zu Sprechereinstellung] {Darstellung der in Experiment 2b verwendeten Bewertungsskala.} 
  \label{Bewertungsskala2}
\end{figure}
%
%
%Da es in Experiment II jedoch nicht mehr um die Skalarität der Ausdrücke ging, sondern der Grad der Bewertung der Aussage erfragt werden sollte, wurde die Aufgabenformulierung an die neue Fragestellung angepasst. Der wesentliche Unterschied bei der zweiten Untersuchung bestand also darin, dass die Versuchspersonen auf einer Bewertungsskala einschätzen sollten, wie bewertend die jeweilige Aussage ist (vgl. Experiment I: Wo lässt sich das im Satz beschriebene Objekt im Vergleich zu einem vergleichbaren normalen Objekt auf der Skala verorten?). Der Leitgedanke einer solchen subjektiven Bewertung wurde den Versuchspersonen mithilfe von zwei Beispielsätzen folgendermaßen verdeutlicht:

%\begin{singlespace*}
%\begin{quotation} \glqq{}Im Folgenden werden Sie kurze Sätze lesen, wie zum Beispiel \textit{Das Auto steht im Hof} oder \textit{Ich habe die Klausur verkackt}. Ihre Aufgabe besteht darin, auf einer Skala einzuschätzen, wie bewertend die Aussage ist. Unter ‚bewertend‘ verstehen wir, dass der Sprecher den Sachverhalt subjektiv beurteilt. Während beide Beispielsätze einen Sachverhalt beschreiben (dass das Auto im Hof steht bzw. dass der Sprecher bei der Klausur schlecht abgeschnitten hat), ist der zweite Satz darüber hinaus bewertend, da der Sprecher zusätzlich seine Verärgerung über den Sachverhalt ausdrückt.\grqq{}\end{quotation}
%\end{singlespace*}

\noindent Wie Abbildung \ref{Bewertungsskala2} zeigt, war die Skala in Experiment 2b am linken Rand mit `gar nicht bewertend' und am rechten Rand mit `maximal bewertend' gelabelt (im Vergleich dazu `normal' bzw. `maximal' in {Experiment 2a}). Durch Betätigen des Reglers mit der Maus färbte sich die bis dahin farblose {Skala} vom Nullpunkt bis zum gewählten Punkt hin ein, wodurch die zugeschriebene Intensität für die Versuchspersonen visuell hervorgehoben wurde. Über der Skala befand sich der zu beurteilende Satz (bspw. \textit{Der Berg ist äußerst hoch}), darunter die Aufgabenstellung: \textit{Geben Sie mit dem {Schieberegler} an, \textbf{wie bewertend} die Aussage ist}. Um die Aufgabe zu erfüllen, war der Schieberegler durch Klicken oder Ziehen mit der Maus an diejenige Skalenposition zu bewegen, die der persönlichen Einschätzung nach den im Satz ausgedrückten Grad an Sprechereinstellung am besten wiedergibt. Kurz gefasst: Erfragte die Skala in Experiment 2a den Grad an Skalarität, ermittelte die in diesem Experiment verwendete Skala den Grad an {Sprechereinstellung}.


Auch dieses Mal wurden die Versuchspersonen in einem ausführlichen Einleitungstext über den Anlass, die Dauer und den Ablauf des Experiments inklusive datenschutzrechtlicher Aspekte informiert. Zudem wurden sie darum gebeten, das Experiment an einem Endgerät mit physischer Tastatur zu bearbeiten, wodurch verhindert werden sollte, dass die Skalen aufgrund unterschiedlicher Bildschirmgröße (Smartphone/Tablet vs. Notebook/Desktop-PC) nicht korrekt angezeigt oder nicht wie vorgesehen bedient werden können.\footnote{Wie schon in Experiment 2a wurde zur Kontrolle die Bildschirmgröße der Versuchspersonen erfasst.} Im Anschluss an die Instruktion startete das Experiment mit den Übungssätzen, die nach der Beantwortung nahtlos und ohne weitere Erklärung in die für die Auswertung bedeutsamen Items {übergingen}.



\subsection{Versuchspersonen}\label{abs:Probanden2}
Am Experiment haben 63 deutsche Muttersprachler:innen im Alter von 18 bis 40 Jahren teilgenommen, die über \textit{Clickworker} angeworben und (bemessen an dem zu der Zeit gültigen deutschen Mindestlohn) mit jeweils {1,70 Euro} {entlohnt} wurden. Die Versuchspersonen erhielten keine Informationen über das Ziel der Untersuchung. Wie zuvor wurde das Experiment mittels {\textit{LimeSurvey}} erstellt und von den Versuchspersonen individuell online durchgeführt. Die Durchführung war vollkommen anony, ohne dass personenbezogene Daten erhoben wurden. Die Dauer der Bearbeitung betrug im Schnitt etwa zehn Minuten.


Auch bei diesem Experiment wurde die Teilnahme von 64 Versuchspersonen angestrebt, um zu gewährleisten, dass von jedem Intensivierer sieben Datenpunkte und somit eine gleichmäßige Verteilung vorliegt. Erneut kam es nach Abschluss des Experiments zu einem Problem von \textit{Clickworker}, infolge dessen nur Daten von 63 Personen exportiert werden konnten. Dadurch lagen letztlich fünf Listen mit Daten von je acht Versuchspersonen, zwei Listen mit Daten von je sieben Versuchspersonen und eine Liste mit Daten von neun Versuchspersonen vor, die bei der Auswertung berücksichtigt wurden.



\subsection{Ausschlusskriterien}\label{abs:Auswertungskriterien2}
Von den 63 Versuchspersonen, deren Antworten exportiert wurden, flossen 62 in die Auswertung mit ein. Die Daten einer Versuchsperson wurden auf Basis von im Vorfeld der Untersuchung definierten Ausschlusskriterien von der Analyse ausgenommen, weil sie die Vermutung nahelegten, dass die Person die Aufgabe nicht aufmerksam bearbeitet, sondern sich lediglich durch das Experiment durchgeklickt hat. Dies äußerte sich in uniformen Angaben ohne Variation in der Skaleneinschätzung. Dieser Verdacht wurde durch eine Begutachtung der durchschnittlichen Reaktionszeit je Item bestätigt: Während die Reaktionszeit aller Versuchspersonen im Schnitt bei etwa 5,76 Sekunden lag {(\textit{SD} = 37,3)}, betrug die der ausgeklammerten Person dagegen bloß 1,38 Sekunden ({\textit{SD} = 0,9}). Die verhältnismäßig niedrige Reaktionszeit und das einheitliche Antwortmuster führten letztlich zu einem Ausschluss dieser Versuchsperson, sodass in der Auswertung bloß repräsentative Ergebnisse {berücksichtigt wurden}.




\subsection{Ergebnisse}\label{abs:Ergebnisse2576}
Nachfolgend werden die Ergebnisse von Experiment 2b präsentiert, das mit der Intention durchgeführt wurde, den Grad an Sprechereinstellung bei Intensivierern mit unterschiedlich wahrnehmbarer Intensität zu ermitteln. {Dadurch} wird die Gültigkeit der Hypothese H6 überprüft, nach der von klasseninternen Unterschieden in Bezug auf den durch die Intensivierer ausgedrückten Intensitätsgrad auszugehen ist, wie von {\citet[87]{Biedermann.1969}} und \citet[408]{Breindl.2007} angeführt wird. Ferner besteht im Sinne von {\citet[133]{Gutzmann.2019a}} die Möglichkeit, dass die Expressivität der Ausdrücke auf deren Skalaritätsgrad einwirkt und letzterer deswegen als stärker wahrgenommen wird. Auch diese Annahme wird getestet.


Begonnen wird mit einer deskriptivstatistischen Analyse, auf die eine inferenzstatistische Analyse folgt. Im Zuge dessen werden die Ergebnisse der beiden {Experimente 2a und 2b} ebenso auf eine mögliche Korrelation zwischen den Variablen \textsc{Skalarität} und \textsc{Sprechereinstellung} hin untersucht. Die inferenzstatistische Analyse wurde auch in diesem Experiment mithilfe von R (Version 1.1.456, \citealt[]{R.2018}) und den Paketen \textit{lme4} ({\citealt[]{Bates.2015}}) und \textit{lmerTest} (\citealt[]{Kuznetsova.2017}) vorgenommen. Die Visualisierung der Daten erfolgte darüber hinaus mit \textit{ggplot2} (\citealt[]{Wickham.2016}).


\subsubsection{Vergleich der Stufen}\label{abs:Deskriptivstatistik02}


%Wie zuvor werden zunächst die Intensiviererklassen einander gegenübergestellt, woraufhin eine genaue Begutachtung der einzelnen Intensivierer folgt.



%\subsubsection{Vergleich der Intensiviererklassen}\label{Intensiv2} $~$ \\
\noindent In Tabelle \ref{SkaWe_E2} sind die gemittelten Skalenwerte und Standardabweichungen für beide Intensiviererklassen, basierend auf 434 Datenpunkten pro Intensivierer, aufgeführt. An dieser Stelle werden die je vier deskriptiven und expressiven Intensivierer zu einer Stufe zusammengefasst. Auch hier ist als Ergänzung der Mittelwert für beide Stufen zusammen angegeben. 


%
\begin{table}
%\setlength{\tabcolsep}{28.75pt}
\begin{tabularx}{.8\textwidth}{Xrrrr}
\lsptoprule
Intensiviererklasse & Skalenwert &  \textit{SD} & Mittelwert &  \textit{SD}\\
\midrule
deskriptiv          & 49,15      &  27,8        & 54,79      &  28,7\\
expressiv           & 60,43      & 28,5         & 54,79      &  28,7\\
\lspbottomrule
\end{tabularx}
\caption[Gemittelte Skalenwerte der Stufen bei Experiment 2b]{Gemittelte Skalenwerte und Standardabweichungen der Intensiviererklassen bei Experiment 2b -- Sprechereinstellung (in alphabetischer Reihenfolge).}
\label{SkaWe_E2}
\end{table}


\noindent Der Übersicht in Tabelle \ref{SkaWe_E2} ist zu entnehmen, dass Sätzen mit expressiven Intensivierern von den Versuchspersonen ein tendenziell höherer Skalenwert zugeschrieben wurde als Sätzen mit deskriptiven Intensivierern {($\Delta$ = ca. 11,30 Punkte)}. Die gemittelten Skalenwerte der Intensiviererklassen unterscheiden sich bei der in Experiment 2b abgefragten Sprechereinstellung also wesentlich stärker voneinander als bei der in Experiment 2a untersuchten Skalarität. Der Wert von letzterer belief sich nämlich bloß auf 4,50 Punkte. Dies legt nahe, dass die Untersuchungsmethode funktioniert hat und die Einteilung der Ausdrücke in deskriptiv und expressiv insgesamt {passend war}.


Den markanten Unterschied zwischen den beiden Intensiviererklassen zeigen auch die in {Abbildung \ref{Box2}} dargestellten Boxplots.


\begin{figure}
\includegraphics[width=1\textwidth]{Rplot_Box_E2}
\caption[Gemittelte Skalenwerte der Stufen bei Experiment 2b] {Gemittelte Skalenwerte der Intensiviererklassen bei {Experiment 2b --} Sprechereinstellung (x-Achse = Intensiviererklasse, y-Achse = Skalenwert).}
\label{Box2}
\end{figure}



\subsubsection{Regressionsanalyse zum Vergleich der Stufen}\label{LiGeMo02}
\largerpage

In diesem Abschnitt werden die Ergebnisse der Regressionsanalyse zur Untersuchung von Sprechereinstellung präsentiert. Wie das vorherige Modell wurde auch dieses durch Rückwärtsselektion und \textit{anova}-Modellvergleich in R berechnet. Wiederholt wurde zunächst ein volles Modell mit den Effekten aller unabhängigen Variablen und deren Zwei-Wege-Interaktionen berechnet, die sich möglicherweise auf die abhängige Variable ausgewirkt haben. Das volle Modell enthielt \textit{Fixed Effects} für den \textsc{Skalenwert} als metrische abhängige Variable sowie die kategorialen Prädiktoren \textsc{Intensiviererklasse} und \textsc{Position} sowie deren Zwei-Wege-Interaktionen. In {Experiment 2b} steht die Variable \textsc{Skalenwert} für den Grad an Sprechereinstellung, den die Versuchspersonen der jeweiligen Intensiviererklasse bzw. dem jeweiligen Intensivierer zugeschrieben haben. Um potenzielle Unterschiede zwischen den Intensiviererklassen zu finden, wurde die nominale Variable \textsc{Intensiviererklasse} dummy-gecoded. Dabei wurde `Deskriptiv' der Wert $0$ und `Expressiv' der Wert $1$ zugewiesen. Diese Ausprägungsformen wurden anschließend miteinander verglichen. Neben den \textit{Fixed Effects} wurden im vollen Modell versuchspersonenweise \textit{Random Intercepts} und \textit{Random Slopes} für die Variablen \textsc{Skalenwert}, \textsc{Intensiviererklasse} und \textsc{Position} miteinbezogen.\footnote{Da das volle Regressionsmodell unter Einbezug der Adjektive nicht konvergierte, konnten bei dieser Berechnung lediglich versuchspersonenweise \textit{Random Intercepts} und \textit{Random Slopes} miteingebunden werden.} Im Rahmen der Datenmodellierung wurden nach und nach all diejenigen Effekte getilgt, die den Modell-Fit nicht signifikant verbesserten und somit keinen Einfluss auf den Skalenwert hatten. Das volle und finale Modell ist im Anhang  \ref{Aja2} angegeben. Letzteres beinhaltet einen signifikanten Haupteffekt für den Prädiktor \textsc{Intensiviererklasse}. Die berechneten Koeffizienten der Regressionsanalyse befinden sich in {Tabelle \ref{LinReg2}}.

%
\begin{table}
%\setlength{\tabcolsep}{13pt}
%\begin{tabular}{p{0.425\textwidth} p{0.1\textwidth} p{0.1\textwidth} p{0.1\textwidth}p{0.1\textwidth} p{0.05\textwidth}}
\begin{tabularx}{.8\textwidth}{Xrrrrr}
\lsptoprule
Prädiktor & \textit{Est.} & \textit{SE} & \textit{$\chi$$^{2}$} & \textit{p} & \\
\midrule

(Intercept) & 49,02 & 2,68 & 39,31 & < 0,001 & $\star$$\star$$\star$\\
Intensiviererklasse & 11,02 & 2,70 & 14,51 & < 0,001 & $\star$$\star$$\star$\\
\lspbottomrule

\end{tabularx}
\caption[Finales Regressionsmodell von Experiment 2b (mit `Deskriptiv' als Referenzlevel)]{Finales Regressionsmodell von Experiment 2b -- Sprechereinstellung (mit `Deskriptiv' als Referenzlevel).}
\label{LinReg2}
\end{table}

\newpage
%
\noindent Den in Tabelle \ref{LinReg2} gelisteten Werten nach übt die unabhängige Variable {\textsc{Intensiviererklasse}} einen signifikanten Effekt auf den Skalenwert aus. Das bedeutet, dass Sätze mit expressiven Intensivierern als signifikant stärker einstellungsausdrückend wahrgenommen werden als Sätze, die deskriptive Intensivierer beinhalten {($\chi$$^{2}$\,(1) = 14,51, \textit{p} < 0,001)}, was wiederum mit der in Experiment 1 getesteten Hypothese H4 im Sinne von {\citet[315]{Lang.1983}} und \citet[57]{Ortner.2014} übereinstimmt (vgl. {Kapitel \ref{Experiment0}}).


Zusammengenommen erhärten die Ergebnisse der Datenauswertung von Experiment 2b den in der theoretischen Literatur angenommenen Unterschied zwischen der deskriptiven und expressiven Intensiviererklasse (vgl. \citealt[315]{Lang.1983}; \citealt[57]{Ortner.2014}), wurden Sätze mit expressiven Intensivierern von den Versuchspersonen doch mit einem signifikant höheren Grad an Sprechereinstellung assoziiert als Sätze mit deskriptiven Intensivierern.


\subsubsection[Einfluss von Sprechereinstellung auf Skalarität]{Beeinflusst die Sprechereinstellung die Skalarität?}


Im nächsten Schritt wird untersucht, ob zwischen der in Experiment 2a erfragten Skalarität und der in Experiment 2b fokussierten Sprechereinstellung eine statistisch signifikante Korrelation besteht. Vor diesem Hintergrund ist vorstellbar, dass ein höherer Grad an Sprechereinstellung bspw. im Falle einer positiven Korrelation gleichermaßen mit einem höheren Grad an Skalarität einhergeht. In dem Fall wären die Variablen voneinander abhängig.


Zur Quantifizierung des möglichen linearen Merkmalszusammenhangs wurde der Korrelationskoeffizient \textit{r} nach Pearson berechnet, dessen \textit{Estimate} wie folgt aussieht: |\textit{r}| = 0,79. Das Ergebnis des zugehörigen \textit{t}-Tests ist nicht signifikant ({\textit{t}\,(2) = 1,82}, \textit{p} = 0,21). Dies spricht gegen einen korrelativen Zusammenhang zwischen Skalarität und Sprechereinstellung -- die beiden Variablen sind in dieser Untersuchung also nicht voneinander abhängig. Dies lässt den Schluss zu, dass die Expressivität der Ausdrücke hier keinen Einfluss auf deren Skalaritätsgrad hatte. Aus welchen Gründen die Aussagekraft dieses Ergebnisses geschwächt ist, wird in Abschnitt \ref{abs:Diskussion2} {diskutiert}.


%Aus welchen Gründen dieses Ergebnis nur bedingt aussagekräftig ist, wird in Abschnitt \ref{abs:Diskussion2} {diskutiert}.


Wie schon bei dem in Kapitel \ref{Experiment0} beschriebenen Experiment 1 wurde auch bei dieser Untersuchung neben der Korrelationsanalyse ein lineares Regressionsmodell gerechnet. Dadurch kann getestet werden, ob der Grad an {Sprechereinstellung} im Experiment den Grad an Skalarität beeinflusste. Dabei ist denkbar, dass die Skalarität mit steigender Sprechereinstellung zunimmt. Es wurde also überprüft, ob die unabhängige Variable \textsc{Sprechereinstellung} einen Effekt auf die abhängige Variable \textsc{Skalarität} hat. Das lineare Modell beinhaltet \textit{Fixed Effects} für die \textsc{Skalarität} und die \textsc{Sprechereinstellung}. Die berechneten Koeffizienten sind in {Tabelle \ref{LM02}} aufgeführt.


\begin{table}
%\setlength{\tabcolsep}{9.4pt}
\begin{tabularx}{\textwidth}{Xrrrr}
\lsptoprule
Prädiktor & \textit{Est.} & \textit{SE} & \textit{F} & \textit{p}  \\
\midrule
(Intercept) & 9,05 & 35,67 & 0,06 & > 0,05  \\
Skalarität vs. Sprechereinstellung & 0,89 & \textcolor{white}{-}0,59 & 2,27 & = 0,27  \\
\lspbottomrule
\end{tabularx}
\caption[Lineares Regressionsmodell zum Einfluss von Sprechereinstellung auf Skalarität bei Experiment 2]{Lineares Regressionsmodell zum Einfluss von Sprechereinstellung auf Skalarität.}
\label{LM02}
\end{table}

\largerpage
\noindent Auch die Werte aus Tabelle \ref{LM02} sprechen gegen einen signifikanten Effekt der Sprechereinstellung auf die Skalarität ({$F$ = 2,27}, \textit{p} > 0,05). Davon ausgehend lässt sich schlussfolgern, dass der Grad an Sprechereinstellung den Grad an Skalarität im Experiment nicht beeinflusst hat. Dieser Befund liefert keine empirische Evidenz für die Annahme im Sinne von \citet[133]{Gutzmann.2019a}.



\begin{figure}[b]
\includegraphics[width=.8\textwidth]{{SkalExpr3}}
\caption [Lineare Regression zum Einfluss von Sprechereinstellung auf Skalarität bei Experiment 2]{Lineare Regression zum Einfluss von Sprechereinstellung auf Skalarität (x-Achse = Sprechereinstellung, y-Achse = Skalarität).}
\label{Korr02}
\end{figure}



Dies ist in Abbildung \ref{Korr02} als lineare Regression grafisch dargestellt. Wiederholt steht die Regressionslinie für die vorhergesagte, geschätzte Sprechereinstellung der in der Untersuchung berücksichtigten expressiven Intensivierer. 



%

%
\subsubsection{Vergleich der Intensivierer}\label{Intensiv4}


\noindent Als Nächstes werden die gemittelten Skalenwerte der einzelnen Intensivierer inspiziert. Durch dieses Vorgehen können die Ausdrücke im Hinblick auf den ihnen zugesprochenen Grad an Sprechereinstellung detaillierter miteinander verglichen werden. Für einen besseren Überblick sind diese Werte inklusive Standardabweichungen in Tabelle {\ref{Einzeln_2}} zusammengetragen. Die Anordnung der Ausdrücke je Klasse erfolgt nach ansteigendem Skalenwert.

%
\begin{table}
%\setlength{\tabcolsep}{27pt}
\begin{tabularx}{\textwidth}{XlYY}
\lsptoprule
Intensiviererklasse & Intensivierer  & Skalenwert &  \textit{SD}\\
\midrule
\rowcolor[gray]{.90}
deskriptiv & \textit{überaus}  & 46,99 & 27,4\\
           & \textit{äußerst}  & 47,40 & 25,7\\
 \rowcolor[gray]{.90}
           & \textit{ziemlich} & 49,06 & 28,2\\
           & \textit{sehr}     & 53,15 & 29,6\\
 \rowcolor[gray]{.90}
expressiv  & \textit{super}    & 57,15 & 28,2\\
           & \textit{voll}     & 57,23 & 27,8\\
 \rowcolor[gray]{.90}
           & \textit{mega}     & 63,67 & 28,0\\
           & \textit{sau}      & 63,69 & 29,4\\
\lspbottomrule
\end{tabularx}
\caption[Gemittelte Skalenwerte der einzelnen Intensivierer von Experiment 2b]{Gemittelte Skalenwerte und Standardabweichungen der einzelnen Intensivierer von Experiment 2b -- Sprechereinstellung (angeordnet nach ansteigendem Skalenwert).}
\label{Einzeln_2}
\end{table}

\noindent Tabelle \ref{Einzeln_2} zeigt ganz klar, dass die gemittelten Skalenwerte der deskriptiven Intensivierer ebenso detailliert betrachtet ausnahmslos unter den Werten der expressiven Intensivierer liegen: Während die deskriptiven Intensivierer auf der Skala im Bereich zwischen 45 und 55 Punkten angesiedelt sind, befinden sich die expressiven Intensivierer im Bereich zwischen 55 und 65 Punkten. Demzufolge wurden Sätze mit den aufgeführten expressiven Intensivierern von den Versuchspersonen den Erwartungen gemäß als stärker einstellungsausdrückend wahrgenommen als Sätze, die die gelisteten deskriptiven Intensivierer beinhalteten. Ebenfalls zu erkennen ist, dass der Unterschied zwischen den beiden Intensiviererklassen grundsätzlich größer ist als der, der hinsichtlich Skalarität zu beobachten war. So gibt es in {Experiment 2b} keinen deskriptiven Intensivierer, dessen Skalenwert an den der expressiven Intensivierer heranreicht. Auffällig ist außerdem, dass sich auch im Kontext der Sprechereinstellung deutlich ausgeprägte Abstufungen innerhalb der Intensiviererklassen herauskristallisieren, die in {Tabelle \ref{stark}} zu sehen sind.


\begin{table}
\small
%\setlength{\tabcolsep}{8.25pt}
%\begin{tabular}{p{0.15\textwidth} p{0.1\textwidth} p{0.15\textwidth} p{0.1\textwidth} p{0.15\textwidth}}
\begin{tabularx}{\textwidth}{X@{} lrr lrr }
\lsptoprule
 &  \multicolumn{6}{c}{\textbf{Intensiviererklasse}}  \\
\cmidrule(lr){2-7}\cmidrule(lr){5-7}
& \multicolumn{3}{c}{{deskriptiv}} & \multicolumn{3}{c}{{expressiv}}  \\
\cmidrule(lr){2-4}\cmidrule(lr){5-7}
{Stärkegrad}  & Intensivierer & Skalenwert& \textit{SD}  & Intensivierer & Skalenwert& \textit{SD}   \\
\midrule
schwach & \textit{überaus} & 46,99& 27,4 & \textit{super} &  57,15 & 28,2 \\
 & \textit{äußerst} & 47,40 & 25,7 & \textit{voll} & 57,23& 27,8   \\
\multicolumn{1}{l}{mittel} & \textit{ziemlich} & 49,06& 28,2 &  &  \\
\multicolumn{1}{l}{stark} & \textit{sehr} & 53,15& 29,6  & \textit{mega} & 63,67 & 28,0\\
\multicolumn{1}{l}{} & \textit{} & &   & \textit{sau} & 63,69& 29,4 \\
\lspbottomrule

\end{tabularx}
\caption[Klasseninterne Abstufungen bei Experiment 2b]{Klasseninterne Abstufungen bei Experiment 2b -- Sprechereinstellung (angeordnet nach ansteigendem Skalenwert).}
\label{stark}
\end{table}


Tabelle \ref{stark} legt innerhalb der deskriptiven Intensiviererklasse in Bezug auf Sprechereinstellung eine dreiteilige Abstufung nahe: {\textit{überaus}|\textit{äußerst} $\prec$ \textit{ziemlich} $\prec$ \textit{sehr}}. Demnach wurden die Ausdrücke \textit{überaus} und \textit{äußerst} von den Versuchspersonen vergleichbar bewertet ($\Delta$ = ca. 0,41 Punkte). Dagegen zeichnen sich innerhalb der expressiven Intensiviererklasse zwei Stufen mit je zwei Intensivierern ab: {\textit{super}|\textit{voll} $\prec$ \textit{mega}|\textit{sau}}. So ähneln sich einerseits die Skalenwerte von \textit{super} und \textit{voll} sehr stark {($\Delta$ = ca. 0,08 Punkte)}, genauso wie andererseits \textit{mega} und \textit{sau} auf der Skala beinahe gleichwertig beurteilt wurden {($\Delta$ = ca. 0,02 Punkte)}. Dieses Ergebnis ist in {Abbildung \ref{Bar2}} in Form eines Säulendiagramms visualisiert. Darüber hinaus ist im {Anhang \ref{Aja3}} in Tabelle \ref{Adj10} eine Übersicht der Ratings der Adjektive im Kontext der Sprechereinstellung zu finden. Darin ist zu sehen, dass deren Mittelwerte in einem ähnlichen Skalenbereich liegen: So befindet sich das am niedrigsten eingeschätzte Adjektiv \textit{breit} auf der Skala im Schnitt bei 49,94 Punkten, wohingegen \textit{heiß} mit durchschnittlich 58,52 Punkten am höchsten verortet wurde. Erneut deutet die nur geringe Differenz von {ca. 8,58 Punkte} darauf hin, dass das Skalenrating als Untersuchungsmethode funktioniert hat. Wäre die Differenz zwischen den Adjektiven an dieser Stelle stärker ausgeprägt, so ließe sich mutmaßen, dass der festgestellte Unterschied zwischen den beiden Intensiviererklassen durch die verwendeten Adjektive verursacht sein könnte. Dies ist augenscheinlich nicht der Fall.



%
\begin{figure}
\includegraphics[width=1\textwidth]{Rplot_Bar_E2}
\caption[Säulendiagramm zu den einzelnen Intensivierern von Experiment 2b] {Säulendiagramm zu den einzelnen deskriptiven und expressiven Intensivierern von Experiment 2b -- Sprechereinstellung (x-Achse = Intensivierer, y-Achse = Skalenwert).}
\label{Bar2}
\end{figure}






\subsubsection{Paarweise Vergleiche der Intensivierer}\label{Paar2}


Wie weiter oben beschrieben, basierten die Paarvergleiche auf der Intention, die Unterschiede zwischen den einzelnen Ausdrücken auf Signifikanz hin zu überprüfen. Dazu wurden die einzelnen deskriptiven und expressiven Intensivierer analog zum Vorgehen bei Experiment 2a unter Zuhilfenahme des statistischen Modellierungsverfahrens \textit{t}-Test gepaart miteinander verglichen. In der daraus resultierenden Matrix bildet das überschneidende Element die Testprüfgröße \textit{t} sowie den jeweiligen Signifikanzwert \textit{p} ab. Auch hier wurde das Signifikanzniveau der $\alpha$-Adjustierung nach Bonferroni unterzogen, um die Alphafehler-Kumulierung zu eliminieren. Im Zuge dessen wurde der {\textit{p}-Wert} jedes kontrastierten Intensiviererpaars mit der Gesamtzahl der Vergleiche multipliziert. Die Gesamtzahl beträgt wie im vorherigen {Experiment 2a} 28, wurden doch alle Intensivierer paarweise miteinander gekreuzt: \textit{überaus} vs. \textit{äußerst}, \textit{überaus} vs. \textit{ziemlich}, \textit{überaus} vs. \textit{sehr}, \textit{überaus} vs. \textit{super}, \textit{überaus} vs. \textit{voll}, \textit{überaus} vs. \textit{mega}, \textit{überaus} vs. \textit{sau}, \textit{äußerst} vs. \textit{ziemlich}, \textit{äußerst} vs. \textit{sehr} usw. Durch die Multiplikation wird der jeweilige \textit{p}-Wert größer, sodass er entweder weniger oder auch überhaupt nicht mehr signifikant ist.


Die aus den Berechnungen resultierende Matrix ist in Tabelle \ref{Sign2_neu} dargestellt. Dabei werden zunächst die Vergleiche der deskriptiven Intensivierer (in ansteigendem Skalenwert) vorgenommen, an die sich die Vergleiche der expressiven Intensivierer (ebenfalls in ansteigendem Skalenwert sortiert) anschließen. Erneut sei klargestellt, dass bei den aufgelisteten \textit{t}-Werten aus Gründen der Darstellbarkeit von einer ausführlichen Dokumentation mit eingeklammerter Angabe der Anzahl der Freiheitsgrade, die auf der Größe der Stichprobe beruht, abgesehen wird {(\textit{t}\,(\textit{\textit{df}}) = \textelp{})}. Der Freiheitsgrad entspricht bei diesen Paarvergleichen dem Wert 3389. Insofern bedeutet die in der Tabelle jeweils aufgeführte verkürzte Form (z.\,B. {\textit{t} = -0,28}, {\textit{p} > 0,05}) streng genommen Folgendes: \textit{t}\,(3389) = \textelp{}, veranschaulicht am eingeklammerten Beispiel also {\textit{t}\,(3389) = -0,28}, {\textit{p} > 0,05}. Die Formel der paarweisen Vergleiche ist außerdem im Anhang \ref{Aja2} angegeben.




%\cellcolor[gray]{0.90}
%\cellcolor[gray]{0.75}
%\cellcolor[gray]{0.59}
%\hfil \ding{55}
%$\star$$\star$$\star$


%
\begin{sidewaystable}
\small
%\setlength{\tabcolsep}{8pt}

\begin{tabularx}{\textwidth}{l@{~~}QQQQ@{}Q@{}QQ@{}Q@{}}
\lsptoprule
Intensivierer & \textit{überaus} & \textit{äußerst} & \textit{ziemlich} & \textit{sehr} & \textit{super} & \textit{voll} & \textit{mega} & \textit{sau} \\
\midrule

\multirow{3}{*}{\textit{überaus}} &  &   \textit{t} = 0,28, & \textit{t} = 1,29, &  \cellcolor[gray]{0.59} \textit{t} = 3,92, & \cellcolor[gray]{0.59} \textit{t} = 6,47, &  \cellcolor[gray]{0.59} \textit{t} = 6,51, & \cellcolor[gray]{0.59}  \textit{t} = 10,62, &  \cellcolor[gray]{0.59} \textit{t} = 10,64, \\ 
& \hfil \ding{55} &   \textit{p} > 0,05   & \textit{p} > 0,05  & \cellcolor[gray]{0.59}  \textit{p} < 0,001 &   \cellcolor[gray]{0.59}  \textit{p} < 0,001  &   \cellcolor[gray]{0.59}  \textit{p} < 0,001   &   \cellcolor[gray]{0.59}  \textit{p} < 0,001    &  \cellcolor[gray]{0.59}  \textit{p} < 0,001    \\ 
 &  &  &  & \cellcolor[gray]{0.59} $\star$$\star$$\star$ & \cellcolor[gray]{0.59} $\star$$\star$$\star$ & \cellcolor[gray]{0.59} $\star$$\star$$\star$ & \cellcolor[gray]{0.59} $\star$$\star$$\star$ &  \cellcolor[gray]{0.59} $\star$$\star$$\star$ \\ 
 
\multirow{3}{*}{\textit{äußerst}} & \textit{t} = -0,28, &   & \textit{t} = 1,01, & \cellcolor[gray]{0.75} \textit{t} = 3,63, & \cellcolor[gray]{0.59} \textit{t} = 6,18, & \cellcolor[gray]{0.59}  \textit{t} = 6,22, &  \cellcolor[gray]{0.59} \textit{t} = 10,32, &  \cellcolor[gray]{0.59}  \textit{t} = 10,34, \\ 
&  \textit{p} > 0,05 &   \hfil \ding{55}   & \textit{p} > 0,05  & \cellcolor[gray]{0.75} \textit{p} < 0,01 &   \cellcolor[gray]{0.59}  \textit{p} < 0,001  &   \cellcolor[gray]{0.59}  \textit{p} < 0,001  &  \cellcolor[gray]{0.59}  \textit{p} < 0,001    &  \cellcolor[gray]{0.59}  \textit{p} < 0,001  \\ 
 &  &  &  & \cellcolor[gray]{0.75} $\star$$\star$ &  \cellcolor[gray]{0.59} $\star$$\star$$\star$ & \cellcolor[gray]{0.59} $\star$$\star$$\star$ & \cellcolor[gray]{0.59} $\star$$\star$$\star$ &  \cellcolor[gray]{0.59} $\star$$\star$$\star$ \\  
 
\multirow{3}{*}{\textit{ziemlich}} & \textit{t} = -1,29, &   \textit{t} = -1,01, &  &  \textit{t} = 2,63, & \cellcolor[gray]{0.59} \textit{t} = 5,18, & \cellcolor[gray]{0.59}  \textit{t} = 5,22, & \cellcolor[gray]{0.59} \textit{t} = 9,33, &  \cellcolor[gray]{0.59} \textit{t} = 9,35, \\ 
& \textit{p} > 0,05  &   \textit{p} > 0,05   & \hfil \ding{55}  & \textit{p} > 0,05 &  \cellcolor[gray]{0.59}  \textit{p} < 0,001  &   \cellcolor[gray]{0.59}  \textit{p} < 0,001  &   \cellcolor[gray]{0.59}  \textit{p} < 0,001  &  \cellcolor[gray]{0.59}  \textit{p} < 0,001   \\ 
 &  &  &  &  & \cellcolor[gray]{0.59} $\star$$\star$$\star$  & \cellcolor[gray]{0.59} $\star$$\star$$\star$ & \cellcolor[gray]{0.59} $\star$$\star$$\star$ &\cellcolor[gray]{0.59} $\star$$\star$$\star$   \\ 
 
\multirow{3}{*}{\textit{sehr}} & \cellcolor[gray]{0.59}  \textit{t} = -3,92, &  \cellcolor[gray]{0.75} \textit{t} = -3,63, & \textit{t} = -2,63, &   & \textit{t} = 2,56, &   \textit{t} = 2,59, & \cellcolor[gray]{0.59} \textit{t} = 6,70, &  \cellcolor[gray]{0.59} \textit{t} = 6,73, \\ 
&  \cellcolor[gray]{0.59}  \textit{p} < 0,001  &   \cellcolor[gray]{0.75} \textit{p} < 0,01    & \textit{p} > 0,05  & \hfil \ding{55}  &   \textit{p} > 0,05 &   \textit{p} > 0,05  &   \cellcolor[gray]{0.59}  \textit{p} < 0,001   &    \cellcolor[gray]{0.59}  \textit{p} < 0,001 \\ 
 & \cellcolor[gray]{0.59} $\star$$\star$$\star$ & \cellcolor[gray]{0.75} $\star$$\star$ &  &  &  &  & \cellcolor[gray]{0.59} $\star$$\star$$\star$ & \cellcolor[gray]{0.59} $\star$$\star$$\star$  \\ 
 
\multirow{3}{*}{\textit{super}} & \cellcolor[gray]{0.59} \textit{t} = -6,47, &  \cellcolor[gray]{0.59} \textit{t} = -6,18, &   \cellcolor[gray]{0.59}  \textit{t} = -5,18, & \textit{t} = -2,56, &  &   \textit{t} = 0,04, &   \cellcolor[gray]{0.59} \textit{t} = 4,14, &  \cellcolor[gray]{0.59}  \textit{t} = 4,17, \\ 
&  \cellcolor[gray]{0.59}  \textit{p} < 0,001  &    \cellcolor[gray]{0.59}  \textit{p} < 0,001    &   \cellcolor[gray]{0.59}  \textit{p} < 0,001   & \textit{p} > 0,05  &   \hfil \ding{55} &   \textit{p} > 0,05  &  \cellcolor[gray]{0.59}  \textit{p} < 0,001  &   \cellcolor[gray]{0.59}  \textit{p} < 0,001  \\ 
 & \cellcolor[gray]{0.59} $\star$$\star$$\star$ & \cellcolor[gray]{0.59} $\star$$\star$$\star$ & \cellcolor[gray]{0.59} $\star$$\star$$\star$ &  &  &  & \cellcolor[gray]{0.59} $\star$$\star$$\star$ &  \cellcolor[gray]{0.59} $\star$$\star$$\star$ \\  
 
\multirow{3}{*}{\textit{voll}} & \cellcolor[gray]{0.59} \textit{t} = -6,51, &  \cellcolor[gray]{0.59}  \textit{t} = -6,22, & \cellcolor[gray]{0.59} \textit{t} = -5,22, &  \textit{t} = -2,59, & \textit{t} = -0,04, &  &  \cellcolor[gray]{0.59} \textit{t} = 4,11, &  \cellcolor[gray]{0.59} \textit{t} = 4,14, \\ 
&  \cellcolor[gray]{0.59}  \textit{p} < 0,001   &  \cellcolor[gray]{0.59}  \textit{p} < 0,001    &  \cellcolor[gray]{0.59}  \textit{p} < 0,001  & \textit{p} > 0,05 &   \textit{p} > 0,05   &  \hfil \ding{55}   & \cellcolor[gray]{0.59}  \textit{p} < 0,001     &    \cellcolor[gray]{0.59}  \textit{p} < 0,001   \\ 
 & \cellcolor[gray]{0.59} $\star$$\star$$\star$ & \cellcolor[gray]{0.59} $\star$$\star$$\star$ & \cellcolor[gray]{0.59} $\star$$\star$$\star$ &  &  &  & \cellcolor[gray]{0.59} $\star$$\star$$\star$ &  \cellcolor[gray]{0.59} $\star$$\star$$\star$ \\ 
  
\multirow{3}{*}{\textit{mega}} & \cellcolor[gray]{0.59} \textit{t} = -10,62, &  \cellcolor[gray]{0.59} \textit{t} = -10,32, & \cellcolor[gray]{0.59} \textit{t} = -9,33, &  \cellcolor[gray]{0.59} \textit{t} = -6,70, & \cellcolor[gray]{0.59} \textit{t} = -4,14, &  \cellcolor[gray]{0.59} \textit{t} = -4,11, &   &   \textit{t} = 0,03, \\ 
&  \cellcolor[gray]{0.59}  \textit{p} < 0,001   &   \cellcolor[gray]{0.59}  \textit{p} < 0,001   &  \cellcolor[gray]{0.59}  \textit{p} < 0,001  & \cellcolor[gray]{0.59}  \textit{p} < 0,001    &  \cellcolor[gray]{0.59}  \textit{p} < 0,001 &   \cellcolor[gray]{0.59}  \textit{p} < 0,001  &   \hfil \ding{55}   &   \textit{p} > 0,05   \\ 
 &  \cellcolor[gray]{0.59} $\star$$\star$$\star$ & \cellcolor[gray]{0.59} $\star$$\star$$\star$ & \cellcolor[gray]{0.59} $\star$$\star$$\star$ & \cellcolor[gray]{0.59} $\star$$\star$$\star$ & \cellcolor[gray]{0.59} $\star$$\star$$\star$ & \cellcolor[gray]{0.59} $\star$$\star$$\star$ &  &   \\ 
 
\multirow{3}{*}{\textit{sau}} & \cellcolor[gray]{0.59} \textit{t} = -10,64, &  \cellcolor[gray]{0.59} \textit{t} = -10,34, & \cellcolor[gray]{0.59} \textit{t} = -9,35, & \cellcolor[gray]{0.59} \textit{t} = -6,73, & \cellcolor[gray]{0.59} \textit{t} = -4,17, &  \cellcolor[gray]{0.59} \textit{t} = -4,14, &  \textit{t} = -0,03, &    \\ 
&  \cellcolor[gray]{0.59}  \textit{p} < 0,001 &   \cellcolor[gray]{0.59}  \textit{p} < 0,001   & \cellcolor[gray]{0.59}  \textit{p} < 0,001   & \cellcolor[gray]{0.59}  \textit{p} < 0,001   & \cellcolor[gray]{0.59}  \textit{p} < 0,001   &    \cellcolor[gray]{0.59}  \textit{p} < 0,001   &  \textit{p} > 0,05   &   \hfil \ding{55}  \\ 
 & \cellcolor[gray]{0.59} $\star$$\star$$\star$ & \cellcolor[gray]{0.59} $\star$$\star$$\star$ & \cellcolor[gray]{0.59} $\star$$\star$$\star$ & \cellcolor[gray]{0.59} $\star$$\star$$\star$ &  \cellcolor[gray]{0.59} $\star$$\star$$\star$  & \cellcolor[gray]{0.59} $\star$$\star$$\star$  &  &   \\
\lspbottomrule

\end{tabularx}
\caption[Ergebnisse der Paarvergleiche der einzelnen Intensivierer von Experiment 2b]{Ergebnisse der Paarvergleiche der einzelnen Intensivierer bei Experiment 2b -- Sprechereinstellung (nach Bonferroni-Korrektur; angeordnet nach ansteigendem Skalenwert).}
\label{Sign2_neu}
\end{sidewaystable}



Die in Tabelle \ref{Sign2_neu} abgebildeten signifikanten Ergebnisse der Paarvergleiche werden nachfolgend knapp umrissen.\footnote{Auch hier wird der signifikante \textit{p}-Wert lediglich bei der ersten Erwähnung eines kontrastierten Intensiviererpaars aufgeführt (z.\,B. \textit{überaus} vs. \textit{sehr}), während er bei der Zweitnennung außen vor gelassen wird (\textit{sehr} vs. \textit{überaus}).} Begonnen wird mit den Vergleichen des am niedrigsten positionierten deskriptiven Intensivierers \textit{überaus}. Hier sind signifikante Unterschiede bei den Kontrasten mit \textit{sehr} ({\textit{t}\,(3389) = 3,92}, \textit{p} < 0,001), \textit{super} (\textit{t}\,(3389) = 6,47, \textit{p} < 0,001), \textit{voll} (\textit{t}\,(3389) = 6,51, {\textit{p} < 0,001}), \textit{mega} (\textit{t}\,(3389) = 10,62, \textit{p} < 0,001) und  \textit{sau} ({\textit{t}\,(3389) = 10,64}, {\textit{p} < 0,001)} auszumachen. Der nächste deskriptive Intensivierer ist \textit{äußerst}. Dieser Ausdruck weist gleichfalls signifikante Unterschiede bei den Vergleichen mit \textit{sehr} (\textit{t}\,(3389) = 3,63, {\textit{p} < 0,01)}, \textit{super} (\textit{t}\,(3389) = 6,18, \textit{p} < 0,001), \textit{voll} ({\textit{t}\,(3389) = 6,22}, {\textit{p} < 0,001}), \textit{mega} {(\textit{t}\,(3389) = 10,32}, \textit{p} < 0,001) und \textit{sau} ({\textit{t}\,(3389) = 10,34}, {\textit{p} < 0,001)} auf. Bei \textit{ziemlich} sind signifikante Unterschiede bei den Kontrasten mit \textit{super} {(\textit{t}\,(3389) = 5,18,} \textit{p} < 0,001), \textit{voll} (\textit{t}\,(3389) = 5,22, {\textit{p} < 0,001}), \textit{mega} (\textit{t}\,(3389) = 9,33, {\textit{p} < 0,001}) und \textit{sau} {(\textit{t}\,(3389) = 9,35,} {\textit{p} < 0,001}) zu beobachten. Als Nächstes werden die Ergebnisse des am höchsten eingestuften deskriptiven Intensivierers \textit{sehr} begutachtet, die bei \textit{überaus}, \textit{äußerst}, \textit{mega} (\textit{t}\,(3389) = 6,70, {\textit{p} < 0,001)} und \textit{sau} ({\textit{t}\,(3389) = 6,73}, {\textit{p} < 0,001}) signifikant sind. Den Resultaten der paarweisen Vergleiche nach zu urteilen, unterscheiden sich also drei deskriptive Intensivierer gänzlich von den untersuchten expressiven Intensivierern, was im Anschluss zu diskutieren ist: \textit{überaus}, \textit{äußerst} {und \textit{ziemlich}}.


Bei den expressiven Intensivierern wird mit den Ergebnissen der Kontraste von \textit{super} begonnen, die bei \textit{überaus}, \textit{äußerst}, \textit{ziemlich}, \textit{mega} ({\textit{t}\,(3389) = 4,14}, {\textit{p} < 0,001})  und \textit{sau} ({\textit{t}\,(3389) = 4,17}, {\textit{p} < 0,001}) signifikant sind. Gleichsam zeigt auch \textit{voll} signifikante Unterschiede bei den Kontrasten mit \textit{überaus}, \textit{äußerst}, \textit{ziemlich}, \textit{mega} ({\textit{t}\,(3389) = 4,11}, {\textit{p} < 0,001}) und \textit{sau} ({\textit{t}\,(3389) = 4,14}, {\textit{p} < 0,001}) auf. Der nächste expressive Intensivierer in der Reihe ist \textit{mega}. Wie bereits beschrieben, sind bei diesem Ausdruck signifikante Unterschiede bei den Vergleichen mit \textit{überaus}, \textit{äußerst}, \textit{ziemlich}, \textit{sehr}, \textit{super} und \textit{voll} zu erkennen, genauso wie sich der höchstplatzierte expressive Intensivierer \textit{sau} signifikant von den soeben genannten Intensivierern unterscheidet.


Von den Ergebnissen der Paarvergleiche ausgehend kann geschlussfolgert werden, dass sich zwei Intensivierer klar im Hinblick auf den ihnen zugeschriebenen Grad an Sprechereinstellung von den restlichen Intensivierern unterscheiden, nämlich \textit{mega} und \textit{sau}. Die Kontraste dieser Ausdrücke sind in sechs von sieben Fällen signifikant. Die inter- und intraklassischen Unterschiede der vorgenommenen Vergleiche werden in Abschnitt \ref{abs:Diskussion2} diskutiert.








%
\subsection{Diskussion}\label{abs:Diskussion2}
Nachfolgend werden die Ergebnisse von Experiment 2b besprochen und mit Blick auf die oben beschriebene Hypothese  H6 bewertet. Nach dieser ist von graduellen Abstufungen innerhalb der Intensiviererklassen auszugehen, die auf die Semantik der Ausdrücke zurückzuführen sind (vgl. {\citealt[87]{Biedermann.1969}}; \citealt[408]{Breindl.2007}). Im Sinne von \citet[133]{Gutzmann.2019a} besteht außerdem die Möglichkeit, dass die Expressivität der Intensivierer auf deren Skalaritätsgrad einwirkt, letzterer infolge der Expressivität als stärker wahrgenommen wird und die beiden Variablen \textsc{Skalarität} und \textsc{Sprechereinstellung} miteinander korrelieren. Da es sinnvoll ist, {Experiment 2a}, das zur Erhebung des Skalaritätsgrades durchgeführt wurde, und das soeben geschilderte Experiment 2b in ihrer Gesamtheit zu betrachten, erfolgt die finale Ergebnisinterpretation erst in Abschnitt \ref{abs:Gesamtdiskussion}.


In Experiment 2b ging es für die Versuchspersonen darum, mithilfe eines Schiebereglers auf einer Skala von `gar nicht bewertend' bis `maximal bewertend' den im Satz enthaltenen Grad an Sprechereinstellung einzuschätzen, bspw. \textit{Wie bewertend ist die Aussage \glqq{}Der Berg ist äußerst hoch\grqq{}?}


Bei der Begutachtung der Intensiviererklassen stellte sich heraus, dass Sätze mit expressiven Intensivierern von den Versuchspersonen mit einem signifikant höheren Grad an Sprechereinstellung assoziiert wurden als Sätze mit deskriptiven Intensivierern (vgl. Tabelle \ref{LinReg2}). Wie schon in dem in Kapitel \ref{Experiment0} besprochenen Experiment 1 erhärtet dieser Befund das Zutreffen der Annahme eines wahrnehmbaren Unterschiedes in der Expressivität von deskriptiven und expressiven Intensivierern (vgl. die Hypothese H4 nach {\citealt[315]{Lang.1983}} und \citealt[57]{Ortner.2014}). Den Ergebnissen der beiden Untersuchungen nach zu urteilen, wurde der Grad an Sprechereinstellung bei expressiven Intensivierern als signifikant stärker empfunden als bei deskriptiven Intensivierern, was als weitere Evidenz für die Gültigkeit der {Hypothese H4} herangezogen werden kann. Dennoch ergibt sich die Frage, wieso die untersuchten deskriptiven Intensivierer trotz fehlender Expressivität nicht mit einem noch niedrigeren Grad an Sprechereinstellung verbunden wurden. So ist gemäß der Definition von \citet[133]{Gutzmann.2019a} anzunehmen, dass deskriptive Intensivierer -- anders als expressive Intensivierer -- überhaupt keine Sprechereinstellung ausdrücken und daher am unteren Ende einer kontinuierlichen Skala verortet werden. Eine solch niedrige Skaleneinschätzung zeigte sich in {Experiment 2b} jedoch nicht. Infolgedessen sind die Ergebnisse insofern bemerkenswert, als sie zu bedenken geben, dass deskriptive Intensivierer durchaus als einstellungsausdrückend wahrgenommen werden, wenn auch nicht dergestalt, dass sie an die expressiven Intensivierer heranreichen.


Auch für diese Untersuchung waren die in der theoretischen Literatur angenommenen intraklassischen Abstufungen zentral, weswegen die Intensiviererklassen im nächsten Analyseschritt aufgetrennt und ihre Mitglieder im Detail begutachtet wurden. Erneut waren klasseninterne Abstufungen auszumachen: \textit{überaus}|\textit{äußerst} $\prec$ \textit{ziemlich} $\prec$ \textit{sehr} sowie \textit{super}|\textit{voll} $\prec$ \textit{mega}|\textit{sau}. Somit gibt es Ausdrücke, die sich nach dem Empfinden der Versuchspersonen in Hinsicht auf Sprechereinstellung gleichen: \textit{überaus} und \textit{äußerst} aufseiten der deskriptiven Intensivierer und \textit{super} und \textit{voll} sowie \textit{mega} und \textit{sau} aufseiten der expressiven Intensivierer. Das bedeutet, dass die Ergebnisse von Experiment 2b im Wesentlichen ebenso für die zu untersuchende {Hypothese H6} sprechen. Trotz der Abstufungen unterscheidet sich die beobachtete Rangfolge nahezu gänzlich von der, die sich im Kontext der in {Experiment 2a} erfragten Skalarität zeigte. Wurde der deskriptive Intensivierer \textit{ziemlich} in {Experiment 2a} mit dem niedrigsten Skalaritätsgrad assoziiert, wurde der mit ihm einhergehende Grad an Sprechereinstellung in Experiment 2b hingegen umso höher beurteilt. Auch wenn der Ausdruck nicht übermäßig skalierend zu sein scheint, wird er dennoch als vergleichsweise stark bewertend wahrgenommen. Gegenteiliges trifft auf den deskriptiven Intensivierer \textit{überaus} zu, dem im Kontext der Skalarität von den Versuchspersonen der höchste Grad zugeschrieben wurde. Die Ergebnisse von {Experiment 2b} weisen jedoch darauf hin, dass der Ausdruck bloß als schwach einstellungsausdrückend empfunden wird. Die Position des deskriptiven Intensivierers \textit{äußerst} ist gleich geblieben, wurde der Ausdruck doch sowohl bezüglich Skalarität als auch Sprechereinstellung an zweiter Stelle verortet. Während das deskriptive \textit{sehr} bereits mit einem hohen Skalaritätsgrad in Verbindung gebracht wurde, wurde der durch ihn ausgedrückte Grad an Sprechereinstellung noch höher eingeschätzt, was dazu führt, dass der Ausdruck diesbezüglich die höchste Skalenposition innerhalb der deskriptiven Intensiviererklasse einnimmt. Daneben hat sich die Rangfolge innerhalb der expressiven Klasse ebenfalls geändert, wenn auch in erster Linie in dem Sinne, dass die Ausdrücke den Platz mit ihrem Nachbarn getauscht haben. Die soeben dargelegten Unterschiede, die in Bezug auf die Einschätzung von Skalarität und Sprechereinstellung auszumachen sind, legen in Summe nahe, dass die Untersuchungsmethode funktioniert hat und die Versuchspersonen die beiden Merkmale unterschiedlich wahrgenommen und {evaluiert haben}.


Aus der Intention heraus die festgestellten Unterschiede auf Signifikanz hin zu überprüfen, wurden die Intensivierer daraufhin gepaart miteinander verglichen (vgl. Tabelle \ref{Sign2_neu}). In diesem Kontext zeigte sich nicht nur eine signifikante Zweiteilung der Ausdrücke nach deskriptiv und expressiv, sondern auch, dass sich fast alle der ähnlich eingestuften Ausdrücke angesichts der Signifikanzen parallel verhalten. Dabei unterscheiden sich vor allem zwei Intensivierer klar im Hinblick auf den ihnen zugesprochenen Grad an Sprechereinstellung, nämlich \textit{mega} und \textit{sau}. Deren Vergleiche sind in sechs von sieben Fällen signifikant: bei allen deskriptiven Intensivierern, d.\,h. bei \textit{überaus}, \textit{äußerst}, \textit{ziemlich} und \textit{sehr}, sowie den expressiven Intensivierern \textit{super} und \textit{voll}. Lediglich die Gegenüberstellung der beiden Ausdrücke ist nicht signifikant, was durch die ähnliche Skaleneinschätzung bedingt ist. Wie sind die signifikanten Unterschiede zu erklären? Der Intensivierer \textit{sau} wurde bereits im Kontext der Skalarität als vergleichsweise hoch empfunden und so wurde er ebenso bezüglich Sprechereinstellung auf der Skala an höchster Position verortet. Bei diesem Ausdruck ist, wie oben beschrieben, denkbar, dass dieses Ergebnis der transparenten Semantik und der damit verbundenen derben Konnotation geschuldet ist, die möglicherweise zur hohen Gradeinschätzung beigetragen haben. Dagegen scheint bei \textit{mega} vor allem die Aktualität des Ausdrucks auf die Wahrnehmung einzuwirken, in deren Folge er gleichermaßen als übermäßig skalar und stark expressiv empfunden wird. Dies wirft die Frage auf, wie die Expressivität von \textit{mega} in ein paar Jahren eingeschätzt wird, sollte der Ausdruck bis dahin inflationär im Gebrauch sein. Ähnliches ist bei den Kontrasten der expressiven Intensivierer \textit{voll} und \textit{super} zu beobachten, deren Vergleiche aufgrund ähnlicher Skalenwerte ebenfalls zu gleichen signifikanten Ergebnissen führten. Diese treten in fünf von sieben Fällen auf: bei den deskriptiven Intensivierern \textit{überaus}, \textit{äußerst} und \textit{ziemlich} sowie den am höchsten eingeschätzten expressiven Intensivierern \textit{mega} und \textit{sau}. Dies deutet darauf hin, dass \textit{super} und \textit{voll} zwar als expressiver wahrgenommen werden als die deskriptiven Intensivierer, der durch sie ausgedrückte Grad an Sprechereinstellung jedoch nicht an \textit{mega} und \textit{sau} heranreicht. Auch innerhalb der deskriptiven Intensiviererklasse zeichnen sich infolge ähnlicher Skaleneinschätzung parallele Signifikanzen ab, sind die Ergebnisse der Paarvergleiche von \textit{äußerst} und \textit{überaus} doch in fünf von sieben Fällen signifikant: bei dem in Bezug auf Sprechereinstellung am höchsten eingeschätzten deskriptiven Intensivierer \textit{sehr} sowie den expressiven Intensivierern \textit{super}, \textit{voll}, \textit{mega} und \textit{sau}. Somit unterscheiden sich die niedrig bewerteten deskriptiven Intensivierer \textit{äußerst} und \textit{überaus} nicht nur signifikant von allen expressiven Intensivierern, sondern auch von dem am höchsten positionierten deskriptiven Intensivierer \textit{sehr}. Weil der Grad an Sprechereinstellung bei \textit{sehr} von den Versuchspersonen im Mittelfeld der Skala angesiedelt wurde, sind seine Vergleiche in vier Fällen signifikant: bei \textit{äußerst} und \textit{überaus} sowie den expressiven Intensivierern \textit{mega} und \textit{sau}. Der im Kontext der Sprechereinstellung vergleichsweise hoch eingeschätzte deskriptive Intensivierer \textit{ziemlich} unterscheidet sich hingegen signifikant von allen expressiven Intensivierern, d.\,h. von \textit{super}, \textit{voll}, \textit{mega} und \textit{sau}.


Auf Grundlage der Ergebnisse von Experiment 2b lässt sich schlussfolgern, dass die zu überprüfende Hypothese H6 im Kern korrekt ist. Denn auch in dieser Untersuchung treten in beiden Intensiviererklassen graduelle Abstufungen bezüglich des mit den Ausdrücken assoziierten Grades an Sprechereinstellung auf: schwach $\prec$ mittel $\prec$ stark. 
%In diesem Sinne können die Intensiviererklassen auf zwei verschiedenen Skalen nebeneinander geordnet werden. Wie bereits bei Skalarität wird dies in Tabelle \ref{stark} verbildlicht. Auch hier erfolgt die Anordnung der Intensiviererklassen in ansteigendem Skalenwert.



%%
%\begin{table}[h]
%
%%\setlength{\tabcolsep}{8.25pt}
%%\begin{tabular}{p{0.15\textwidth} p{0.1\textwidth} p{0.15\textwidth} p{0.1\textwidth} p{0.15\textwidth}}
%\begin{tabularx}{.8\textwidth}{Xlclc}
%\lsptoprule
% & & \multicolumn{2}{c}{\textbf{Intensiviererklasse}}  \\ [0.5em]
%\multirow{2}{*}{Stärkegrad} & \multicolumn{2}{c}{\textsc{deskriptiv}} & \multicolumn{2}{c}{\textsc{expressiv}}  \\ 
% & Intensivierer & Skalenwert & Intensivierer & Skalenwert  \\ 
%  &  & (\textit{SD})  &  & (\textit{SD})  \\
%\midrule 
%\\
%schwach & \textit{überaus} & 46,99 & \textit{super} & 57,15  \\
% & & (27.4) & & (28.2)  \\
% & \textit{äußerst} & 47,40 & \textit{voll} & 57,23 \\
%& & (25.7) &  & (27.8)\\
%\multicolumn{1}{l}{mittel} & \textit{ziemlich} & 49,06 & \textit{mega} & 63,67\\
%  & & (28.2) & & (28.0)\\
%stark & \textit{sehr} & 53,15 & \textit{sau} & 63,69  \\
% & & (29.6) & & (29.4)  \\
%\end{tabularx}
%\caption[Klasseninterne graduelle Abstufungen bei Experiment 2b]{Klasseninterne graduelle Abstufungen bei Experiment 2b -- Sprechereinstellung (in ansteigendem Skalenwert).}
%\label{stark}
%\end{table}









%
%\subsubsection{Paarweise Vergleiche der Intensivierer}\label{abs:Intensiviererpaarvergl} $~$ \\
%Die aus den Ergebnissen der Paarvergleiche resultierende Matrix deutete auf den ersten Blick auf eine signifikante Zweiteilung der Intensivierer nach deskriptiv und expressiv hin (vgl. Tabelle \ref{Sign2}). So zeigten nahezu alle Intensivierer signifikante Unterschiede bei den Vergleichen mit der gegenteiligen Intensiviererklasse, d.\,h. deskriptiv vs. expressiv und expressiv vs. deskriptiv. Die einzige Ausnahme stellte der deskriptive Intensivierer \textit{sehr} dar: Der gemittelte Skalenwert dieses Ausdrucks lag mit 53,15 Punkten nahe beim am niedrigsten positionierten expressiven Intensivierer \textit{super}, der auf der Skala im Schnitt bei 57,15 Punkten angesiedelt wurde ($\Delta$ = ca. 4,00 Punkte). Infolgedessen war der Unterschied zwischen \textit{sehr} und \textit{super} ebenso wie zwischen \textit{sehr} und \textit{sau} nicht signifikant. Durch den vergleichsweise hohen Skalenwert unterschied sich \textit{sehr} überdies signifikant von den deskriptiven Intensivierern {\textit{äußerst} (47,40 Punkte)} und \textit{überaus} (46,99 Punkte), dessen denkbare Ursachen oben diskutiert wurden. Ähnliches war innerhalb der expressiven Intensiviererklasse auszumachen, in der sich die beiden höchstpositionierten Ausdrücke {\textit{mega} (63,67 Punkte)} und \textit{sau} (63,69 Punkte) signifikant von den niedrigpositionierten Ausdrücken \textit{super} (57,15 Punkte) und \textit{voll} {(57,23 Punkte)} unterschieden. Die möglichen Gründe hierfür wurden ebenfalls im vorherigen Abschnitt beschrieben.


%Zusammengenommen zeigten die Ergebnisse der paarweisen Vergleiche im Hinblick auf die in Experiment 2b untersuchte Sprechereinstellung einen stark ausgeprägten Unterschied zwischen den Intensiviererklassen.






%Diese werden im Anschluss der Reihe nach erörtert. Begonnen wird mit dem deskriptiven Intensiviererpaar \textit{überaus} (46,99 Punkte) und \textit{äußerst} (47,40 Punkte), das transparente Ausdrücke beinhaltet, deren lexikalischer Ursprung deutlich wahrnehmbar ist {($\Delta$ = ca. 0,41 Punkte)}. Semantisch denotieren die Intensivierer einen bestimmten Punkt, der entweder angenähert oder überschritten wird. Überraschend ist insbesondere der gemittelte Skalenwert von \textit{überaus}, handelte es sich in Experiment 2a dabei doch um denjenigen deskriptiven Intensivierer, dessen Skalaritätsgrad von den Versuchspersonen auf der Skala am höchsten bewertet wurde. Dieses Ergebnis wurde auf eine niedrige Gebrauchsfrequenz zurückgeführt, die sich offenbar in einer höheren Einstufung des Intensivierers auf der Skala niedergeschlagen hat. Daraus lässt sich ableiten, dass der Ausdruck zwar mit einem hohen Grad an Skalarität assoziiert wird, der durch ihn ausgedrückte Grad an Sprechereinstellung allerdings als verhältnismäßig niedrig empfunden wird. Gegenläufig dazu wurde der Skalaritätsgrad von \textit{äußerst} in {Experiment 2a} niedriger eingestuft, was durch eine hohe Gebrauchsfrequenz erklärt wurde. Während sich \textit{überaus} im Korpus mit 10\,534 absoluten Treffern belegen lässt, ist \textit{äußerst} mit 19\,723 Treffern darin fast doppelt so häufig vertreten. Dessen ungeachtet wurden die beiden Ausdrücke in {Experiment 2b} im Hinblick auf Sprechereinstellung auf der Skala vergleichbar evaluiert, was zeigt, dass sie von den Versuchspersonen gleichermaßen als nur wenig einstellungsausdrückend wahrgenommen wurden. Genauso unerwartet war das Ergebnis von \textit{sehr}, dem innerhalb der deskriptiven Intensiviererklasse der höchste Skalenwert zugeschrieben wurde. Dies lässt den Schluss zu, dass \textit{sehr} nicht, wie in der theoretischen Literatur mitunter behauptet (bspw. \citealt[7]{Mueller.1899}; {\citealt[365]{Claudi.2006}}), notwendig nur moderat intensivierend ist.\footnote{Hiervon ist natürlich das aus pragmatischen Gründen reduplizierte \textit{sehr} zu trennen, das ohnehin einen stärkeren Intensitätsgrad ausdrückt als das einfache \textit{sehr} (vgl. u.\,a. {\citealt[16f.]{Hauschild.1899}}; \citealt[266]{Becher.1907}; \citealt[191]{Staffeldt.2018}).} Stattdessen deutet die Einschätzung von \textit{sehr} vielmehr auf einen ungewöhnlich hohen Grad an Sprechereinstellung hin, dessen Ursache möglicherweise in der Opazität des Ausdrucks begründet ist. So legten die Ergebnisse der deskriptiven Intensivierer nahe, dass semantische Transparenz für die kommunizierte Sprechereinstellung eher hemmend ist, während semantische Opazität dagegen mutmaßlich zu einem höheren Skalenwert führt. Innerhalb der expressiven Intensiviererklasse zeichnete sich hinsichtlich Sprechereinstellung eine klare zweiteilige Abstufung ab (vgl. Abbildung \ref{Bar2}). Diese äußerte sich darin, dass das Intensiviererpaar {\textit{super} (57,15 Punkte)} und \textit{voll} {(57,23 Punkte)} von den Versuchspersonen ähnlich bewertet wurde {($\Delta$ = ca. 0,08 Punkte)}, ebenso wie auch \textit{mega} (63,67 Punkte) und \textit{sau} (63,69 Punkte) eine nahezu gleichwertige Skalenposition zugewiesen wurde {($\Delta$ = ca. 0,02 Punkte)}. Während \textit{voll} bereits in Experiment 2a in Bezug auf Skalarität niedrig gerankt wurde, wies \textit{super} erst bei der in {Experiment 2b} untersuchten Sprechereinstellung einen vergleichsweise niedrigen Skalenwert auf. Demnach deuten die Ergebnisse darauf hin, dass \textit{super} nach Einschätzung der Versuchspersonen zwar mit einem mittleren Skalaritätsgrad einhergeht, der durch ihn ausgedrückte Grad an Sprechereinstellung jedoch eher niedrig ist. Dies ist schätzungsweise dem starken Lexikalisierungsgrad und der dementsprechend hohen Gebrauchsfrequenz geschuldet, die dazu führten, dass \textit{super} mit 391,26 Treffern pro Million Textwörter der häufigste Intensivierer im Korpus ist. Dadurch ist es naheliegend, dass der Ausdruck inzwischen als nicht mehr übermäßig expressiv wahrgenommen wird und er in der Folge mit Blick auf Sprechereinstellung von den Versuchspersonen auf der Skala niedrig verortet wurde. Ähnliches ist bei \textit{voll} vorstellbar: Auch wenn der Ausdruck im Korpus lediglich 36,65 Treffer pro Million Textwörter aufweist, ist er aufgrund seines langen Gebrauchs eventuell zu abgenutzt, um als stark expressiv beurteilt zu werden. Das Ergebnis von \textit{voll} steht zudem im Einklang mit den Resultaten des in {Kapitel \ref{Experiment0}} vorgestellten Experiments: Auch darin wurde dem intensivierenden \textit{voll} von den Versuchspersonen weder in Bezug auf Skalarität noch auf Sprechereinstellung ein hoher Grad zugesprochen (vgl. die {Tabellen \ref{SkaWe2} und \ref{SkaWey}}). Neben \textit{super} und \textit{voll} wurden die im Korpus niedrigfrequenten Intensivierer \textit{mega} und \textit{sau} in {Experiment 2b} fast äquivalent bewertet. Dieses Paar bildete die zweite Stufe innerhalb der expressiven Intensiviererklasse. Festzustellen war außerdem, dass beide Ausdrücke schon im Hinblick auf Skalarität ähnlich hoch lokalisiert wurden (vgl. {Tabelle \ref{Skalenwert3}}), was bedeutet, dass sie von den Versuchspersonen sowohl im Kontext von Skalarität als auch von Sprechereinstellung mit einem hohen Grad assoziiert wurden. Dieser Befund liefert weitere Evidenz für die zu testenden Hypothesen, nach denen sich Expressivität in Form eines höheren Grades an Skalarität und Sprechereinstellung manifestiert (vgl. {\citealt[315]{Lang.1983}}; {\citealt[22]{Paradis.1997}}; {\citealt[57]{Ortner.2014}}; \citealt[133]{Gutzmann.2019a}). Bei \textit{mega} wurde der hohe Skalenwert durch die anhaltende Aktualität erklärt, die gegebenenfalls dazu führt, dass der Ausdruck als ausdrucksstärker wahrgenommen wird als bspw. der dem Anschein nach abgenutzte Intensivierer \textit{voll}. Ursächlich für den hohen Skalenwert von \textit{sau} war potenziell die klar erkennbare Semantik und die damit verbundene derbe Konnotation, die bereits bei dem in Kapitel \ref{Experiment0} geschilderten Experiment zu einem vergleichbaren Ergebnis bzw. einem als hoch eingeschätzten Grad an Sprechereinstellung beigetragen hat.
%Angesichts der bereits gewonnenen Ergebnisse ist die Zweiteilung, die sich innerhalb der expressiven Intensiviererklasse beobachten lässt, wenig überraschend und erhärtet diese nachdrücklich.





%\subsubsection{Paarweise Vergleiche der Intensivierer}\label{abs:Intensiviererpaarvergl} $~$ \\
%Die aus den Ergebnissen der Paarvergleiche resultierende Matrix deutete auf den ersten Blick auf eine signifikante Zweiteilung der Intensivierer nach deskriptiv und expressiv hin (vgl. Tabelle \ref{Sign2}). So zeigten nahezu alle Intensivierer signifikante Unterschiede bei den Vergleichen mit der gegenteiligen Intensiviererklasse, d.\,h. deskriptiv vs. expressiv und expressiv vs. deskriptiv. Die einzige Ausnahme stellte der deskriptive Intensivierer \textit{sehr} dar: Der gemittelte Skalenwert dieses Ausdrucks lag mit 53,15 Punkten nahe beim am niedrigsten positionierten expressiven Intensivierer \textit{super}, der auf der Skala im Schnitt bei 57,15 Punkten angesiedelt wurde ($\Delta$ = ca. 4,00 Punkte). Infolgedessen war der Unterschied zwischen \textit{sehr} und \textit{super} ebenso wie zwischen \textit{sehr} und \textit{sau} nicht signifikant. Durch den vergleichsweise hohen Skalenwert unterschied sich \textit{sehr} überdies signifikant von den deskriptiven Intensivierern {\textit{äußerst} (47,40 Punkte)} und \textit{überaus} (46,99 Punkte), dessen denkbare Ursachen oben diskutiert wurden. Ähnliches war innerhalb der expressiven Intensiviererklasse auszumachen, in der sich die beiden höchstpositionierten Ausdrücke {\textit{mega} (63,67 Punkte)} und \textit{sau} (63,69 Punkte) signifikant von den niedrigpositionierten Ausdrücken \textit{super} (57,15 Punkte) und \textit{voll} {(57,23 Punkte)} unterschieden. Die möglichen Gründe hierfür wurden ebenfalls im vorherigen Abschnitt beschrieben.


%Zusammengenommen zeigten die Ergebnisse der paarweisen Vergleiche im Hinblick auf die in Experiment 2b untersuchte Sprechereinstellung einen stark ausgeprägten Unterschied zwischen den Intensiviererklassen.






 





\section{Ergebnisinterpretation}\label{abs:Gesamtdiskussion}
Im Anschluss werden die Ergebnisse der beiden Experimente zusammengefasst, miteinander verzahnt und hinsichtlich der zentralen {Hypothese H6} evaluiert. Wie in Abschnitt \ref{abs:theo} beschrieben, wurde die Untersuchung mit dem Ziel vorgenommen, möglichen Unterschieden zwischen Intensivierern erstmals mit experimentellen Methoden auf den Grund zu gehen. In diesem Zusammenhang wurde die Überprüfung einer Annahme angestrebt, die sich aus den von {\citet[87]{Biedermann.1969}} und \citet[408]{Breindl.2007} postulierten intraklassischen Abstufungen ergibt. Von diesem Standpunkt aus ist anzunehmen, dass sich graduelle Abstufungen in den Stärkegraden von Intensivierern zeigen und sich die Mitglieder einer Intensiviererklasse im Hinblick auf den durch sie ausgedrückten Grad an Skalarität und Sprechereinstellung voneinander unterscheiden. Um die Korrektheit der {Hypothese H6} empirisch zu testen, wurden zwei Rating-Experimente mit parallelem Untersuchungsaufbau durchgeführt: Während {Experiment 2a} den Grad an Skalarität erfragte, zielte {Experiment 2b} auf den Grad an {Sprechereinstellung ab}.


%\subsection{Experiment 2a und Experiment 2b}\label{abs:Gesamtdiskuss}
Sowohl die Ergebnisse von Experiment 2a wie auch die von {Experiment 2b} sprechen in erster Linie klar für einen wahrnehmbaren Unterschied zwischen deskriptiven und expressiven Intensivierern. So positionierten die Versuchspersonen Sätze mit expressiven Intensivierern in beiden Kontexten signifikant höher als Sätze mit deskriptiven Intensivierern. Dies bekräftigt die grundlegende Annahme, dass Expressivität gleichermaßen mit einem höheren Grad an Skalarität und Sprechereinstellung einhergeht (vgl. {\citealt[315]{Lang.1983}}; \citealt[22]{Paradis.1997}; {\citealt[57]{Ortner.2014}}; {\citealt[133]{Gutzmann.2019a}}). Darüber hinaus stimmt es mit den Ergebnissen überein, die durch die in Kapitel \ref{Experiment0} vorgestellte Untersuchung in Bezug auf die darin getesteten Hypothesen H3 und H4 zutage gebracht werden konnten.


Da in dieser Untersuchung vor allem die potenziellen klasseninternen Abstufungen zentral waren, wurden die beiden Intensiviererklassen im nächsten Schritt ausdifferenziert und ihre Mitglieder detailliert begutachtet. Die Klassen werden im Folgenden der Reihe nach inspiziert. Begonnen wird mit den deskriptiven Intensivierern. Für einen besseren Überblick sind die gemittelten Skalenwerte und Standardabweichungen für die {Experimente 2a und 2b} in {Tabelle \ref{deIn}} zusammengefasst.
%Die Anordnung der Intensivierer erfolgt jeweils in ansteigendem Skalenwert.


\begin{table}
\small
%\setlength{\tabcolsep}{8.25pt}
%\begin{tabular}{p{0.15\textwidth} p{0.1\textwidth} p{0.15\textwidth} p{0.1\textwidth} p{0.15\textwidth}}
\begin{tabularx}{\textwidth}{X@{} lrr lrr }
\lsptoprule
 &  \multicolumn{6}{c}{\textbf{Deskriptive Intensivierer}}  \\
\cmidrule(lr){2-7}\cmidrule(lr){5-7}
& \multicolumn{3}{c}{{Skalarität}} & \multicolumn{3}{c}{{Sprechereinstellung}}  \\
\cmidrule(lr){2-4}\cmidrule(lr){5-7}
{Stärkegrad}  & Intensivierer & Skalenwert& \textit{SD}  & Intensivierer & Skalenwert& \textit{SD}   \\
\midrule
schwach & \textit{ziemlich} & 51,22& 23,6 & \textit{überaus} & 46,99 & 27,4  \\
 & & &  &    \textit{äußerst} & 47,40& 25,7 \\
mittel & \textit{äußerst} & 59,46& 23,4  & \textit{ziemlich} & 49,06& 28,2 \\
 & \textit{sehr} & 59,69& 23,3 &  & \\
stark & \textit{überaus} & 62,36& 22,0 & \textit{sehr} & 53,15 & 29,6 \\
\lspbottomrule
\end{tabularx}
\caption[Gemittelte Skalenwerte und Standardabweichungen der deskriptiven Intensivierer bei den {Experimenten 2a und 2b}]{Gemittelte Skalenwerte und Standardabweichungen der deskriptiven Intensivierer bei den Experimenten 2a und 2b (angeordnet nach ansteigendem Skalenwert).}
\label{deIn}
\end{table}

\noindent Den Werten aus Tabelle \ref{deIn} nach lassen sich die deskriptiven Intensivierer sowohl bei der Untersuchung von Skalarität als auch von Sprechereinstellung grundsätzlich in drei Stärkegrade einteilen, was mit der zu testenden {Hypothese H6} übereinstimmt. Von den gemittelten Skalenwerten ausgehend können sie somit unter Berücksichtigung beider Merkmale als schwach, mittel oder stark klassifiziert werden, wobei die beobachtete Klassifikation partiell der in der theoretischen Literatur vertretenen Einschätzung entspricht (vgl. {\citealt[87]{Biedermann.1969}} und \citealt[408]{Breindl.2007}). Dass sich die Rangfolge der Abstufungen in Abhängigkeit des untersuchten Merkmals, d.\,h. Skalarität oder Sprechereinstellung, unterscheidet, deutet zudem darauf hin, dass die Versuchspersonen sie als verschieden wahrgenommen und dementsprechend unterschiedlich evaluiert haben. Dieses Ergebnis gibt Hinweise darauf, dass das Skalenrating wie intendiert funktioniert hat. Auffällig ist außerdem die Skaleneinschätzung bei der in Experiment 2b abgefragten Sprechereinstellung. Wie in {Abschnitt \ref{abs:Diskussion2}} erläutert, wurden die deskriptiven Intensivierer in diesem Kontext nicht so niedrig bewertet, wie man es aufgrund der fehlenden Expressivität eigentlich erwarten würde. So ist nach der Definition von {\citet[133]{Gutzmann.2019a}} davon auszugehen, dass die deskriptiven Intensivierer bezüglich Sprechereinstellung weitgehend am unteren Ende der Skala positioniert werden. Diese Mutmaßung hat sich durch die Untersuchung jedoch nicht bestätigt. Infolgedessen lässt das Ergebnis den Schluss zu, dass auch deskriptive Intensivierer zu einem gewissen Grad als einstellungsausdrückend wahrgenommen werden, weswegen ihnen in {Experiment 2b} von den Versuchspersonen ein vergleichsweise hoher Grad an Sprechereinstellung zugesprochen wurde. Im Unterschied zu expressiven Intensivierern kann die Expressivität bei diesen Ausdrücken nicht als Erklärung für die hohe Skaleneinschätzung in Erwägung gezogen werden, sodass deren tatsächlichen Gründe an dieser Stelle nicht zweifelsfrei abzuschätzen sind. Dessen ungeachtet stellt der als hoch eingeschätzte Grad an Sprechereinstellung ein unerwartetes und bemerkenswertes Ergebnis dar.


Als Nächstes werden die expressiven Intensivierer besprochen, deren gemittelten Skalenwerte und Standardabweichungen in {Tabelle \ref{exIn} gelistet sind.}


% (vgl. Tabelle \ref{exIn}).

%. Deren gemittelten Skalenwerte und Standardabweichungen sind für die beiden Experimente 2a und 2b in Tabelle \ref{exIn} aufgeführt. 
%Auch hier sind die Ausdrücke jeweils in ansteigendem Skalenwert gelistet.

%
\begin{table}
\small
%\setlength{\tabcolsep}{8.25pt}
%\begin{tabular}{p{0.15\textwidth} p{0.1\textwidth} p{0.15\textwidth} p{0.1\textwidth} p{0.15\textwidth}}
\begin{tabularx}{\textwidth}{X@{} lrr lrr }
\lsptoprule
 &  \multicolumn{6}{c}{\textbf{Expressive Intensivierer}}  \\
\cmidrule(lr){2-7}\cmidrule(lr){5-7}
& \multicolumn{3}{c}{{Skalarität}} & \multicolumn{3}{c}{{Sprechereinstellung}}  \\
\cmidrule(lr){2-4}\cmidrule(lr){5-7}
{Stärkegrad}  & Intensivierer & Skalenwert& \textit{SD}  & Intensivierer & Skalenwert& \textit{SD}   \\
\midrule

schwach & \textit{voll} & 56,38& 23,4  & \textit{super} & 57,15& 28,2   \\
 &  &  & & \textit{voll} & 57,23 & 27,8\\
\multicolumn{1}{l}{mittel} & \textit{super} & 63,15& 22,2 & \textit{} & \\
  \multicolumn{1}{l}{} & \textit{sau} & 63,87& 24,4 & & \\
stark & \textit{mega} & 67,27& 24,4  & \textit{mega} & 63,67& 28,0   \\
  &  & &  & \textit{sau} & 63,69 & 29,4  \\
\lspbottomrule
\end{tabularx}
\caption[Gemittelte Skalenwerte und Standardabweichungen der expressiven Intensivierer bei den {Experimenten 2a und 2b}]{Gemittelte Skalenwerte und Standardabweichungen der expressiven Intensivierer bei den Experimenten 2a und 2b (angeordnet nach ansteigendem Skalenwert).}
\label{exIn}
\end{table}

Die Werte in Tabelle \ref{exIn} deuten auf graduelle Abstufungen zwischen den untersuchten expressiven Intensivierern hin. So lassen sich diese Ausdrücke in {Experiment 2a} in Abhängigkeit ihres Skalaritätsgrades als schwach, mittel oder stark klassifizieren. Bei der in {Experiment 2b} erfragten {Sprechereinstellung} können sie zu zwei Stufen mit je zwei Mitgliedern zusammengefasst werden, die in der Tabelle als schwach und stark gekennzeichnet sind.\footnote{Ob man den Stärkegrad von \textit{mega} und \textit{sau} als mittel oder stark klassifiziert, spielt meines Erachtens für die Interpretation der Daten keine Rolle. Um den Unterschied zwischen den beiden Stärkegruppen deutlicher zu machen, habe ich mich hier für `stark' entschieden. Außerdem sei erwähnt, dass die Abstufungen zum Teil nur marginal sind und die Klassenbildung nicht in jedem Fall offensichtlich ist, die Klassifikation also keineswegs als trennscharf anzusehen ist. Nichtsdestotrotz habe ich die Ausdrücke in Abhängigkeit ihrer Bepunktung und somit in Bezug auf die relative Nähe zu ihren Gruppenmitgliedern eingeordnet.} Dieser Befund liefert die erste empirische Evidenz für die von \citet[87]{Biedermann.1969} und {\citet[408]{Breindl.2007}} angenommenen klasseninternen Unterschiede und untermauert die zur Diskussion stehende {Hypothese H6}.


In Anbetracht der Einschätzung des Skalaritätsgrades expressiver Intensivierer ist jedoch ebenso wie bei dem in Kapitel \ref{Experiment0} vorgestellten {Experiment 1} zu bedenken, dass die beiden Bedeutungsdimensionen in einer solchen experimentellen Erhebung methodisch nicht trennbar sind. Da sich der expressive Gehalt eines Intensivierers womöglich auch hier auf die Bewertung seiner skalaren Grundbedeutung ausgewirkt hat, sind realistische Aussagen über die zu testende Hypothese H3 im Sinne von \citet[22]{Paradis.1997} und {\citet[133]{Gutzmann.2019a}} nur eingeschränkt möglich. Aus diesem Grund ist das tatsächliche Ausmaß der Auswirkung von Expressivität trotz der beobachteten Daten auch nach der Untersuchung nicht abschließend geklärt.


Zusätzlich wurde im Rahmen der Datenauswertung untersucht, ob es einen statistisch signifikanten Zusammenhang zwischen der in {Experiment 2a} abgefragten Skalarität und der in Experiment 2b fokussierten Sprechereinstellung gibt. Das Ergebnis der Korrelationsanalyse ist nicht signifikant, was bedeutet, dass die Variablen \textsc{Skalarität} und \textsc{Sprechereinstellung} in diesem Experiment einander nicht bedingen. Dies überrascht insoweit, als nach der Definition von {\citet[133]{Gutzmann.2019a}} von einer Korrelation auszugehen ist, nimmt er doch an, dass ein erhöhter Grad an Sprechereinstellung gleichsam von einem höheren Grad an Skalarität begleitet wird. Dass die erwartete Korrelation an dieser Stelle nicht nachgewiesen werden kann, ist jedoch mit Vorsicht zu interpretieren. Aufgrund des Umstands, dass es sich bei den Versuchspersonen von Experiment 2a und {Experiment 2b} um zwei unterschiedliche Personengruppen handelte, wurde die Korrelationsanalyse nicht auf Grundlage der einzelnen Ratings, sondern auf Grundlage der Mittelwerte der Intensivierer vorgenommen. Darüber hinaus  wurden in die Analyse auch lediglich die Mittelwerte der vier expressiven Intensivierer miteinbezogen, während die Mittelwerte der deskriptiven Intensivierer aufgrund der fehlenden Expressivität zur Gänze außen vor gelassen wurden. Eine Korrelationsanalyse mit bloß vier Werten ist aber nur eingeschränkt aussagekräftig: Nur weil die beobachteten Daten keine Rückschlüsse auf einen Zusammenhang zwischen den Merkmalen zulassen, bedeutet dies nicht, dass es grundsätzlich keinen gibt. Vor diesem Hintergrund ist also vorstellbar, dass die Analyse mit einer größeren Datenmenge durchaus zu einer signifikanten Korrelation geführt hätte. Dies lässt sich vom jetzigen Standpunkt aus indes nicht evaluieren, sodass ihre weitere Untersuchung nachfolgenden Experimenten vorbehalten bleiben muss.



%dass die untersuchten deskriptiven Intensivierer im Kontext der in Experiment 2a erfragten Skalarität von den Versuchspersonen auf der Skala höher verortet wurden als im Kontext der in Experiment 2b ermittelten Sprechereinstellung. Demzufolge wurden sie insgesamt vielmehr als skalierend denn als einstellungsausdrückend wahrgenommen, was auf den ersten Blick wenig überrascht. Nichtsdestotrotz wurden die Ausdrücke bezüglich Sprechereinstellung nicht so niedrig bewertet, wie man es aufgrund der fehlenden Expressivität eigentlich erwarten würde. So war nach der Definition von {\citet[133]{Gutzmann.2019a}} davon auszugehen, dass die deskriptiven Intensivierer in Experiment 2b weitgehend am unteren Ende der Skala positioniert werden. Diese Vermutung hat sich durch die Untersuchung jedoch nicht bestätigt. Infolgedessen lässt das Ergebnis den Schluss zu, dass auch deskriptive Intensivierer zu einem gewissen Grad als einstellungsausdrückend bzw. bewertend wahrgenommen werden, weswegen ihnen in {Experiment 2b} von den Versuchspersonen ein vergleichsweise hoher Grad an Sprechereinstellung zugesprochen wurde. Im Unterschied zu expressiven Intensivierern kann die Expressivität bei diesen Ausdrücken nicht als Erklärung für die hohe Skaleneinschätzung infrage kommen, sodass deren tatsächlichen Gründe an dieser Stelle nicht zweifelsfrei abzuschätzen sind. Dessen ungeachtet stellt der als hoch eingeschätzte Grad an Sprechereinstellung ein unerwartetes und durchaus bemerkenswertes Ergebnis dar. Weiterhin war festzustellen, dass sich die deskriptiven Intensivierer sowohl bei der Untersuchung zu Skalarität als auch zu Sprechereinstellung grundsätzlich in drei Stärkegrade einteilen ließen, was mit der zu testenden {Hypothese H6} übereinstimmt, die ebensolche klasseninternen Abstufungen annimmt. Von den gemittelten Skalenwerten ausgehend konnten die deskriptiven Intensivierer somit unter Berücksichtigung beider Merkmale als schwach, mittel oder stark klassifiziert werden, wobei die beobachtete Klassifikation partiell ebenso der in der theoretischen Literatur vertretenen Einschätzung entsprach (vgl. bspw. \citet[87]{Biedermann.1969} und \citet[408]{Breindl.2007}). Dass sich die Rangfolge der Abstufungen in Abhängigkeit des untersuchten Merkmals, d.\,h. Skalarität oder Sprechereinstellung, unterschied, deutete zudem darauf hin, dass sie die Versuchspersonen als verschieden wahrgenommen und dementsprechend unterschiedlich evaluiert haben. Dieses Ergebnis lieferte Hinweise darauf, dass das Skalenrating wie intendiert funktioniert hat.





%Diesbezüglich fiel bei Experiment 2a auf, dass das Ergebnis von sechs der insgesamt acht Intensivierer insofern erwartungsgemäß war, als sie eine klare Trennung zwischen deskriptiver und expressiver Klasse anzeigten. Anders verhielten sich dagegen das deskriptive \textit{überaus} und das expressive \textit{voll}: Während \textit{überaus} einen vergleichsweise hohen Skalenwert aufwies (62,36 Punkte), lag der von \textit{voll} deutlich darunter (56,38 Punkte). Dieses Ergebnis wurde auf unterschiedliche Interpretationsmöglichkeiten der {Hypothese H4} zurückgeführt: Die erste, starke Interpretation fußte auf der Annahme, dass der stärkste deskriptive Intensivierer grundsätzlich über einen niedrigeren Skalenwert verfügt als der schwächste expressive Intensivierer, es also eine klare Trennung der Ausdrücke nach deskriptiv und expressiv gibt. Entsprechend wäre von Abstufungen in den Stärkegraden auszugehen, die unabhängig von der Intensiviererklasse sind. Nach dieser Interpretation würde es sich bei \textit{überaus} und \textit{voll} um Ausreißer handeln, die sich wider Erwarten anders verhalten als die restlichen Gruppenmitglieder. Eine Erklärung für dieses abweichende Verhalten wurde in den Gebrauchsfrequenzen von \textit{überaus} und \textit{voll} gesehen, die gegebenenfalls zum eingeschätzten Skalaritätsgrad beigetragen haben. Denkbar ist jedoch auch eine gemäßigtere Interpretation, die eine Tripartition des mit den Ausdrücken verbundenen Skalenwerts annimmt und von schwach über mittel bis hin zu stark reicht. Nach dieser können die deskriptiven Intensivierer in Bezug auf Skalarität schwach, mittel oder stark sein, ebenso wie auch die expressiven Intensivierer dahingehend schwach, mittel oder stark sein können. Beide Intensiviererklassen würden gleichrangig nebeneinander stehen. Durch parallel verschiebbare Grenzen wären Überlappungen möglich, wie in dem Fall bei \textit{überaus} und \textit{voll} zu verzeichnen ist. Die Einzelbetrachtung von {Experiment 2b} stützte dagegen die zuvor getroffene Beobachtung in Gänze, denn auch hier zeichnete sich eine klare Trennung der Intensiviererklassen ab. So wurde der hinsichtlich Sprechereinstellung am stärksten bewertete deskriptive Intensivierer {\textit{sehr} (53,15 Punkte)} auf der Skala im Schnitt 4,00 Punkte niedriger verortet als der am schwächsten evaluierte expressive Intensivierer \textit{super} (57,15 Punkte).


%Im Rahmen der inferenzstatistischen Datenanalyse wurden die Intensivierer mittels \textit{t}-Tests paarweise miteinander verglichen, um zu sehen, ob die beobachteten Unterschiede statistisch signifikant waren. In diesem Zusammenhang war bei der in Experiment 2a erfragten Skalarität festzustellen, dass sich vordergründig drei Intensivierer klar von den anderen unterschieden: der deskriptive Intensivierer \textit{ziemlich} sowie die expressiven Intensivierer \textit{mega} und \textit{voll}. Das Ergebnis von \textit{ziemlich} war insofern wenig überraschend, als es auf der Skala in Bezug auf Skalarität am niedrigsten positioniert wurde (51,22 Punkte). Erklärt wurde der verhältnismäßig niedrige Skalenwert durch die Semantik des Ausdrucks, die kein Höchstmaß der bezeichneten Eigenschaft, sondern einen eher gemäßigten Grad ausdrückt. {\citet[7]{Nouwen.2020}} führt ebendiesen gemäßigteren Intensitätsgrad auf die inhärent positive Polarität des Ausdrucks zurück, durch die die Verstärkung in Form einer moderaten Annäherung an einen bestimmten Punkt realisiert wird. Aufseiten der expressiven Intensivierer zeigte \textit{mega} signifikante Unterschiede beim Vergleich mit sechs der restlichen Ausdrücke, was mit der anhaltenden Aktualität und dem dadurch als hoch eingeschätzten Skalaritätsgrad erklärt wurde (67,27 Punkte). Dies lieferte außerdem weitere Indizien dafür, dass ein neuer und wenig frequenter Intensivierer als expressiver evaluiert wird als ein älterer und frequenter Intensivierer (vgl. H1-Temp und H1-Freq). Wie oben beschrieben, deckt sich dieses Ergebnis mit dem von {\citet[]{Calpestrati.2017}}. Er beobachtet im Rahmen seiner vergleichenden Studie zu deutschen und italienischen Intensivierern, dass \textit{super} (neben \textit{extra} und \textit{maxi}) von den Versuchspersonen mit Blick auf Intensität als am schwächsten eingestuft wird, während \textit{mega} (neben \textit{hyper}) mit einem hohen Wert in Verbindung gebracht wird. Ebenso auffällig war das Ergebnis der Paarvergleiche von \textit{voll}, die in fünf von sieben Fällen signifikant waren. Diesbezüglich wurde gemutmaßt, dass der Intensivierer aufgrund seines langen Gebrauchs schon zu abgenutzt ist, was möglicherweise den niedrigen Skalenwert bedingt hat (56,38 Punkte). Auf die gleiche Weise wurden auch die im Rahmen von Experiment 2b ermittelten Skalenwerte zur Sprechereinstellung mehreren Paarvergleichen mittels \textit{t}-Tests unterzogen. Anhand der Ergebnisse war zu erkennen, dass sich die beiden Intensiviererklassen nahezu vollständig voneinander unterschieden. So zeigten mit Ausnahme von \textit{sehr} alle Ausdrücke signifikante Unterschiede beim Vergleich mit der gegenteiligen Intensiviererklasse an. Das Ergebnis von \textit{sehr} wurde auf dessen hohen Lexikalisierungsgrad zurückgeführt, der sich anscheinend in einer vergleichsweise hohen Skalenpositionierung niedergeschlagen hat.


%Weiterhin wurden für beide Experimente lineare gemischte Modelle berechnet. Dem Regressionsmodell zu Experiment 2a nach übten die Prädiktoren \textsc{Intensiviererklasse} und \textsc{Position} sowie deren Zwei-Wege-Interaktion Einfluss auf den Skalenwert aus. Der Unterschied zwischen den deskriptiven und expressiven Intensivierern, der sich in Bezug auf Skalarität abzeichnete, war folglich statistisch signifikant ($\chi$$^{2}$\,(1) = 8,81, \textit{p} < 0,01). Hinzu kam ein signifikanter Positionseffekt, nach dem ein Item mit fortschreitender Position auf der Skala niedriger verortet wurde ($\chi$$^{2}$\,(1) = 16,20, \textit{p} < 0,001). Die signifikante Interaktion aus beiden Prädiktoren deutete obendrein auf einen Gewöhnungseffekt hin, demgemäß sich die eingeschätzten Skalenwerte der Intensiviererklassen mit fortschreitender Position einander annäherten ({$\chi$$^{2}$\,(1) = 4,17}, {\textit{p} < 0,05}). Ein ähnliches Ergebnis zeigte sich bei dem im Rahmen von {Experiment 2b} berechneten Modell zu Sprechereinstellung: Auch hier wirkte sich der Prädiktor \textsc{Intensiviererklasse} signifikant auf den Skalenwert aus {($\chi$$^{2}$\,(1) = 14,51}, \textit{p} < 0,001). Der festgestellte Unterschied zwischen deskriptiven und expressiven Intensivierern war also nicht zufallsbedingt.


%Abschließend lässt sich resümieren, dass die Ergebnisse der Untersuchung die Gültigkeit der Hypothesen H4 und H5 erhärten. Entsprechend ist davon auszugehen, dass expressive Intensivierer \textit{de facto} gleichermaßen mit einem höheren Grad an Skalarität und Sprechereinstellung einhergehen als deskriptive Intensivierer. Dies liefert die erste empirische Evidenz für die theoretischen Annahmen von {\citet[22]{Paradis.1997}} und {\citet[133]{Gutzmann.2019a}} sowie {\citet[315]{Lang.1983}} und {\citet[57]{Ortner.2014}}.



%unterschied sich die Einschätzung der expressiven Intensivierer stark von der der deskriptiven Intensivierer. Während letztere im Hinblick auf Skalarität höher bewertet wurden als im Hinblick auf Sprechereinstellung, zeichnete sich bei den expressiven Intensivierern dahingehend ein heterogenes Muster ab. So wurde \textit{voll} im Kontext von Sprechereinstellung ein höherer Grad zugeschrieben, wohingegen der eingeschätzte Skalenwert von \textit{sau} bei beiden Merkmalen beinahe identisch war. Anders sah dies bei \textit{super} und \textit{mega} aus, deren Grad an Skalarität den Grad an Sprechereinstellung überragte. Darauf basierend stellt sich die Frage, inwiefern die beobachteten Ergebnisse zu der zu überprüfenden Hypothese passen. Wie mehrfach beschrieben, nimmt {\citet[133]{Gutzmann.2019a}} an, dass die Expressivität eines Ausdrucks eine Anhebung seines Skalaritätsgrades nach sich zieht. Insofern wäre zu erwarten, dass der Grad an Skalarität bei expressiven Intensivierern zwar als hoch evaluiert wird, er jedoch nicht an den Grad an Sprechereinstellung heranreicht, letzterer also prinzipiell als stärker empfunden wird. Dies war bei \textit{super} und \textit{mega} augenscheinlich nicht der Fall. Liefern die Daten dieser Ausdrücke möglicherweise Gegenevidenz, die Gutzmanns Hypothese entkräftet? Diese Frage lässt sich nicht zweifelsfrei beantworten. Dagegen wurde der Grad an Sprechereinstellung bei \textit{voll} höher eingeschätzt als der Grad an Skalarität. Obwohl der Ausdruck als nicht übermäßig expressiv empfunden wurde, schien er gewissermaßen mehr einstellungsausdrückend denn als skalierend empfunden worden zu sein. Inwiefern die Expressivität des Intensivierers einen Anteil an der Einschätzung seines Skalaritätsgrades hatte, lässt sich auf Grundlage der Daten allerdings nicht abschätzen. So war es in Experiment 2a nicht möglich, die Skalarität ohne jeglichen Einfluss der expressiven Komponente zu erheben und beide Bedeutungsdimensionen isoliert zu untersuchen. Dies bedeutet, dass die skalare Bedeutung bei einem expressiven Intensivierer experimentell nicht als solche identifizierbar ist und die Expressivität unweigerlich auf die Beurteilung der Skalarität einwirkt. Insofern bleibt festzuhalten, dass das Untersuchungsdesign eine gute -- und die bislang {einzige --} Approximation zur Erhebung der Grade beider Bedeutungsdimensionen darstellt, gleichwohl keine isolierte Betrachtung ebendieser erfolgen kann. Auffällig war außerdem, dass \textit{sau} in beiden Kontexten von den Versuchspersonen nahezu gleichwertig wahrgenommen wurde. Dabei handelt es sich bei {\citet[133]{Gutzmann.2019a}} um den klassischen Stellvertreter der expressiven Intensiviererklasse. Aus dieser Sichtweise ist es überraschend, dass der Ausdruck sowohl bei Skalarität als auch bei Sprechereinstellung auf der Skala beinahe identisch verortet wurde. Doch auch hierbei stellt sich wiederum die Frage, welcher Anteil rein auf Skalarität zukommt und wieviel der Expressivität des Ausdrucks geschuldet ist, was sich, wie soeben erläutert, nicht ohne Weiteres beantworten lässt. Als weiteres Ergebnis konnten die gemäß der {Hypothese H6} vermuteten graduellen Abstufungen beobachtet werden. So ließen sich die expressiven Intensivierer in Experiment 2a in Abhängigkeit ihres Skalaritätsgrades als schwach, mittel oder stark klassifizieren. In Experiment 2b konnten sie zu zwei Stufen mit je zwei Mitgliedern zusammengefasst werden, die in der Tabelle als schwach und stark gekennzeichnet waren. Dieser Befund liefert die erste empirische Evidenz für die von \citet[87]{Biedermann.1969} und {\citet[408]{Breindl.2007}} angenommenen klasseninternen Unterschiede und untermauert die zur Diskussion stehende {Hypothese H6}.





%\subsection[Korrelation Skalarität und Sprechereinstellung]{Sind Skalarität und Sprechereinstellung korreliert?}   \label{ZuSa}




%Dieses Resultat lässt unterschiedliche Interpretationen zu. Die erste Interpretation bezieht sich auf die im Experiment angewandte Methode: Da das Ergebnis nicht so klar ausfällt, wie auf Grundlage der Annahmen zu erwarten wäre, ist vorstellbar, dass die Unkorreliertheit einem methodischen Defizit geschuldet ist. So war der beobachtete Unterschied zwischen den beiden Intensiviererklassen trotz statistischer Signifikanz nicht so stark ausgeprägt, wie infolge der Klassifikation der Intensivierer zu erwarten war. In diesem Zusammenhang besteht die Möglichkeit, dass die Versuchspersonen den Skalaritätsgrad der deskriptiven Intensivierer fälschlicherweise als Sprechereinstellung interpretiert haben, was sich eventuell in einem verhältnismäßig hohen Skalenwert niedergeschlagen hat. Der Skalenwert könnte aber auch durch die Aufgabenstellung bedingt gewesen sein, wurde in Experiment 2b doch \textit{expressis verbis} nach dem Grad an Bewertung statt nach dem Grad an Sprechereinstellung wie in dem in Kapitel \ref{Experiment0} dargelegten Experiment gefragt. Dass es sich beim Ergebnis um ein Artefakt der Daten handelt, ist in Anbetracht der hohen Teilnehmerzahl auszuschließen. Eine andere Interpretation betrifft die Klassifikation der Intensivierer. Vor diesem Hintergrund ist anzunehmen, dass die vorgenommene Klasseneinteilung nach deskriptiv und expressiv nur bedingt korrekt war. Demnach ist durchaus vorstellbar, dass der als expressiv klassifizierte Intensivierer \textit{super} aufgrund seines hohen Lexikalisierungsgrades und der damit einhergehenden hohen Frequenz inzwischen an Expressivität eingebüßt hat. Gleichsam ist der ebenso als expressiv eingeschätzte Intensivierer \textit{voll} durch seinen persistenten Gebrauch vermutlich zu abgenutzt und daher nicht mehr ausdrucksstark genug, um als expressiv empfunden zu werden. Entsprechend ist denkbar, dass die Untersuchung mit anderen, eventuell auch aktuelleren Intensivierern an dieser Stelle zu anderen Ergebnissen geführt hätte. Die Frage nach den tatsächlichen Gründen für die Skaleneinschätzung der besagten Intensivierer, lässt sich allerdings nicht zweifelsfrei beantworten und bleibt somit über die Untersuchung hinaus offen.

 


%\subsection[Pseudowörter versus reale Wörter]{Pseudowörter versus reale Wörter}   \label{Pseure}
%Außerdem wurde untersucht, ob es zwischen den Mittelwerten der im vorangegangenen Experiment 1 (vgl. Kapitel \ref{Experiment0}) und der in Experiment 2 verwendeten realen Wörtern einen Unterschied gab. In diesem Zusammenhang wurden die gemittelten Skalenwerte der Schnittmenge bestehend aus fünf Intensivierern (d.\,h. \textit{mega}, \textit{sau}, \textit{sehr}, \textit{super} und \textit{voll}) einem \textit{t}-Test unterzogen, mit dessen Hilfe die Mittelwerte zweier abhängiger Stichproben auf statistische Signifikanz hin überprüft werden kann {\citep[vgl.][16]{Sauer.2019}}. Das Ergebnis deutete im Hinblick auf beide Merkmale zunächst auf einen signifikanten Unterschied hin: Skalarität (\textit{t}\,(4) = 34,44, \textit{p} < 0,001) und Sprechereinstellung (\textit{t}\,(4) = 19,98, \textit{p} < 0,001). Dabei war jedoch auffällig, dass sich insbesondere die Skalenwerte des deskriptiven Intensivierers \textit{sehr} bei beiden Experimenten stark voneinander unterschieden, während bei den anderen Ausdrücken keine derart großen Unterschiede zu verzeichnen waren. In der Folge wurde \textit{sehr} vom \textit{t}-Test ausgeschlossen, sodass nur noch die Mittelwerte der expressiven Intensivierer miteinander kontrastiert wurden. Das Ergebnis dieses Tests war in beiden Kontexten nicht signifikant: Skalarität ({\textit{t}\,(3) = 28,12}, {\textit{p} > 0,05}) und Sprechereinstellung (\textit{t}\,(3) = 15,54, {\textit{p} > 0,05}). Somit ist davon auszugehen, dass es für die expressiven Intensivierer absolut unerheblich war, ob sie im Experiment mit Pseudowörtern oder realen Wörtern kombiniert wurden. Daran schließt sich jedoch die Frage an, wieso dies beim deskriptiven Intensivierer \textit{sehr} offenbar anders war, wies dieser doch im Kontext der Pseudowörter einen durchschnittlichen Skalenwert von 63,90 Punkten bei Skalarität und 42,34 Punkten bei Sprechereinstellung sowie im Kontext der realen Wörter einen durchschnittlichen Skalenwert von 59,69 Punkten bei Skalarität und 53,15 Punkten bei Sprechereinstellung auf. Demnach unterschied sich die Skaleneinschätzung von \textit{sehr} in Abhängigkeit der verwendeten Wörter vor allem bei den Untersuchungen zur Sprechereinstellung merklich, was sich möglicherweise auf die bereits umschriebenen Schwierigkeiten beim Verständnis der Aufgabenstellung zurückführen lässt. Den Ergebnissen von {Experiment 2b} nach ist möglich, dass die Versuchspersonen die Skalarität der deskriptiven Intensivierer fälschlicherweise als Sprechereinstellung aufgefasst und den Ausdrücken in der Konsequenz diesbezüglich einen vergleichsweise hohen Skalenwert zugeschrieben haben. Dieses Problem schien bei dem in {Kapitel \ref{Experiment0}} vorgestellten Experiment 1 aufgrund geänderter Fragestellung nicht aufgetreten zu sein, sodass \textit{sehr} in diesem auf der Skala erheblich niedriger positioniert wurde als im Kontext der realen Wörter.\footnote{An dieser Stelle sei erneut auf die geänderte Reihenfolge der Präsentation der Experimente hingewiesen, wurde das in Kapitel \ref{Experiment0} geschilderte Experiment in der Realität im Anschluss an die in diesem Kapitel vorgestellte Untersuchung durchgeführt und konnte in Bezug auf die Fragestellung optimiert werden.}





\section{Fazit}\label{abs:Fazit1025}
Die in diesem Kapitel präsentierte Untersuchung zielte darauf ab, eine weiterführende Hypothese zu testen, die sich implizit aus der Annahme eines wahrnehmbaren Unterschiedes zwischen deskriptiven und expressiven Intensivierern ergibt. Auf Basis der theoretischen Literatur (vgl. {\citealt[87]{Biedermann.1969}}; \citealt[408]{Breindl.2007}) folgt, dass es innerhalb der Intensiviererklassen Unterschiede bzw. graduelle Abstufungen in Bezug auf den durch die Intensivierer ausgedrückten Intensitätsgrad gibt und sich die Ausdrücke in Abhängigkeit ihrer Semantik diesbezüglich skalar anordnen lassen. Zur Überprüfung der Hypothese H6 wurden zwei getrennte Experimente mit parallelem Aufbau und identischem lexikalischen {Material} durchgeführt. Wurde in {Experiment 2a} der Grad an Skalarität erhoben, erfragte {Experiment 2b} den Grad an Sprechereinstellung. In der Untersuchung wurde eine Stichprobe aus jeweils vier deskriptiven und expressiven Intensivierern mit unterschiedlich wahrnehmbarer Intensität mit Blick auf Skalarität und Sprechereinstellung kontrastiert. Ihr Ziel lag einerseits darin, dem Unterschied zwischen deskriptiven und expressiven Intensivierern zum ersten Mal experimentell nachzugehen und zu überprüfen, ob letztere hinsichtlich Skalarität und Sprechereinstellung als stärker wahrgenommen werden, wie von zahlreichen Autor:innen postuliert wird (vgl. {\citealt[315]{Lang.1983}}; {\citealt[22]{Paradis.1997}}; {\citealt[57]{Ortner.2014}}; {\citealt[133]{Gutzmann.2019a}}). Andererseits wurden die vermuteten klasseninternen Abstufungen getestet und Rückschlüsse auf das Zutreffen der oben beschriebenen {Hypothese H6} nach {\citet[87]{Biedermann.1969}} und \citet[408]{Breindl.2007} gezogen.


Experiment 2a diente dazu, den Skalaritätsgrad der Intensivierer zu ermitteln, um festzustellen, ob expressive Intensivierer, wie nach \citet[22]{Paradis.1997} und \citet[133]{Gutzmann.2019a} zu erwarten ist, dahingehend als stärker beurteilt werden. Das Ergebnis weist auf die Gültigkeit dieser Annahme hin, werden Sätze mit expressiven Intensivierern auf der Skala doch signifikant höher evaluiert als Sätze, die deskriptive Intensivierer beinhalten. Daneben zeigt sich sowohl innerhalb der deskriptiven wie auch der expressiven Intensiviererklasse eine dreiteilige Abstufung, gemäß derer sich die Ausdrücke als schwach, mittel oder stark skalierend klassifizieren lassen. Dies kann als erste empirische Evidenz für die Korrektheit der {Hypothese H6} gewertet werden.


Im Rahmen des parallel aufgebauten Experiments 2b wurde der durch die Intensivierer ausgedrückte Grad an Sprechereinstellung bestimmt. Auch hier wurde auf Grundlage der theoretischen Literatur angenommen, dass sich beide Intensiviererklassen diesbezüglich auf der Skala signifikant voneinander unterscheiden (vgl. \citealt[315]{Lang.1983}; {\citealt[57]{Ortner.2014}}). Ausgehend vom Ergebnis der Untersuchung kann diese Vermutung für korrekt befunden werden: Sätze mit expressiven Intensivierern werden im Hinblick auf Sprechereinstellung auf der Skala signifikant höher verortet als Sätze mit deskriptiven Intensivierern. Erneut sind klasseninterne Abstufungen zu beobachten, in deren Folge die Ausdrücke als unterschiedlich stark einstellungsausdrückend klassifiziert werden können. In diesem Zusammenhang stellte sich die Frage, ob sich die beobachtete Rangfolge der expressiven Intensivierer mit der Rangfolge von Experiment 2a deckt, was allerdings bloß bedingt der Fall ist. Der Unterschied in der Wahrnehmung spricht dafür, dass das Skalenrating als Untersuchungsmethode funktioniert hat.


Ein weiteres Ziel der Untersuchung lag in der Überprüfung des potenziellen Zusammenhangs zwischen Skalarität und Sprechereinstellung. Wenngleich die Experimente separat durchgeführt, beide Merkmale also getrennt abgefragt wurden, konnten die Ergebnisse final zusammengebracht und miteinander in Beziehung gesetzt werden. Durch dieses Vorgehen wurde getestet, ob es einen korrelativen Zusammenhang zwischen den Variablen {\textsc{Skalarität}} und \textsc{Sprechereinstellung} gibt, wie \citet[133]{Gutzmann.2019a} nahelegt. Wie an entsprechender Stelle beschrieben, findet sich keine Korrelation, was in Anbetracht der geringen Datenmenge wenig überrascht. So wurden in die Analyse lediglich die Mittelwerte der expressiven Intensivierer miteinbezogen, wohingegen die Mittelwerte der deskriptiven Intensivierer aufgrund fehlender Expressivität außen vor gelassen wurden. Dadurch ist die Aussagekraft der Ergebnisse eingeschränkt. Auch wenn die beobachteten Daten keinen Schluss auf eine Korrelation zwischen den Merkmalen Skalarität und Sprechereinstellung zulassen, ist eine solche jedoch nicht grundsätzlich auszuschließen. Dieser Punkt sollte in weiterführenden experimentellen Untersuchungen daher unbedingt berücksichtigt werden.


Auch wenn die genauen Gründe der Skaleneinschätzung auf Grundlage der Daten nicht zweifelsfrei abzuschätzen sind, liefern die durchgeführten Experimente dennoch die erste empirische Evidenz für den in der theoretischen Literatur vertretenen Unterschied zwischen deskriptiven und expressiven Intensivierern. Den Ergebnissen der Untersuchungen nach zu urteilen, werden expressive Intensivierer in der Tat mit einem höheren Grad an Skalarität und Sprechereinstellung assoziiert als deskriptive Intensivierer. Dies spricht klar für die Gültigkeit der in diesem Kontext untersuchten Hypothesen H3 und H4 nach \citet[22]{Paradis.1997} und \citet[133]{Gutzmann.2019a} sowie {\citet[315]{Lang.1983}} und \citet[57]{Ortner.2014}. Ferner zeichnen sich innerhalb der beiden Intensiviererklassen die vermuteten graduellen Abstufungen ab, nach denen die Ausdrücke in Abhängigkeit des durch sie ausgedrückten Grades an Skalarität und Sprechereinstellung als schwach, mittel oder stark klassifiziert werden können. Gleichwohl die Stufen partiell anders aussehen als im Vorfeld der Untersuchung angenommen und in manchen Fällen zudem nicht übermäßig stark ausgebildet sind,\footnote{An dieser Stelle stellt sich die Frage nach der Aufgabenstellung, wurde in Experiment 2 doch expressis verbis nach dem Grad an Bewertung statt dem Grad an Sprechereinstellung wie bei dem in Kapitel \ref{Experiment0} dargelegten Experiment 1 gefragt. Letzteres wurde erst nach der in diesem Kapitel präsentierten Untersuchung durchgeführt und konnte in Bezug auf die Fragestellung optimiert werden. Insofern besteht die Möglichkeit, dass die Ergebnisse mit einer geänderten Aufgabenstellung zu einer klareren Ausdifferenzierung der Stärkegruppen und dadurch deutlicheren Abstufungen geführt hätten, als es hier offenbar der Fall ist.} kann die {Hypothese H6} im Sinne von \citet[87]{Biedermann.1969} und {\citet[408]{Breindl.2007}} ebenfalls als zutreffend befunden werden, auch wenn es aufgrund der niedrigen Datenmenge keinen signifikanten Zusammenhang zwischen Skalarität und Sprechereinstellung gibt, wie {\citet[133]{Gutzmann.2019a}} in seiner Definition eigentlich zu bedenken gibt. Zusammengenommen legen die durch die Experimente erbrachten Ergebnisse nahe, dass sich auch das schwer fassbare Konzept der Expressivität mit experimentellen Methoden untersuchen lässt, was zu weiterer Forschung auf diesem Gebiet einlädt. Dass das Thema Expressivität in der empirischen Sprachwissenschaft von wenigen Ausnahmen abgesehen bisher stiefmütterlich behandelt wurde, kann nicht oft genug betont werden.




%. Wie in Abschnitt \ref{ZuSa} diskutiert, fand sich keine Korrelation, wofür zwei mögliche Erklärungsansätze angeführt wurden. Die erste Erklärung bestand darin, dass die gewählte Methode nicht in intendiertem Maße funktioniert hat, was hauptsächlich der Fragestellung von Experiment 2b geschuldet sein könnte. In diesem Teil der Untersuchung wurden die Versuchspersonen dazu aufgefordert, mithilfe des Schiebereglers anzugeben, wie bewertend der Sprecher der Aussage ist. Dabei ist jedoch möglich, dass die Versuchspersonen nicht nur die expressiven, sondern auch die deskriptiven Intensivierer als bewertend verstanden und letzteren in der Folge einen verhältnismäßig hohen Skalenwert zugewiesen haben. Somit ist vorstellbar, dass der verzeichnete Unterschied trotz statistischer Signifikanz im Falle einer anderen Fragestellung stärker ausgefallen wäre als das Ergebnis von Experiment 2b suggeriert. Die zweite Möglichkeit bezog sich auf die Klassifikation der Intensivierer. In diesem Kontext wurde gemutmaßt, dass die vorgenommene Klasseneinteilung nach deskriptiv und expressiv nicht in Gänze zutreffend war, was sich vor allem bei den gemittelten Skalenwerten der expressiven Intensivierer \textit{super} und \textit{voll} beobachten ließ. Bei diesen Ausdrücken ist denkbar, dass sie aufgrund hohen Lexikalisierungsgrades und damit einhergehender hoher Frequenz bzw. persistentem Gebrauch inzwischen zu abgenutzt sind, um als expressiv wahrgenommen zu werden. Daher ist nicht auszuschließen, dass die Untersuchung mit anderen Intensivierern an der Stelle zu eindeutigeren Ergebnissen und gegebenenfalls auch zu einem korrelativen Zusammenhang zwischen Skalarität und Sprechereinstellung geführt hätte. Dessen ungeachtet legt die Parallelität bei den Einschätzungen der expressiven Intensivierer nahe, dass die Ergebnisse der Untersuchung ebenso repräsentativ sind wie die durch Pseudowörter erzielten Resultate (vgl. Kapitel \ref{Experiment0}).


%Gleichwohl die genauen Gründe der Skaleneinschätzung auf Grundlage der Daten nicht zweifelsfrei abzuschätzen sind, lieferten die durchgeführten Experimente dennoch die erste empirische Evidenz für den postulierten Unterschied zwischen deskriptiven und expressiven Intensivierern, wodurch die eingangs formulierten Hypothesen faktisch als gültig erachtet werden können. Zudem legen die erbrachten Ergebnisse nahe, dass sich auch das schwer fassbare Konzept der Expressivität mit experimentellen Methoden untersuchen lässt, was zu weiterer Forschung auf diesem Gebiet einlädt.

%footnote{Das Symbol X-Bar ($\bar{x}$) gibt hier und im Folgenden den Mittelwert an.}



%%Die Studie, in der deskriptive Intensivierer und expressive Intensivierer konstrastiert wurden...
%%Damit wird/werden die eingangs formulierten Hypothesen (nicht) / erstmal empirisch gestützt
%%Darüber hinaus legt die Studie nahe, dass
%%Damit bleibt die Frage bestehen, worauf sich... zurückführen lässt / zurückzuführen ist
%%Es ist (jedoch) evident, dass...




%Ein Teil der Ergebnisse der Einzelbegutachtung deckt sich zudem mit den Resultaten, die \citet[]{Calpestrati.2017} aus seinen Untersuchungen zum Deutschen und Italienischen gewinnen konnte. Eine seiner Intentionen lag darin, die Unterschiede in der Wahrnehmung der Intensität deutscher und italienischer Steigerungspartikeln im Vergleich zu erörtern. Ohne näher auf die Details der Untersuchung einzugehen, zeigten seine Daten, dass \textit{super} (neben \textit{extra} und \textit{maxi}) von den Probanden als am wenigsten intensiv eingeschätzt wurden, während \textit{mega} (neben \textit{hyper}) als am stärksten beurteilt wurde. Diese Tendenzen zeigten sich auch in meinen Experimenten, bei der \textit{super} zwar nicht den niedrigsten, aber dennoch einen deutlich geringeren Skalenwert aufwies als \textit{mega}, das als am stärksten bewertet wurde.
